
\documentclass[12pt]{article}
%%%%%%%%%%%%%%%%%%%%%%%%%%%%%%%%%%%%%%%%%

%PACKAGES
%--------
\usepackage{amscd}     % nice arrows
\usepackage{amsmath}   % math stuff
\usepackage{amssymb}   % special symbols for math
\usepackage{amsthm}    % gives us the \newtheorem
\usepackage{xr-hyper}
\usepackage{hyperref}  % creates hyper referencing
\usepackage{bookmark}  % creates automatic bookmarks in the pdf by sections
\usepackage{titlesec}
\usepackage[title]{appendix} 
%\usepackage{showkeys} % shows labels on PDF
\setcounter{secnumdepth}{4}
\usepackage{mathtools} % all sorts of math widgets
%\usepackage{epstopdf}  % enables usage of png and eps in the same file
\usepackage{verbatim}  % if you want latex to read it exactly as you write it use verbatim environment
\usepackage{graphicx}  % standard for graphicx includes
\usepackage{color}     % enables font coloring
\usepackage{soul}
\usepackage{scalerel,stackengine}
\usepackage[rightcaption]{sidecap}
\usepackage[skip=2pt,font={small, it}]{caption} % enables caption customizations

\usepackage{algorithm} % Algorithms environment
\usepackage{algorithmicx}
\usepackage{multirow}

\usepackage{algpseudocode}
\usepackage[normalem]{ulem}

%\usepackage[useregional]{datetime2}
\usepackage[UKenglish]{babel}
%\usepackage[a4paper, total={7in, 10in}]{geometry}
\usepackage{authblk}

\usepackage{geometry}
\geometry{
     a4paper,
     total={170mm,257mm},
     left=20mm,
     top=30mm,
     bottom=20mm,
 }

\usepackage{tikz}
	\usetikzlibrary{shapes,arrows}
	\usetikzlibrary{positioning}
	\usetikzlibrary{calc,shapes, positioning}
	\usetikzlibrary{patterns}
	
\DeclareMathAlphabet{\mathpzc}{OT1}{pzc}{m}{it}


% NEW THEOREMS DEFINITIONS
%-------------------------
\newtheorem{definition}{Definition}
\newtheorem{problem}{Problem}
    \newtheorem{Theorem}{Theorem}[section]
\newtheorem{Lemma}[Theorem]{Lemma}
\newtheorem{note}[Theorem]{Note}
\newtheorem{corollary}[Theorem]{Corollary}
\newtheorem{assumption}[Theorem]{Assumption}
\newtheorem{conjecture}[Theorem]{Conjecture}
\newtheorem{prop}[Theorem]{Proposition}
\theoremstyle{remark}
\newtheorem{remark}[Theorem]{Remark}

% --- begin CommandDefinitions.tex ---

% NEW COMMANDS & OPERATORS
%--------------------------
\newcommand{\RR}{\mathbb{R}}
\newcommand{\CC}{\mathbb{C}}
\newcommand{\ZZ}{\mathbb{Z}}
\newcommand{\EE}{\mathbb{E}}
\newcommand{\MM}{\mathcal{M}}
\newcommand{\DD}{\mathcal{D}}
\newcommand{\NN}{\mathcal{N}}
\newcommand{\OO}{\mathcal{O}}
\newcommand{\UU}{\mathcal{U}}
\newcommand{\WW}{\mathcal{W}}
\newcommand{\KK}{\mathcal{K}}
\newcommand{\LL}{\mathcal{L}}
\newcommand{\GG}{\mathcal{G}}
\newcommand{\C}{\mathcal{C}}
\newcommand{\HH}{\mathcal{H}}
\newcommand{\N}{\mathbf{N}}
\newcommand{\PP}{\mathcal{P}}
\newcommand{\XX}{\mathscr{X}}
\newcommand{\JJ}{\mathcal{J}}
\newcommand{\ZZZ}{\mathscr{Z}}
\newcommand{\YY}{\mathscr{Y}}
\newcommand{\LLL}{\mathscr{L}}

\newcommand{\eps}{\varepsilon}
\newcommand{\defeq}{\stackrel{\bigtriangleup}{=}}
\newcommand{\eqdef}{=\vcentcolon}
\newcommand{\ceil}[1]{\left\lceil #1 \right\rceil}
\newcommand{\abs}[1]{\left\vert #1 \right\vert}
\newcommand{\norm}[1]{\left\Vert #1 \right\Vert}
\newcommand{\pdev}[2]{\frac{\partial #1}{\partial #2}}
\newcommand{\dist}{\operatorname{dist}}
\newcommand{\wtilde}[1]{\widetilde{#1}}
\newcommand{\dbtilde}[1]{\stackrel{\approx}{#1}}
\newcommand{\lrangle}[1]{\langle #1 \rangle}
\newcommand{\lrbrackets}[1]{\left( #1 \right)}
\newcommand{\emphbrackets}[1]{\emph{(}#1\emph{)}}
\newcommand{\maxangle}{\angle_{\max}}
\newcommand{\at}[2]{\left. #1 \right\rvert_{#2}}
%\newcommand{\nchoosek}[2]{(\begin{array}{c} #1 \\ #2 \end{array})}
\DeclareMathOperator*{\argmin}{\arg\!\min}

\DeclareMathOperator{\vol}{Vol}
\DeclareMathOperator{\Ima}{Im}
\DeclareMathOperator{\cov}{cov}


% \stackMath
% \newcommand\reallywidehat[1]{%
% \savestack{\tmpbox}{\stretchto{%
%   \scaleto{%
%     \scalerel*[\widthof{\ensuremath{#1}}]{\kern.1pt\mathchar"0362\kern.1pt}%
%     {\rule{0ex}{\textheight}}%WIDTH-LIMITED CIRCUMFLEX
%   }{\textheight}% 
% }{2.4ex}}%
% \stackon[-6.9pt]{#1}{\tmpbox}%
% }
% \parskip 1ex

\newrobustcmd*{\mysquare}[1]{\tikz{\filldraw[draw=#1,fill=#1] (0,0)
rectangle (0.2cm,0.2cm);}}

\newrobustcmd*{\mycircle}[1]{\tikz{\filldraw[draw=#1,fill=#1] (0,0) circle [radius=0.1cm];}}


%%%%%%%%%%%%%%%%%%%%%%%%%%%%%%%%%%%%%%
%%%% Deep lists
%%%%%%%%%%%%%%%%%%%%%%%%%%%%%%%%%%%%%%%
\usepackage{enumitem}
\setlistdepth{9}

\newlist{myEnum}{enumerate}{9}
\setlist[myEnum,1]{label=\arabic*.}
\setlist[myEnum,2]{label=(\alph*)}
\setlist[myEnum,3]{label=\roman*.}
\setlist[myEnum,4]{label=\Alph*.}
\setlist[myEnum,5]{label=(\arabic*)}
\setlist[myEnum,6]{label=(\Roman*)}
\setlist[myEnum,7]{label=(\Alph*)}
\setlist[myEnum,8]{label=(\roman*)}
\setlist[myEnum,9]{label=(\arabic*)}


%%%%%%%%%%%%%%%%%%%%%%%%%%%%%%%%%%%%%%
%%%% Colors
%%%%%%%%%%%%%%%%%%%%%%%%%%%%%%%%%%%%%%%
\definecolor{myPurple}{rgb}{0.4940, 0.1840, 0.5560}
\definecolor{myGreen}{rgb}{0, 0.5, 0}


\allowdisplaybreaks



\newcommand{\mycomment}[1]{
	\textcolor{red}{Comment: #1}%
}
\newcommand{\barak}[1]{
	\textcolor{blue}{#1}%
}% --- end CommandDefinitions.tex ---

\graphicspath{{./Figures/}}


%%%%%%%%%%%%%%%%%%%%%%%%%%%%%%%%%%%%%%%%


\begin{document}
\title{Non-Parametric Estimation of Manifolds from Noisy Data}
\author[1]{Yariv Aizenbud\footnote{authors contributed equally}}
\author[2]{Barak Sober$^*$}
\affil[1]{Department of Mathematics, Yale University}
\affil[2]{Department of Statistics and Data Science, The Hebrew University of Jerusalem}


\maketitle

% --- begin main_abstract.tex ---
\begin{abstract}
A common observation in data-driven applications is that high dimensional data has a low intrinsic dimension, at least locally. 
In this work, we consider the problem of estimating a $d$ dimensional sub-manifold of $\mathbb{R}^D$ from a finite set of noisy samples.
Assuming that the data was sampled uniformly from a tubular neighborhood of $\mathcal{M}\in \mathcal{C}^k$, a compact manifold without boundary, we present an algorithm that takes a point $r$ from the tubular neighborhood and outputs $\hat p_n\in \mathbb{R}^D$, and $\widehat{T_{\hat p_n}\mathcal{M}}$ an element in the Grassmanian $Gr(d, D)$.
We prove that as the number of samples $n\to\infty$ the point $\hat p_n$ converges to $p\in \mathcal{M}$ and $\widehat{T_{\hat p_n}\mathcal{M}}$ converges to $T_p\mathcal{M}$ (the tangent space at that point) with high probability.  
Furthermore, we show that the estimation yields asymptotic rates of convergence of $n^{-\frac{k}{2k + d}}$ for the point estimation and $n^{-\frac{k-1}{2k + d}}$ for the estimation of the tangent space.
These rates are known to be optimal for the case of function estimation.
\end{abstract}

% --- end main_abstract.tex ---

% --- begin main_text.tex ---

\section{Introduction}
Differentiable manifolds are an indispensable language in modern physics and mathematics. 
As such, there is a plethora of analytic tools designed to investigate manifold-based models (e.g., connections, differential forms, curvature tensors, parallel transport, bundles). 
In order to facilitate these tools, one normally assumes access to a manifold's atlas of charts (i.e., local parametrizations). 
Over the past few decades, manifold-based modeling has permeated into data analysis as well (e.g., see  \cite{hastie1984principal,roweis2000LLE, scholkopf1998KPCA}), usually to avoid working in high dimensions, due to the known curse of dimensionality \cite{stone1980optimal}. 
However, in these data-driven models, charts are not accessible and the only information at hand are the samples themselves. 
As a result, a common practice in Manifold Learning is to embed the data in a lower dimensional Euclidean domain, while maintaining some notion of distance (e.g., geodesic or diffusion).
Subsequently, the embedded data is being processed using linear methods on the low dimensional domain; to mention just a few of this body of literature, see  \cite{belkin2003laplacian,belkin2006manifold,coifman2006diffusion,saul2003LLE}. Some of the approaches have robustness guarantees~\cite{ding2020phase,dunson2021spectral,el2016graph,shen2020scalability}.
The main drawback of such dimensionality reduction approaches is that they inevitably lose some of the information in the process of data projection.



In recent years there have been a growing interest in the problem of manifold estimation.
The aim of these approaches is to reconstruct an underlying manifold $\widehat{\MM}$, approximating the sampled one $\MM$, based upon a given discrete sample set.
The first attempt (neglecting the literature dealing with approximations of curves and surfaces \cite{dey2006curve,wendland2004scattered}) at this problem was probably made by Cheng et. al. at 2005 \cite{cheng2005manifold} who present an algorithm that outputs a simplicial complex, homeomorphic to the original manifold $\MM$ and is proven to be close to it in the Hausdorff sense.
However, this algorithm is deemed intractable by the authors as it is based on the creation of Delaunay complexes through Voronoi diagrams in the ambient space. 
Harvesting the idea of tangential Delaunay complexes, Boissonat and Ghosh \cite{boissonnat2014manifold} have provided a method reconstructing a simplicial complex which is computationally tractable (i.e., its complexity has linear dependency in the ambient dimension).
In this approach there is an underlying assumption that the local tangent at each sampled point is given, and they recommend using a local Principle component analysis (PCA) to find these tangents.
The choice of local PCA as an approximating tangent is shown to be a valid one in the analysis given in \cite{aamari2018stability,kaslovsky2011optimal,kaslovsky2014non,singer2012vector}. Furthermore, \cite{aamari2018stability} shows that the estimate given by Boissonat and Ghosh~\cite{boissonnat2014manifold} achieves optimal minmax rates of convergence in case of noiseless samples, with respect to a certain class of manifolds.
Some other works aims at learning multiscale dictionaries to describe the manifold data \cite{allard2012multi}. 

In parallel, meshless methods for the reconstruction of manifolds from point sets have been developed.
Niyogi et. al. \cite{niyogi2008finding} present such an approach through a union of $\varepsilon$-balls around the samples.
They show that this approximant can be homologous to $\MM$ under certain conditions. 
This approach is somewhat similar to the one proposed by Fefferman et. al. in \cite{fefferman2018fitting}, where an analysis of convergence under a Gaussian noise model is given as well. 
Furthermore, Fefferman et. al. \cite{fefferman2019fitting} %\cite{fefferman2016testing}
propose another meshless way of approximating manifolds from point sets up to arbitrarily small Hausdorff distance in the noiseless case, and such that the approximant itself is a smooth manifold of the same intrinsic dimension as $\MM$ (\cite{mohammed2017manifold} uses this framework to provide two more algorithms of such properties).
Faigenbaum-Golovin and Levin uses a generalization of the $L_1$ median to approximate manifolds from meshless data \cite{faigenbaum2020manifold}.
Sober and Levin \cite{sober2016MMLS} give an approximation scheme based upon a generalization of the the Moving Least-Squares (MLS) \cite{levin2004mesh,mclain1976two} that provides a smooth manifold with optimal convergence rates in the noiseless case as well (this approach is referred below as the Manifold-MLS).
Their approximation is built through a two-stage procedure, first estimating a local coordinate system and then building a local polynomial regression over it. 
This framework is extended to deal with approximations of functions over manifolds \cite{sober2017approximation} as well as geodesic distances \cite{sober2020Geodesics}.
Aamari and Levrard \cite{aamari2019nonasymptotic} provide a different algorithm, which is shown to be optimal in the noiseless case as well. 
Differently from Sober and Levin's approach, this algorithm estimates a tangent along with a polynomial estimation above the tangent domain at once.
However, in their practical implementation Aamari and Levrard propose a two step solution (first perform PCA to achieve a tangent and then a polynomial regression above it). Note, that although there are results regarding the convergence of local PCA to the tangent space of some manifold these works assume that the localization is around a point on the manifold itself, which is not given in the current problem setting.


Upper bounds on the minimax rates of convergence were first introduced for smooth manifolds by Genovesse et. al. \cite{genovese2012minimax,genovese2012manifold}. Later, in \cite{kim2015tight}, the same rates were shown to be optimal.
These results were later refined to a class of H\"older-like smooth manifolds by Aamari and Levrard \cite{aamari2019nonasymptotic}.
They come to the conclusion that the optimal rate of convergence for such $k$-times smooth manifold estimation is $O(n^{-k/d})$ for the noiseless case, and is bounded from below by $O((\frac{\sigma}{n})^{k/(k+d)})$ in an orthogonal noise model, where $d$ is the intrinsic dimension of $\MM$, $n$ is the number of samples and $\sigma$ is a bound on the noise level.
However, they do not show that their bound in the noisy case is achievable. 

Note, that in all previous manifold reconstruction algorithms the convergence under noise assumptions is either not analyzed at all, or is guaranteed only when the noise level decays to zero as the sample size $n$ tends to infinity. 
In the current paper, we assume a sample of size $n$ drawn from the uniform distribution on $\MM_\sigma$, a $\sigma$-tubular neighborhood of the manifold $\MM$.
Then, for a given $ r  \in \MM_\sigma $, we present an algorithm that outputs a point $ \hat p_n\in \RR^D$ and $ \widehat{T_{\hat p_n}\MM}\in Gr(d, D) $, a $d$-dimensional linear subspace of $\RR^D$, which estimate $ p\in \MM $ and $ T_{ p}\MM $, the subspace tangent to $\MM$ at $p$.
We prove, in Theorem \ref{thm:MainResult}, that with high probability $ \norm{\hat p_n -  p} = O(n^{-k/(2k+d)}) $ and that $ \maxangle(\widehat{T_{\hat p_n}\MM}, T_{ p}\MM) = \wtilde O(n^{-(k-1)/(2k+d)}) $, where $ \wtilde O $ neglects dynamics weaker than polynomial order (e.g., $ \ln(n) $ and $ \ln(\ln(n)) $).
These achieved convergence rates coincide with the optimal rates of non-parametric estimation of functions \cite{stone1980optimal}.
We note that, in its current formulation, the theorems and proofs in this paper are valid only for noise level $\sigma$  bounded away from zero (the Region Of Interest in the algorithm will approach an empty set as $\sigma \to 0$). 
We believe that this can be fixed, however, this will complicate the proofs which are already complicated enough.
Be that as it may, our approach builds upon the Manifold-MLS \cite{sober2016MMLS}, but differs from it as explained below.
Nevertheless, the analysis performed in \cite{sober2016MMLS} suggests that the convergence rates in case of clean samples are $ \wtilde O(n^{-k/d}) $ and $ \wtilde O(n^{-(k-1)/d}) $ for the point and tangent respectively.

The algorithm presented in this paper is divided into two steps. In step 1 we find a local coordinate system. It is proved in Theorem \ref{thm:Step1} that this coordinate system is a ``reasonable" approximation to the tangent of the manifold at some point. Next, in step 2, we improve the estimation of step 1 in an iterative manner to get an accurate estimation of a point on the manifold and its tangent.  We prove, in Theorem \ref{thm:Step2}, that these iterations indeed converge to an accurate estimate, and show the convergence rates mentioned above. The formal problem setting, along with the algorithm's description are presented in Section \ref{sec:Problem_Setting}. In Section \ref{sec:mainResults} the formal results are presented, where Theorem \ref{thm:MainResult} is the main result of this paper. The theorems of Section \ref{sec:mainResults} are proved in Section \ref{sec:Proofs}. Finally, in Section \ref{sec:applications} we present one possible application of the presented method. Although there are many possible applications (e.g.  denoising, trajectory reconstruction, etc.) , we chose one which can easily be demonstrated visually. The code for the algorithm in this paper, along with examples, can be found in  \textcolor{blue}{\href{https://github.com/aizeny/manapprox}{https://github.com/aizeny/manapprox}}.





\section{Problem Setting and Algorithm Description}
\label{sec:Problem_Setting}
% --- begin Section_ProblemSetting.tex ---
Throughout the paper we limit the discussion to estimation of $\MM\in\C^k$, $k$-times smooth, compact submanifolds of $\RR^D$ without boundary.
This limitation is important for the analysis, but the algorithm we present is local and thus from a practical perspective, a local version of these assumptions should suffice.
A key concept in the analysis of manifold estimates is the \textit{reach} of a manifold (e.g., see the analyses at \cite{aamari2018stability,aamari2019nonasymptotic,fefferman2018fitting,niyogi2008finding}), which was introduced by Federer in \cite{federer1959curvature} and is defined as the maximal distance for which there exists a unique projection onto $\MM$.
\begin{definition}[Reach \cite{federer1959curvature}]\label{def:reach}
The reach of a subset $A$ of $\RR^D$, is the largest $\tau$ \emph{(}possibly $\infty$\emph{)} such that for any $x\in\RR^D$ that maintains $dist(A,x)\leq \tau$, there exists a unique point $P_A(x)\in A$, nearest to $x$. 
\end{definition}
Using the reach we can bound both the local behavior($1/\tau$ bounds all sectional curvatures of the manifold) and global behavior  of a manifold (i.e., it measures how close a manifold can get to itself). 
Thus, the reach provides a good way of expressing our limitations in the problem of manifold estimation (in \cite{niyogi2008finding} the same concept is defined as the \textit{condition number} of a manifold).
For example, if the reach is too small and the sampling density is not fine enough, we would expect that small features could not be recovered.
Moreover, through the reach of a manifold we can define the acceptable levels of noise that do not obscure the geometrical shape.
In accordance with that, we limit our discussion to manifolds with reach bounded away from zero (notice that in case of flat manifolds the reach is infinite) and with a noise model which limits the noise level from above by the reach.

Our noise model in the analysis is as follows: 
We assume that we are given a finite set of samples $\{r_i\}_{i=1}^n$ drawn independently from 
\begin{equation}\label{eq:Msigma}
    \MM_\sigma \defeq \{q\in \RR^D ~|~ dist(q, \MM) < \sigma \}
,\end{equation}
a tubular neighborhood of $\MM$.
Explicitly, we assume that $r_i \sim \operatorname{Unif}(\MM_\sigma)$, which is the uniform distribution on $\MM_\sigma$; i.e., the normalized Lebesgue measure with respect to $\RR^D$.



Finally, let $p\in\MM$, we wish to describe $W_p \subset \MM$ a neighborhood of $p$ as a graph of a function 
\begin{equation}
\varphi_p:W_{T_p\MM}\to \MM \label{eq:LocalChart_clean}    
\end{equation}
where $W_{T_p\MM} = P_{T_p\MM}(W_p)$, the projection of $W_p$ onto the tangent, and $\varphi_p$ is defined by
\begin{equation}
    \varphi_p(x) = p+(x,\phi_p(x))_{T_p\MM}\label{eq:phi_def_clean}
,\end{equation}
where $x\in \RR^d \simeq T_p\MM$, $\phi_p(x)\in \RR^{D-d} \simeq T_p\MM^\perp $, and $(x,y)_{T_p\MM} \in \RR^D$ denotes that $x\in\RR^d$ and $y\in\RR^{D-d}$ represent a point in some basis of $T_p\MM$ and $T_p\MM^\perp$ correspondingly.
Then, we define the graph of $\phi_p$ to be
\begin{equation}\label{eq:FunctionGraph_init}
\Gamma_{\phi_p, W_p} \defeq \{p + (x, \phi_p(x))_{T_p\MM} | x\in W_{T_p\MM} \}  
.\end{equation}
For simplicity, throughout the paper we identify the graph of $\phi_p$ with $\Gamma_{\phi_p,W_p}$. That is, we refer to $\MM$ as locally a graph of $\phi_p$ (see Figure \ref{fig:SkewedH}).

We would like to stress that throughout the paper there is a slight misuse of notation.
Explicitly, when we refer to $T_p\MM$ and $T_p\MM^\perp$, we sometimes look at it as elements in the Grassmannian $Gr(d,D)$ and $Gr(D-d, D)$; i.e., subspaces of $\RR^D$ with dimensions $d$ and $D-d$ correspondingly.
On the other hand, in some other occasions (as in \eqref{eq:FunctionGraph_init}) we neglect the fact that these are subsets of $\RR^D$ which is equivalent to choosing some basis and working in it.

\begin{SCfigure}
		\centering
		\includegraphics[width=0.6\linewidth]{./Figures/SkewedH-3D-v3.pdf}
		\caption{Illustration of $ \varphi_{p} $: The $ xy $-plane is $ H(p) $; The local origin $ q(p) $, which is mapped by $ \phi_{p} $ to $ p $ is marked by the red dot; the vector $ x $ represents a ``tangential" movement; and $ \phi_{p}(x) $ is a normal movement. }
		\label{fig:SkewedH}
\end{SCfigure}


\subsection{Summary of manifold and sampling assumptions}\label{sec:SamplingAssumptions}
Throughout this paper we assume that the (unknown) manifold $\MM$ and the samples $\{r_i\}_{i=1}^n$ satisfy the following requirements: 
\begin{enumerate}
    \item $\MM\in\C^k$ is a compact $d$-dimensional sub-manifold of $\RR^D$ without boundary
    \item $M = \frac{\tau}{\sigma}$ is large enough, where $\tau$ is the reach of $\MM$ and $\sigma>0$ is the noise level. 
    \label{assume:noise}
    \item $\{r_i\}_{i=1}^n$ are samples drawn independently and uniformly from $\MM_\sigma$ (i.e., $r_i\sim \operatorname{Unif}(\MM_\sigma)$).\label{assume:sampling}
\end{enumerate}
% --- end Section_ProblemSetting.tex ---

\subsection{Algorithm Description}\label{sec:Algorithm}
% --- begin Section_AlgorithmDescription.tex ---
As explained above, given a point $r\in\MM_\sigma$ we aim at providing a procedure $\PP(r)$ that will estimate a point $ p \in \MM$. 
This is performed through an altered version of the Manifold-MLS that was introduced in \cite{sober2016MMLS}. 
The Manifold-MLS is constructed through a two-step procedure. 
First, an estimate of a local coordinate system is computed.
Second, above this local coordinate system a local polynomial regression is performed, by which we derive the estimate for the projection onto $\MM$ as well as the tangent domain $T_p\MM$.
Below we show that the first step of the Manifold-MLS yields a reasonable estimate to the tangent even in the presence of noise.
However, performing the second step, which is just a local polynomial regression, above the slightly tilted domain results with a biased estimate.
That is, if we try to estimate the manifold locally as a function of a this approximated domain the noise has bias (see Figure \ref{fig:TiltedBias}).
To account for the bias, in our altered version of the algorithm, we perform the second step iteratively, taking the tangent estimate at each iteration as an improved coordinate system.
We show that in the limit, as the number of samples $n$ approaches $\infty$, our estimate projects $r$ onto $\MM$ and the estimated tangent coincides with the tangent at that projected point.

\begin{SCfigure}
	\centering
	\includegraphics[width=0.5\textwidth]{./Figures/Tilt1-new.pdf}
	\caption{Illustration of a manifold $\MM$ (marked by the blue line) along with its tubular neighborhood $\MM_\sigma$. Assuming uniform sampling in $\MM_\sigma$ we mark the point $p\in\MM$ by \mycircle{myGreen} and the expected value with respect to the given coordinate system by \mysquare{myPurple}. \emph{(}a\emph{)} The coordinate system is aligned with the tangent. \emph{(}b\emph{)} The coordinate system is tilted with respect to the tangent. As can be seen, in (a) the two points coincide, whereas in (b) the expected value differ from the point we wish to estimate.}
	\label{fig:TiltedBias}
\end{SCfigure}

\subsubsection{Step 1 - The Initial Coordinate System}
Given a point $r\in\MM_\sigma$ we limit the the region of interest (ROI) to:
\begin{equation}\label{eq:ROI_clean}
    U_{\textrm{ROI}} = {\{r_i|\dist(r_i,r)<\sqrt{\sigma\tau}\}}
,\end{equation}
and denote the number of samples in the ROI by $N$.
Then, we define the relevant coordinate system as the pair $(q, H)\in \RR^D\times Gr(d, D)$, which minimizes the functional:
\begin{equation}\label{eq:J1_clean}
   J_1(r; q, H) = \frac{1}{N}\sum_{r_i\in U_{\textrm{ROI}}} \dist^2 (r_i - q,H)
\end{equation}
under the constraints
\begin{enumerate}
    \item Orthogonality: $ r - q\perp H $.
    \item Region of interest: $ r_i \in U_{\textrm{ROI}} $.
    \item Search region: $\norm{r-q} < 2\sigma$.
\end{enumerate}
Explicitly, we denote
\begin{equation}\label{eq:argmin1_clean}
    q^*(r), H^*(r) = \argmin_{\substack{ q, H \\ r-q \perp H \\ \|r- q\|<2\sigma}} J_1(r; q, H)
\end{equation}

Note, that Constraint 2 limits our region of interest in accordance with sampling assumptions \ref{assume:sampling} and \ref{assume:noise}.
Furthermore, since $r\in\MM_\sigma$ we know that the true projection onto the manifold is in the search region defined in Constraint 3 (if $q=P_\MM(r)$, the projection of $r$ onto $\MM$ and $H= T_{q}\MM$).
Finally, Constraint 1 extends the notion of orthogonal projection onto manifolds.
As discussed in \cite{sober2016MMLS}, this constraint is responsible for having a unique minimizer for \eqref{eq:argmin1_clean} given enough samples. 
The aforementioned minimization problem is summarized in Algorithm \ref{alg:alg1_clean}.


\begin{algorithm}[H]
	\caption{Step 1: Find an initial coordinate system}
	\label{alg:alg1_clean}
	\begin{algorithmic}[1]
		\Statex {\bfseries Input:}\begin{tabular}[t]{ll}
			$\{r_i\}_{i=1}^n\subset \MM_\sigma$ & noisy samples of a $d$ dimensional manifold $ \MM $.\\
			$r\in \MM_\sigma$ & a point in a tubular neighborhood of $\MM$.
			\end{tabular}
		\Statex {\bfseries Output:}\begin{tabular}[t]{ll}
				$q^*$ & crude estimation of $p = P_\MM(r)$.\\
				$H^*$ & estimation of $T_p\MM$.\\
		\end{tabular}
		\State Disregard points outside of $U_{\textrm{ROI}}$.
		\State Find $q^*, H^*$ minimizing \eqref{eq:argmin1_clean} subject to $\norm{q^* - r}<2\sigma$ and $r - q^* \perp H^*$. 
	\end{algorithmic}
\end{algorithm}


\subsubsection{Step 2 - The iterated projection}
Given $(q, H)\in \RR^D\times Gr(d,D)$ we define the following minimization scheme, known as local polynomial regression (e.g., \cite{cleveland1979robust,mclain1976two}):
Find $\pi\in \Pi_{k-1}^{d\mapsto D}$ a polynomial of total degree $deg(\pi)\leq k-1$ from $\RR^d$ to $\RR^{D-d}$ which minimizes
\begin{equation}\label{eq:Step2_clean}
    J_2(\pi ~|~ q, H) = \frac{1}{N_{q, H}}\sum_{r_i \in U_{\textrm{ROI}}^n} \norm{r_i - (x_i, \pi(x_i))_H}^2
,\end{equation}
where $x_i\in \RR^d$ are the projections of $r_i - q$ onto $H$, $(x, y)_H\in \RR^d\times\RR^{D-d}$ are coefficients in a basis of $H\times H^\perp$, $U_{\textrm{ROI}}^n(q,H)$ is defined through a bandwidth $\epsilon_n$ as
\begin{equation}\label{eq:ROI_ell_clean}
    U_{\textrm{ROI}}^n(q,H) = {\{r_i\in U_\textrm{ROI}~|~\norm{x_i}<\epsilon_n\}}
,\end{equation}
and $N_{q, H}$ denotes the number of samples in $U_{\textrm{ROI}}^n(q, H)$.
Explicitly, the local polynomial regression is defined through
\begin{equation}\label{eq:argmin2_clean}
    \pi^*_{q, H} = \argmin_{\pi\in \Pi_{k-1}^{d\mapsto D}} J_2(\pi ~|~ q, H)
.\end{equation}
As required to ensure convergence in probability for local polynomial regression \cite{aizenbud2021VectorEstimation,stone1980optimal}, we demand that the bandwidth $\epsilon_n\to 0$ as $n\to\infty$ is such that 
\begin{equation}\label{eq:bandwidthDecay_clean}
    0<\lim_{n\rightarrow\infty}n^{1/(2k+1)}\cdot \epsilon_n < \infty
.\end{equation}

Unfortunately, for local polynomial regression, there are no results that relate the probability achieving the required error bound, and the number of samples needed. Such results appear in Theorem 3.2 of \cite{aizenbud2021VectorEstimation} for a slight variant of local polynomial regression, namely some sort of ``median trick"~\cite{alon1999space} on \eqref{eq:argmin2_clean}. For simplicity of notations, we abuse the definition of $ \pi^*_{q, H}$ in \eqref{eq:argmin2_clean}, and define
\begin{equation}\label{eq:pi_star_median_trick}
     \pi^*_{q, H} = \mbox{Algorithm 2 of \cite{aizenbud2021VectorEstimation}} (x_i,y_i),
\end{equation}
where $y_i$ are the projection of $r_i-q$ onto $H^\perp$. 
Any derivative of $\pi^*_{q, H}$ can also be estimated by means of Algorithm 2 of~\cite{aizenbud2021VectorEstimation}. For simplicity of presentation, throughout the paper when we write $\partial \pi^*_{q, H}$ or $\DD_{\pi^*_{q_\ell,H_{\ell}}}$ what we actually mean is the estimate of the derivative rather than the derivative of $\pi^*_{q, H}$.


We begin by setting $q_{-1} = q^*$ (a crude approximation of a point on the manifold) and $H_0 = H^*$ (the initial tangent estimate) resulting from Algorithm \ref{alg:alg1_clean}.
Then, in order to have the origin $q$ closer to $\MM$ we begin by updating $q_0 = q_{-1} + (0, \pi^*_{q_{-1}, H_0}(0))_{H_0}$.
From here on we start updating iteratively the directions of $H$ with respect to a tangent of $\MM$, as well as the point $ q $.
Explicitly, in iteration $\ell$, we define $H_{\ell+1}$ to be the subspace coinciding with $ \Ima(\DD_{\pi^*_{q_{\ell}, H_{\ell}}}[0]) $ the image of the differential of $(\textrm{Id},\pi^*_{q_{\ell},H_{\ell}}):\RR^d\to\RR^D$ at $ 0 $ (i.e., the tangent to the graph of $ \pi^*_{q_{\ell  }, H_{\ell }} $).
In other words, we look at the manifold as a local graph of a function 
\begin{equation}
f_{\ell}:\RR^d\simeq H_\ell \to \RR^{D-d} \simeq H_\ell ^\perp 
.\end{equation}
That is, we define a local patch of the manifold through the graph 
\begin{equation}\label{eq:Graph_fl}
\Gamma_{f_\ell, q_\ell, H_\ell} = \{q_\ell + (x, f_\ell(x))_{H_\ell} | x\in B_{H_\ell}(0,\rho) \}\subset\MM,   
\end{equation}
where $\rho$ is some radius where this function is defined (see Lemma \ref{lem:M_is_locally_a_fuinction_clean} for more details regarding the existence of such $\rho > 0$). 
Then, we estimate the first order differential of $f_\ell$ through taking the differential of the local polynomial regression estimate $\pi^*_{q_{\ell},H_{\ell }}$.
The image of the differential determines $d$-directions in $\RR^D$ (i.e., an element in the Grassmannian $Gr(d,D)$), by which we define $H_{\ell+1}$. 
Following this, we define 
\begin{equation}\label{eq:fl12_def}
    f_{\ell+1/2}:(q_\ell, H_{\ell+1}) \to H_{\ell+1} ^\perp \simeq \RR^{D-d}
\end{equation} 
using the graph $\Gamma_{f_{\ell+1/2}, q_\ell, H_{\ell+1}}$, defined similarly to \eqref{eq:Graph_fl}.
In other words, $f_\ell:(q_\ell, H_\ell)\to H_\ell^\perp$ is defined from a coordinate system with origin at $q_\ell$ and directions $H_\ell$, and $f_{\ell + 1/2}:(q_\ell, H_{\ell+1})\to H_{\ell+1}^\perp$ is defined from a coordinate system with origin at $q_\ell$ and directions $H_{\ell+1}$.
Finally, we update $q_\ell$ by taking 
\begin{equation}\label{eq:ql_definition}
q_{\ell+1} = q_\ell + (0, \pi^*_{q_{\ell }, H_{\ell+1}}(0))_{H_{\ell + 1}} ,
\end{equation}
and then, we have 
\begin{equation}\label{eq:fl_def}
f_{\ell+1}:(q_{\ell+1}, H_{\ell+1}) \to H_{\ell+1} ^\perp \simeq \RR^{D-d}    
\end{equation}

Therefore, the difference between $f_{\ell+1}$ and $f_{\ell +1/2}$ is merely in the location of the origin ($q_{\ell+1}$ instead of $q_\ell$). That is, $q_{\ell+1} - q_\ell \perp H_{\ell+1}$ and
\begin{equation}\label{eq:fell_def}
    f_\ell = f_{\ell + 1/2} - \pi^*_{q_\ell, H_{\ell+1}}(0)
.\end{equation}

We note that the estimate for the first derivative in case of scalar valued functions was analyzed by \cite{stone1980optimal} (as well as others) and was shown to converge to the true derivative with optimal rates in case of unbiased noise. The results of \cite{stone1980optimal} are generalized to vector valued functions in \cite{aizenbud2021VectorEstimation}. 
A core assumption in these results is that $\EE(Y | X=x)$, the expected value of the samples, aligns with the estimated function.
However, in our case this assumption does not hold, since the noise model is tubular with respect to the manifold and unless the coordinate system is aligned with the tangent, the expected value of $Y$ given $X=x$ does not equal to $f_\ell(x)$ (see Figure \ref{fig:TiltedBias}); i.e., the samples have bias.
We show below that the iterations described above improves the maximal angle with respect to a true tangent. 
Thus, eliminating the problem of bias iteratively. 

Finally, after performing $\kappa$ iterations we get the estimate for $ p\in \MM$ and $ T_{ p}\MM $ by 
\begin{equation}
    \hat p_n \defeq q_\kappa ,\quad \widehat{T_{\hat p_n}\MM} \defeq H_\kappa
.\end{equation}
Below we show for a specific value of $\kappa$ that with probability tending to 1 (as the number of samples tend to $\infty$), $ \norm{\hat p_n -  p} =  O(n^{-k/(2k+d)}) $ and that $ \maxangle(\widehat{T_{\hat p_n}\MM}, T_{ p}\MM) = \wtilde O(n^{-(k-1)/(2k+d)}) $, where $ \wtilde O $ neglects dynamics weaker than polynomial order (e.g., $ \ln(n) $ and $ \ln(\ln(n)) $). The aforementioned algorithm is summarized in Algorithm \ref{alg:step2_clean}.

\begin{algorithm}
	\caption{Step 2: estimating the manifold from a good initial guess}
	\label{alg:step2_clean}
	\begin{algorithmic}[1]
		\State {\bfseries Input:}\begin{tabular}[t]{ll}
            $\{r_i\}_{i=1}^n\subset \MM_\sigma$ & noisy samples of a $d$ dimensional manifold $ \MM $.\\
            $q_{-1}$ & a crude approximation of $p = P_\MM(r)$ (normally initialized with $q^*$ of \eqref{eq:argmin1_clean})\\
			%&(normally initialized with $q^*$ of \eqref{eq:argmin1_clean})\\
			$H_0$ & Initial approximation of $T_{p}\MM$ s.t. $\angle_{\max}(H_0,T_{p}\MM) < \alpha_0$ (normally initialized with $H^*$ of \eqref{eq:argmin1_clean})\\
			%&(normally initialized with $H^*$ of \eqref{eq:argmin1_clean})\\
		\end{tabular}
		\State {\bfseries Output:}\begin{tabular}[t]{ll}
				$\hat p_n$ & Estimation of some $ p\in\MM$.\\
				$\widehat{T_{\hat p_n}\MM}$ & Estimation of   $T_{ p}\MM$.\\
		\end{tabular}
		
		\State Compute $\pi^*_{q_{-1}, H_0}$ through a version of linear least-squares minimization of \eqref{eq:pi_star_median_trick}.
        \State $q_0 = q_{-1} + (0,\pi^*_{q_{-1}, H_0}(0))^T$ 
        \For{$\ell=0$ to $\kappa-1$}
		\State Compute  $\pi^*_{q_{\ell},H_{\ell}}$ through a version of linear least-squares minimization  \eqref{eq:pi_star_median_trick}.
		\State Compute $ \DD_{\pi^*_{q_\ell,H_{\ell}}}[0] $ the differential of  $\pi^*_{q_\ell,H_{\ell}}$ in some basis at zero.
		\State $H_{\ell+1} = \Ima(\DD_{\pi^*_{q_\ell,H_{\ell}}}[0]) $ \Comment{this is the column space of $ \DD_{\pi^*_{q_\ell,H_{\ell}}}[0] $}.
		\State Compute $\pi^*_{q_{\ell}, H_{\ell+1}}$ through a version of linear least-squares minimization of \eqref{eq:pi_star_median_trick}.
        \State $q_{\ell+1} = q_\ell + (0, \pi^*_{q_{\ell}, H_{\ell+1}}(0))^T$ 
		\EndFor 
		\State $\hat p_n = q_\kappa$.
		\State $\widehat{T_{\hat p_n}\MM} = H_\kappa $
		\State\Return $\hat p_n$ and $\widehat{T_{\hat p_n}\MM}$.
	\end{algorithmic}
\end{algorithm}
% --- end Section_AlgorithmDescription.tex ---
\subsection{Practical Considerations}\label{sec:practical}
% --- begin Section_PracticalConsiderations.tex ---
\subsubsection{Implementation details}
The minimization problem portrayed in \eqref{eq:argmin1_clean} is non-linear since we optimize for both $q$ and $H$ at the same time; note that if we fix $q$ this amounts to the Principal Component Analysis (which is also related to the iterated linear least-squares problem motivating our algorithms -- see \cite{aizenbud2019approximating}).
This problem has already been studied in \cite{sober2017approximation,sober2016MMLS} and we recommend using the iterative scheme presented in Algorithm \ref{alg:step1Inpractice_clean} to solve it (which is a slight adaptation of the algorithm proposed in \cite{sober2016MMLS}). We note that as the initial estimation of the tangent (step \ref{algstep:U_init_in_practice} in Algorithm \ref{alg:step1Inpractice_clean}) we use the local PCA, which was utilized in many other works and shown to be of merit \cite{aamari2019nonasymptotic,aizenbud2015OutOfSample,kaslovsky2014non}. However, if one wishes to improve the computational complexity, the initialization step can be done in a randomized manner as well \cite{aizenbud2016SVD,halko2011algorithm}.
Algorithm \ref{alg:step1Inpractice_clean} can be shown to converge in theory to a local minimizer of \eqref{eq:argmin1_clean}. 
As explained at length in \cite{sober2016MMLS} under some conditions this minimization has a unique minimum. 
Furthermore, in practical implementations we experienced very fast convergence to a minimum.

As for the practical implementation of Step 2, we note that the derivatives of $f_{\ell + 1/2}$ identify with those of $f_{\ell +1}$.
Finally, the number of iterations $\kappa$ in Algorithm \ref{alg:step2_clean} can be computed explicitly to obtain the rates of convergence as explained in the proofs of Theorem \ref{thm:MainResult}.
However, for the practical implementation, given a specific sample we suggest to iterate until convergence.
See Algorithm \ref{alg:step2InPractice_clean} for the adapted implementation.

\begin{algorithm}
\caption{Step 1: in practice}
\label{alg:step1Inpractice_clean}
\begin{algorithmic}[1]
\State {\bfseries Input:} $\lbrace r_i \rbrace_{i=1}^N, r, \epsilon$
\State{\bfseries Output:}\begin{tabular}[t]{ll}
                         $q$ - an $n$ dimensional vector \\
                         $U$ - an $n\times d$ matrix whose columns are $\lbrace u_j \rbrace_{j=1}^d$ 
                         \end{tabular}
                         \Comment{$H = q + Span\lbrace u_j \rbrace_{j=1}^d$}
\State define $R$ to be an $n\times N$ matrix whose columns are $r_i$
\State initialize $U$ with the first $d$ principal components of the spatially weighted PCA \label{algstep:U_init_in_practice}
\State $q\leftarrow r$
\Repeat
    \State $q_{prev} = q$
    \State $\tilde{R} = (R - repmat(q,1,N))\cdot \Theta$ \Comment{where $\Theta$ is an indicator for points in $U_{\textrm{ROI}}$}
    \State $X_{N\times d} = \tilde{R}^T U$ \Comment{find the representation of $r_i$ in $Col(U)$}
    \State define $\tilde{X}_{N\times (d+1)} = \left[(1,...,1)^T, X\right]$
    \State solve $\tilde{X}^T\tilde{X}\alpha = \tilde{X}^T \tilde{R}^T$ for $\alpha \in M_{(d+1)\times n}$ \Comment{solving the LS minimization of $\tilde{X}\alpha \approx \tilde{R}^T$}
    \State $\tilde{q} = q + \alpha(1,:)^T$
    \State $Q, \hat{R} = qr(\alpha(2:end, :)^T)$ \Comment{where $qr$ denotes the QR decomposition}
    \State $U \leftarrow Q$
    \State $q = \tilde{q} + U U^T (r-\tilde{q})$
\Until {$\|q-q_{\text{prev}}\|<\epsilon$}
\end{algorithmic}
\end{algorithm}

\begin{algorithm}[H]
	\caption{Step 2: in practice}
	\label{alg:step2InPractice_clean}
	\begin{algorithmic}[1]
		\State {\bfseries Input:}\begin{tabular}[t]{ll}
            $\{r_i\}_{i=1}^n\subset \MM_\sigma$ \\
            $q_{-1}$ & a crude approximation of $p_r = P_\MM(r)$ (normally initialized with $q^*$ of \eqref{eq:argmin1_clean})\\
			%&\\
			$H_0$ & Initial approximation of $T_{p}\MM$ s.t. $\angle_{\max}(H_0,T_{p}\MM) < \alpha_0$ (normally initialized with $H^*$ of \eqref{eq:argmin1_clean})\\
			%&(normally initialized with $H^*$ of \eqref{eq:argmin1_clean})\\
		\end{tabular}
		\State {\bfseries Output:}\begin{tabular}[t]{ll}
				$\hat p_n$ & Estimation of some $p\in\MM$.\\
				$\widehat{T_{\hat p_n}\MM}$ & Estimate of  $T_p\MM$.\\
		\end{tabular}
		
		\State Compute $\pi^*_{q_{-1}, H_0}$
		\State $q_0 = q_{-1} + (0,\pi^*_{q_{-1}, H_0}(0))_{H_0}$ 
        \Repeat %{$\ell=0$ to $\kappa-1$\mycomment{until convergence?}}
        \State Compute $\pi^*_{q_\ell,H_{\ell}}$ through the least-squares minimization of \eqref{eq:argmin2_clean}.
        \State Compute $\DD_{\pi^*_{q_\ell,H_{\ell}}}[0]$ in some basis.
        \State Set $H_{\ell + 1} = \Ima(\DD_{\pi^*_{q_\ell,H_{\ell}}}[0])$
        \State Set $q_{\ell+1} = q_\ell + (0, \pi^*_{q_\ell,H_{\ell}}(0))_{H_\ell}$
		\Until {$|q_{\ell} - q_{\ell +1}| \leq \epsilon$}
		\State \Return $\hat p_n = q_{\ell+1}$,  $\widehat{T_{\hat p_n}\MM} = H_{\ell+1} $
		%\State\Return $\hat p_n$ and $\widehat{T_{\hat p_n}\MM}$.
	\end{algorithmic}
\end{algorithm}

% \subsubsection{\barak{Smoothness and manifoldness}}
% \mycomment{I want to remove this subsubsection as well}

% The algorithm presented here differs from the Manifold-MLS \cite{sober2016MMLS} in yet another aspect other than the iterative application of Step 2.
% Namely, the Manifold-MLS localizes the problem by means of smooth weighting functions applied in Steps 1 and 2, whereas in the present algorithm we use an indicator weight function in both steps.
% We chose this discontinuous weighting scheme to simplify the statistical analysis.
% However, by doing so, the fact that $\widehat \MM$ is a smooth $d$-dimensional manifold shown in \cite{sober2016MMLS} could not be inherited into our altered algorithm directly.
% Be that as it may, the proofs portrayed below can be reproduced using smooth and compactly supported weight functions as well (or even non-compactly supported but fast decaying).
% Thus, in practice, if a smooth localization scheme is used in both steps one achieves both smoothness and our estimation rates at the same time.

% --- end Section_PracticalConsiderations.tex ---

\section{Main Results}\label{sec:mainResults}
% --- begin Section_MainResults.tex ---
The main result reported in this paper is Theorem \ref{thm:MainResult}.
For convenience we wish to reiterate the sampling assumptions presented above in Section \ref{sec:SamplingAssumptions}, as they are relevant for all the following theorems.
Namely, we assume that 
\begin{enumerate}
    \item[i] $\MM\in\C^k$ is a compact $d$-dimensional sub-manifold of $\RR^D$ without boundary.
    \item[ii] $M = \frac{\tau}{\sigma}$ is large enough, where $\tau$ is the reach of $\MM$ and $\sigma>0$ is the noise level. 
    \item[iii] $\{r_i\}_{i=1}^n$ are samples drawn independently and uniformly from $\MM_\sigma$ (i.e., $r_i\sim \operatorname{Unif}(\MM_\sigma)$).
\end{enumerate}



\begin{Theorem}\label{thm:MainResult}
    Assuming $M > C_\tau \sqrt{D\log D}$ for some constant $C_\tau$ independent of $\tau$, and let $r\in \MM_\sigma$.
    Then, for any $\delta>0$ arbitrarily small, there exists $ N $ such that for any number of samples $ n > N $,
    applying Algorithm~\ref{alg:step2_clean} with inputs $q_{-1}, H_0$ being the outputs of Algorithm~\ref{alg:alg1_clean}, and with number of iterations $\kappa$ dependent on $n,d, \delta, k$, we get $ \hat p_n, \widehat{T_{\hat p_n}\MM} $, for which 
	\[
	\norm{\hat p_n - \hat p} \leq \frac{C\ln\left(\frac{1}{\delta}\right)}{n^{r_0}}
	,\]
	and 
	\[
	\angle_{\max}(\widehat{T_{\hat p_n}\MM},T_{\hat p}\MM) \leq  C_{d}\ln\left(\frac{1}{\delta}\right) \left(\frac{ n}{\ln\left(\ln(n) \right)^2}\right)^{-r_1} = \wtilde{\mathcal{O}}(n^{-r_1}) 
	\]
	with probability of at least $ 1 - \delta $, where $ r_0 = \frac{k}{2k +d} $, $ r_1 = \frac{k-1}{2k +d} $ and $\hat p$ is some point in $\MM$.
\end{Theorem}


We derive this result by showing that Algorithm \ref{alg:alg1_clean} yields a ``reasonable" estimation for the tangent directions (Theorem \ref{thm:Step1}), and the fact that Algorithm \ref{alg:step2_clean} yields estimates that converge  to a point and its tangent on the original manifold as $n \to \infty$ (Theorem \ref{thm:Step2}). Accordingly, Theorem \ref{thm:MainResult} can be proven directly from Theorems \ref{thm:Step1} and \ref{thm:Step2}.


\begin{Theorem}\label{thm:Step1}
	Let $ (q^*(r), H^*(r)) $, the output of Algorithm \ref{alg:alg1_clean}, and let $p = P_\MM(q^*)$. 
	Denote $\alpha = \sqrt{C_M/M}$ for some constant $C_M$ (independent of $\alpha$ and $M$).
	Then, for any $\delta>0$ arbitrarily small, there exists $ N_{\delta} $ such that for all $ n > N_{\delta} $ 
	\[\angle_{\max}(H^*, T_p\MM)\leq \alpha\]
	with probability of at least $ 1 - \delta $.
	Furthermore, we have
	\[
	\norm{p - q^*} \leq 3 \sigma
	\]
\end{Theorem}

The proof of this theorem can be found in Section \ref{sec:Step1Analysis}.
\begin{Theorem}\label{thm:Step2}
	Assume that $M>C_\tau\sqrt{D\log D}$.  
	Let $(q, H)$ be  a coordinate system, for which $\norm{q - p}\leq 3 \sigma$ and $ \maxangle (H, T_p\MM) \leq \alpha_0$ for some $p \in \MM$ and $\alpha_0 = \sqrt{C_M/M}$.
	For any $\delta>0$ arbitrarily small, denote by $ \hat p_n, \widehat{T_{\hat p_n}\MM} $ the estimates derived from Algorithm \ref{alg:step2_clean} initialized with $(q,H)$ with the number of iterations $\kappa$ specified in Lemma \ref{lem:comute_kappa}. Then, there are $ C_{\delta,d,k} $ and $ N_{\delta} $ such that for all $ n > N_{\delta}  $ there is $\mathbf{p}\in \MM$ for which 
	\begin{equation}\label{eq:thm3.3_p_bound}
	\norm{\hat p_n - \mathbf p} \leq \frac{C \ln\left(\frac{1}{\delta}\right)}{n^{r_0}}
	,\end{equation}
	and 
	\begin{equation}\label{eq:thm3.3_tangent_ang}
	\angle_{\max}(\widehat{T_{\hat p_n}\MM},T_{\mathbf p}\MM) \leq C_{d}\ln\left(\frac{1}{\delta}\right) \left(\frac{ n}{\ln\left(\ln(n) \right)^2}\right)^{-r_1} = \wtilde{\mathcal{O}}(n^{-r_1})    
	\end{equation}
	with probability of at least $ 1 - \delta $, where $ r_0 = \frac{k}{2k +d} $ and where $ r_1 = \frac{k-1}{2k +d} $.
\end{Theorem}
The proof of this theorem can be found in Section \ref{sec:Step2Analysis}.


% --- end Section_MainResults.tex ---

\section{Proofs}
\label{sec:Proofs}
% --- begin Section_Proofs.tex ---
Theoretically, if we had known the tangent bundle of the sampled manifold at every point, we could have utilized it as a moving frame for the ``$x$-domain" to simply perform  a Moving Least-Squares approximation.
In this case, the convergence analysis would have been similar to standard local polynomial regression \cite{stone1980optimal} (with a varying coordinate system), as the sample bias issue described above would not have occurred. 
Thus, the first part of our investigation is focused on proving that $(q^*(r), H^*(r))$, the solution to the minimization problem of Equation \eqref{eq:argmin1_clean}, yields crude approximations to a tangent of the manifold.

Then, we refine the coordinate system in order to prevent bias introduced by the fact that $H^*(r)$ is tilted with respect to $T_{p}\MM$.
Yet, as we show below, one of the keys to unlocking the convergence rates are the known rates for local polynomial regression.


% --- end Section_Proofs.tex ---
\subsection{Proof of Theorem \ref{thm:Step1}}
\label{sec:Step1Analysis}
% --- begin Section_ProofStep1.tex ---
\begin{proof}[proof of Theorem \ref{thm:Step1}]
The proof can be described by the following three arguments which are proven in Lemmas \ref{lem:J1pTp_clean}, \ref{lem:J_1qH_clean}.
\begin{enumerate}
    \item[Arg. 1:]\label{outline:Step1_1_clean}
    Denote $p_r = P_\MM(r)$.
    Then, since $r - p_r\perp T_{p_r}\MM $ and $r\in\MM_\sigma$, we have that $q = p_r$ along with $H= T_{p_r}\MM$ are in the search space defined by the constraints of \eqref{eq:argmin1_clean}.
    
    \item[Arg. 2:]\label{outline:Step1_2_clean} 
    From Lemma \ref{lem:J1pTp_clean} it follows that for large enough $M = \tau/\sigma$ such that $ \sqrt{\frac{\sigma}{\tau}}  + \frac{\sigma}{\tau} < \frac{1}{2} $ we have $J_1(r; p_r; T_{p_r}\MM) \leq 50\cdot \sigma^2$. 
    According to Assumption \ref{assume:noise} in Section \ref{sec:SamplingAssumptions}, we have that $M$ is large enough, and Lemma \ref{lem:J1pTp_clean} hold.
    Due to the definition of $q^*, H^*$ in \eqref{eq:argmin1_clean}, we achieve that $J_1(r; q^*, H^*) \leq 50 \cdot \sigma^2$ as well.
    
    \item[Arg. 3:]\label{outline:Step1_3_clean} From Lemma \ref{lem:J_1qH_clean}, we have that  
    for $\alpha = \sqrt{C_M/M}$, where $M = \frac{\tau}{\sigma}$, and $C_M$ is a constant
    the following holds:
    For any $\delta >0$ arbitrarily small there is $ N_{\delta} $ sufficiently large (independent of $\alpha$) such that for all $ n>N_{\delta} $, \textbf{all} pairs $(q, H)$ in the search space of \eqref{eq:argmin1_clean} such that $\angle_{max}(H,T_{p_r}\MM) > \alpha$, the score $J_1(r; q, H) \geq 100 \sigma^2 $ with probability of at least $1 - \delta$.
\end{enumerate}
Combining Arguments 2 and 3 we have that 
for $\alpha = \sqrt{C_M/M}$, where $M = \frac{\tau}{\sigma}$, and $C_M$ is a constant
the following holds:
	For any $\delta>0$ arbitrarily small, there exists $ N_{\delta} $ (independent of $\alpha$) such that for all $ n > N_{\delta} $ 
	\[\angle_{\max}(H^*, T_p\MM)\leq \alpha\]
	with probability of at least $ 1 - \delta $. Additionally, since the search space of \eqref{eq:argmin1_clean} requires that $\|r-q^*\| <2\sigma$, we have that $\|p-q^*\|\leq 3\sigma$, and the proof is concluded.
\end{proof}

\begin{Lemma}\label{lem:J1pTp_clean}
Let the sampling assumption of \ref{sec:SamplingAssumptions} hold, let ${p_r} = P_\MM(r)$ be the orthogonal projection of $r$ onto $\MM$, and let $ T_{p_r}\MM $ be the tangent to $ \MM $ at $ {p_r} $.
Then, for $M$ (of Assumption \ref{assume:noise}) large enough
    \begin{equation}\label{eq:TangentScore_clean}
       J_1(r; {p_r}, T_{{p_r}}\MM) \leq 50 \cdot \sigma^2 
    \end{equation}
\end{Lemma}
\textbf{Idea of the proof: } Since all the sampled points are $\sigma$-close to the manifold which is linearly approximated by the tangent, the mean squared distance to the tangent is of the order of $\OO(\sigma^2)$.
The proof is given in Appendix \ref{subsec:proof_lem_J1pTp_clean}


We show, in Lemma \ref{thm:J1pH_clean} that given a coordinate system $(p_r, H)$ with $\angle_{max}(H,T_{{p_r}}\MM) > \alpha$ with yields a large score of of our cost $J_1$ with high probability.
This will be generalized to a coordinate system $(q, H)$ around any origin $q$ in the search space of \eqref{eq:argmin1_clean} in Lemma \ref{lem:J_1qH_clean}.

\begin{Lemma}\label{thm:J1pH_clean}
    Let the sampling assumptions of Section \ref{sec:SamplingAssumptions} hold.
    Let ${p_r} = P_\MM(r)$ be the projection of $r$ onto $\MM$, and $ T_{p_r}\MM $ be the tangent to $ \MM $ at $ {p_r} $. 
    For $\alpha = \sqrt{C_M/M}$, where $M = \frac{\tau}{\sigma}$, and $C_M$ is a constant,
    the following holds: For any $\delta >0$ there is $N_\delta$ (independent of $\alpha$) such that $\forall n > N_\delta$ all linear sub-spaces $H\in Gr(d,D)$ with $\angle_{max}(H,T_{{p_r}}\MM) > \alpha$, yield a score 
    \[J_1(r; {p_r}, H) \geq 109 \sigma^2 ,\]
    with probability of at least $1 - \delta$.
\end{Lemma}

\begin{proof}
We first wish to denote by $\beta_j$ the principal angles between $H$ and $T_{p_r}\MM$ and their matching principal pairs $(u_j, w_j)\in T_{p_r}\MM\times H$ (see Definition \ref{def:principal_angles_clean}).  
Throughout the proof we work on the sectional planes defined by $\LL_j = Span\{u_j,w_j\}$.
Thus, we can define the orthogonal complement of $u_j$ and $w_j$ on $\LL_j$ by $y_j$ and $\tilde y_j$ correspondingly.
That is, both $\{u_j, y_j\}$ and $\{w_j, \tilde y_j\}$ are orthogonal bases of $\LL_j$.
Since for any $i \neq j$ we have that $u_i \perp w_j$, $u_i \perp u_j$ and $w_i \perp w_j$, we have that both $\{\tilde y_j\}_{j=1}^{d}$ and $\{y_j\}_{j=1}^{d}$ are orthonormal sets.
Then, complete the sets $\{\tilde y_j\}_{j=1}^{d}$ and $\{y_j\}_{j=1}^{d}$  to an orthonormal basis of $H^\perp$  and $T_{p_r}\MM^\perp$ through adding the orthonormal sets $\{\tilde y_j\}_{j=d+1}^{D-d}$ and $\{y_j\}_{j=d+1}^{D-d}$ correspondingly.
Explicitly, we know that for all $j' = 1\ldots, d$
$$\textrm{Span}\{\tilde y_j\}_{j=d+1}^{D-d}\perp \LL_{j'} ~~,~~ \textrm{Span}\{ y_j\}_{j=d+1}^{D-d}\perp \LL_{j'}.$$ 
Thus, for $j=d+1, \ldots, D-d$ we have both $\tilde y_j , y_j\in H^\perp \cap T_{p_r}\MM^\perp$ and without limiting the generality we can choose $\tilde y_j = y_j$ for such $j$.
Using this notation we get that for any point $x\in U_\textrm{ROI}$
\[
\dist^2(x - {p_r}, H) = 
\sum_{j=1}^{D-d} \lrangle{x - {p_r}, \tilde y_j}^2 =
\sum_{j=1}^{D-d} \lrangle{x - {p_r}, \tilde y_j}^2 + \sum_{j=1}^{D-d} \left[\lrangle{x - {p_r},  y_j}^2 - \lrangle{x - {p_r},  y_j}^2\right] =
\]
\[
=\sum_{j=1}^{D-d} \lrangle{x - {p_r}, y_j}^2 + \sum_{j=1}^{D-d}\left[ \lrangle{x - {p_r},  \tilde y_j}^2 - \lrangle{x - {p_r},  y_j}^2\right]
,\]
and since $\tilde y_j = y_j$ for $j=d+1,\ldots,D-d$ we have
\begin{align}\label{eq:EJ_1_clean}
	\dist^2(x - {p_r}, H) &=
	\dist^2(x - {p_r}, T_{{p_r}}\MM) + \sum_{j=1}^{d}\left[ \langle x - {p_r}, \tilde{y}_j \rangle^2 - \langle x - {p_r}, {y}_j \rangle^2\right]
.\end{align}

The remainder of the proof is achieved through the following set of claims:
\begin{enumerate}
    \item\label{thm:claim_J1H_clean} From \eqref{eq:EJ_1_clean} it follows that
    \[J_1(r; {p_r}, H) =  J_1(r;{p_r},T_{p_r}\MM) +  R_1(r;{p_r},H),\]
    where
    \begin{equation}\label{eq:sum_k_clean}
    R_1(r; {p_r}, H) = \sum\limits_{j=1}^{d} \frac{1}{\#|U_\textrm{ROI}|}\sum\limits_{r_i\in U_\textrm{ROI}}  \langle r_i - {p_r}, \tilde{y}_j \rangle^2 - \langle r_i - {p_r}, {y}_j \rangle^2
    .\end{equation}
    \item Thus, in order to bound $J_1(r; {p_r}, H)$ from below we can focus on bounding $R_1(r; {p_r}, H)$.
    We then consider separately two sets of indices $\KK', \KK''$ such that $\KK' \cup \KK'' = \{1,\ldots, d\}$, where for all $j\in \KK'$ we have $\beta_j > \alpha$ and for all $j\in \KK''$ we have $\beta_j\leq \alpha$.
    Notice that since $\angle_{max}(H, T_{p_r}\MM) > \alpha$ we, get that $\KK' \not = \emptyset$ and that $\#\KK'' \leq d-1$.
    Writing this explicitly we get
    \begin{equation}\label{eq:RpRpp_clean}
        R_1(r; {p_r}, H) = R_1'(r; {p_r}, H)  +  R_1''(r; {p_r}, H)
    \end{equation}
    where 
    \begin{equation}\label{eq:R1p_def_clean}
    R_1'(r; {p_r}, H) = \sum\limits_{j\in\KK'} \frac{1}{\#|U_\textrm{ROI}|}\sum\limits_{r_i\in U_\textrm{ROI}}  \langle r_i - {p_r}, \tilde{y}_j \rangle^2 - \langle r_i - {p_r}, {y}_j \rangle^2
    ,\end{equation}
    and
    \begin{equation}\label{eq:R1pp_def_clean}
    R_1''(r; {p_r}, H) = \sum\limits_{j\in\KK''} \frac{1}{\#|U_\textrm{ROI}|}\sum\limits_{r_i\in U_\textrm{ROI}}  \langle r_i - {p_r}, \tilde{y}_j \rangle^2 - \langle r_i - {p_r}, {y}_j \rangle^2
    .\end{equation}
    \item\label{thm:claim_R1pp_bound_clean} We show in \ref{sec:proof_claim_R1pp_bound_clean} that
    \begin{equation}\label{eq:R1pp_bound_clean}
    R_1''(r; {p_r}, H) \geq
    -9 \sigma^2
    .\end{equation}
    \item\label{thm:claim_ri_tyk_bound_clean} Then, when we focus on $j\in\KK'$ we show that for $a_1 = \frac{1}{4}$ and $a_2 = \frac{1}{8}$ and $x_0\in T_{p_r}\MM$ such that $\norm{{p_r} - x_0} = a_1\sqrt{\sigma\tau}$ and $B_{T_{p_r}\MM}(x_0, a_2 \sqrt{\sigma\tau} + \sigma)\subset B_D({p_r}, \sqrt{\sigma \tau} - \sigma) \subset U_\textrm{ROI}$, and 
    \begin{equation}\label{eq:ri_tyk_bound_clean}
        \lrangle{r_i - {p_r}, \tilde y_j}^2 \geq \frac{1}{64}\sigma \tau sin\alpha - \frac{1}{2}\sigma^{3/2}\tau^{1/2}
    .\end{equation}
    \item\label{thm:claim_useStochasticLemma} From Assumption \ref{assume:noise} in Section \ref{sec:SamplingAssumptions} we know that $ \sqrt{\frac{\sigma}{\tau}} < \frac{1}{2} $. 
    In addition, setting $ \rho = a_2 \sqrt{\sigma \tau} $ we know that $ B_{T_{p_r}\MM}(x_0, \rho+\sigma)\subset B_{T_{p_r}\MM}({p_r}, \sqrt{\sigma \tau} - \sigma) $, so we can use Lemma \ref{lem:num_of_samples_in_a_ball_clean}.
    Explicitly, let $\mathrm{V}_d = \frac{\pi^{d/2}}{\Gamma(d/2+1)}$ be the volume of a $ d $-dimensional unit ball, and denote 
     \begin{align*}
    \mu_{\min} &= \mathrm{V}_{D-d}\sigma^{D-d} \min_{\substack{p\in \MM\\ x\in B_{T_p\MM}(0,\sqrt{\sigma\tau}- \sigma)}}  \sqrt{det(G_p(x))}\\
    \mu_{\max} &= \mathrm{V}_{D-d}\sigma^{D-d}\max_{\substack{p\in \MM \\ x \in B_{T_p\MM}(0,\sqrt{\sigma\tau}- \sigma)}}  \sqrt{det(G_p(x))}
    ,
    \end{align*}
    where $ \vol(B_{{T_p\MM}^\perp}(\sigma)) $ is the volume of a $ D-d $ dimensional ball with radius $ \sigma $ and $ G_p $ is the matrix representing the Riemannian metric at $ p $ in the chart $ \varphi_p $ of \eqref{eq:phi_def_clean}.
    Then, from Lemma \ref{lem:num_of_samples_in_a_ball_clean}, we get that for any $ \eps, \delta $ there is $ N $ such that for all $ n > N $
    \begin{align*}
    &\#\{r_i|P_{T_p\MM}(r_i) \in B_{T_p\MM}(x_0,\rho) \}
    \leq
    n(2\cdot \mu_{min}\cdot \mathrm{V}_d \cdot\rho^d+\eps) \\
    &\#\{r_i|P_{T_p\MM}(r_i) \in B_{T_p\MM}(x_0,\rho+\sigma) \}
    \geq
    n \left(\frac{\mu_{max}}{2}\cdot \mathrm{V}_{d}\cdot \rho^d - \eps\right)
    \end{align*}
    with probability of at least $1-\delta$.
    By setting $ \eps $ appropriately we get 
    \begin{align*}
    &\#\{r_i|P_{T_p\MM}(r_i) \in B_{T_p\MM}(x_0,\rho) \}
    \leq
    3 n \cdot\mu_{min}\cdot \mathrm{V}_d \cdot\rho^d \\
    &\#\{r_i|P_{T_p\MM}(r_i) \in B_{T_p\MM}(x_0,\rho+\sigma) \}
    \geq
    \frac{n}{3}  \cdot\mu_{max}\cdot \mathrm{V}_{d} \cdot \rho^d 
    .\end{align*}
    \item\label{thm:claim_R1p_bound_clean} Rephrasing \ref{thm:claim_ri_tyk_bound_clean}-\ref{thm:claim_useStochasticLemma} as one statement, there is a large subset of $U_\textrm{ROI}$ with large values of $\lrangle{r_i - p, \tilde y_j}^2$, with probability of at least $1-\delta$.
    Accordingly, we show that for any $ \delta $ there is $ N_\delta $ such that for all $ n>N_\delta $ 
    \begin{align}
    R_1'(r; p, H) \geq  &\frac{\mu_{max}}{\mu_{min} }\cdot\frac{1}{9} \left(\frac{1-\sqrt{\sigma/\tau}}{ 8}\right)^d\cdot \left(\frac{1}{64}\sigma \tau sin\alpha - \frac{1}{2}\sigma^{3/2}\tau^{1/2} \right)  -9 \sigma^2
    \label{eq:R1p_bound_clean},\end{align}
    with probability of at least $1-\delta$.
    \item Combining \eqref{eq:RpRpp_clean} with \eqref{eq:R1pp_bound_clean} and \eqref{eq:R1p_bound_clean}, we get that for $\alpha = \sqrt{C_M/M}$, where $M = \frac{\tau}{\sigma}$, and $C_M$ is a constant the following holds:
    For any $ \delta > 0$ there is $ N $ such that for all $ n > N $ 
    \begin{equation}
        R_1(r;p,H) \geq 109 \sigma^2
    ,\end{equation}
    with probability of at least $1-\delta$.
    \item From Claim \ref{thm:claim_J1H_clean} above, since $J_1(r;{p_r},T_{p_r}\MM) \geq 0$ we achieve that with probability of at least $1-\delta$ there is $N_\delta$ large enough such that for all $ n > N_\delta $
    \[
       J_1(r; p, H) \geq 109 \sigma^2  
    ,\]
    as required.
    
\end{enumerate}



\begin{figure}
	\centering
	\includegraphics[width=0.3\textwidth]{Lk.pdf}
	\caption{ Illustration of $\mathcal{L}_{\tilde j}$ and the sections $L_1,L_2,L_3,L_4$. $U_{bad}$ is marked in light green and $U_{good}$ is marked in light pink. The angles $\angle (\tilde y_{\tilde j}, y_{\tilde j}), \angle (w_{\tilde j},u_{\tilde j})$ which equal to $\beta_{\tilde j}$ are marked in black. The bisectors of these angles are marked in dotted red lines.}
	\label{fig:Lk_clean}
\end{figure}
For clarity of presentation we wish to first prove Claim \ref{thm:claim_R1p_bound_clean} and only then show the correctness of the formula \eqref{eq:ri_tyk_bound_clean} presented in Claim 
\ref{thm:claim_ri_tyk_bound_clean}.
\paragraph*{Proof of Claim \ref{thm:claim_R1p_bound_clean}:}
For all $j\in \KK'$ we have $\beta_j>\alpha$.
Let us assume that there is only one index $\tilde j$ in $\KK'$ (otherwise we can treat each index separately and arrive at the same conclusion), then \eqref{eq:R1p_def_clean} can be rewritten as
\[
R_1'(r; p, H) = \frac{1}{\#|U_\textrm{ROI}|}\sum\limits_{r_i\in U_\textrm{ROI}}  \langle r_i - p, \tilde{y}_{\tilde j} \rangle^2 - \langle r_i - p, {y}_{\tilde j} \rangle^2
.\]
Thus, the only property which affect the score of $R_1'$ is the difference between the measurements $\langle r_i - p, \tilde{y}_{\tilde j} \rangle^2$ and $\langle r_i - p, {y}_{\tilde j} \rangle^2$, both on the 2D plane $\LL_{\tilde j}$.
Accordingly, using the bisector of $\angle(\tilde y_{\tilde  j}, y_{\tilde  j})$ and its orthogonal complement, we can split $\LL_{\tilde j}$ into four regions $L_1, L_2, L_3, L_4$ (see Figure \ref{fig:Lk_clean}), where in two regions ($L_2$ and $L_4$ in Figure \ref{fig:Lk_clean}) $\lrangle{r_i - p, \tilde y_{\tilde j}} \geq \lrangle{r_i - p, y_{\tilde j}}$ and in the other two regions ($L_1$ and $L_3$ in Figure \ref{fig:Lk_clean}) $\lrangle{r_i - p, \tilde y_{\tilde j}} \leq \lrangle{r_i - p, y_{\tilde j}}$.
By denoting 
$$U_{bad} = \{r_i \in U_\textrm{ROI} ~|~ P_{\LL_{\tilde j}}(r_i - p)\in L_1 \cup L_3\},$$
and
$$U_{good} = \{r_i \in U_\textrm{ROI} ~|~ P_{\LL_{\tilde j}}(r_i - p)\in L_2 \cup L_4\}$$
we get that
\[
R_1'(r; p, H) = \frac{1}{\#|U_\textrm{ROI}|}\left[\sum\limits_{r_i\in U_{good}}\left[  \langle r_i - p, \tilde{y}_{\tilde j} \rangle^2 - \langle r_i - p, {y}_{\tilde j} \rangle^2\right] + \sum\limits_{r_i\in U_{bad}} \left[ \langle r_i - p, \tilde{y}_{\tilde j} \rangle^2 - \langle r_i - p, {y}_{\tilde j} \rangle^2\right]\right]
\]
\[
R_1'(r; p, H) \geq \frac{1}{\#|U_\textrm{ROI}|}\left[\sum\limits_{r_i\in U_{good}}  \langle r_i - p, \tilde{y}_{\tilde j} \rangle^2 - \sum\limits_{r_i\in U_\textrm{ROI}} \langle r_i - p, {y}_{\tilde j} \rangle^2\right]
.\]
Similar to \eqref{eq:ri_yk_bound_s_square_clean}, 
\[
\frac{1}{\#|U_\textrm{ROI}|}\sum\limits_{r_i\in U_\textrm{ROI}} \langle r_i - p, {y}_{\tilde j} \rangle^2 \leq 9 \sigma^2
,\]
and thus
\[
R_1'(r; p, H) \geq \frac{1}{\#|U_\textrm{ROI}|}\sum\limits_{r_i\in U_{good}}  \langle r_i - p, \tilde{y}_{\tilde j} \rangle^2
-9\sigma^2
.\]
Therefore, all we need to show is that given $n$ large enough, there are enough samples in $U_{good}$ for which the value $\lrangle{r_i - p, \tilde y_{\tilde j}}^2$ is large enough.
Using Lemma \ref{lem:num_of_samples_in_a_ball_clean}, as described in Claim \ref{thm:claim_useStochasticLemma}, since $U_\textrm{ROI}\subset B_D(P_{\MM}(r), \sqrt{\sigma \tau}+\sigma)$, then for any $\delta$ there is $N$ large enough such that for all $n>N$  with probability of at least $1 - \delta$ 
\[
\#|U_\textrm{ROI}| < 3n\cdot\mu_{min}\cdot \mathrm{V}_{d}\cdot (\sqrt{\sigma\tau}+\sigma)^d 
.\]
Thus,
\[
R_1'(r; p, H) \geq \frac{1}{3n\cdot\mu_{min}\cdot V\cdot ( \sqrt{\sigma\tau} +\sigma)^d }\sum\limits_{r_i\in U_{good}}  \langle r_i - p, \tilde{y}_{\tilde j} \rangle^2
-9\sigma^2
.\]
Below in the proof of Claim \ref{thm:claim_ri_tyk_bound_clean} we show that \eqref{eq:ri_tyk_bound_clean} holds (the proof below is independent of the current one, but utilizes the notion of $U_{good}$ defined above).
Explicitly, for $ r_i \in U_{good} $
\begin{equation*}
\lrangle{r_i - p, \tilde y_j}^2 \geq \frac{1}{64}\sigma \tau \sin\alpha_0 - \frac{1}{2}\sigma^{3/2}\tau^{1/2}
.\end{equation*}
Combining this with Lemma \ref{lem:num_of_samples_in_a_ball_clean} we get that for any $\delta$ there is $N$ large enough such that for all $n>N$  with probability of at least $1-\delta$
\begin{align*}
R_1'(r; p, H) 
&\geq 
\frac{\frac{n}{3}\cdot \mu_{max}\cdot \mathrm{V}_d \cdot(a_2\sqrt{\sigma\tau})^d}{3n\cdot\mu_{min}\cdot V\cdot ( \sqrt{\sigma\tau}+\sigma)^d}\left(\frac{1}{64}\sigma\tau sin\alpha_0 -  \frac{1}{2}\sigma^{3/2}\tau^{1/2} \right)
-9\sigma^2 \\
&=
\frac{ a_2^d\cdot\mu_{max}}{9\cdot\mu_{min}}\left(\frac{ 1}{1+\sqrt{\sigma/\tau}}\right)^d\left(\frac{1}{64}\sigma\tau sin\alpha_0 - \frac{1}{2}\sigma^{3/2}\tau^{1/2} \right)
-9 \sigma^2
.\end{align*}
Since $1/(1+x) \geq 1-x$ for sufficiently small $x$ we have that, for large enough $M$ of Assumption \ref{assume:noise}, of section \ref{sec:SamplingAssumptions},
\[
\left(\frac{ 1}{1+\sqrt{\sigma/\tau}}\right)^d \leq (1-\sqrt{\sigma/\tau})^d
\]
holds.
Thus, we have 
\[
R_1'(r; p, H) \geq \frac{ a_2^d\cdot\mu_{max}}{9\cdot\mu_{min}}(1-\sqrt{\sigma/\tau})^d\left(\frac{1}{64}\sigma\tau sin\alpha_0 - \frac{1}{2}\sigma^{3/2}\tau^{1/2} \right)
-9 \sigma^2,
\]
Since $a_2 = 1/8$ we have 
\[
R_1'(r; p, H) \geq \frac{ \mu_{max}}{\mu_{min}}\cdot\frac{1}{9}\left(\frac{1-\sqrt{\sigma/\tau}}{ 8}\right)^d \left(\frac{1}{64}\sigma\tau sin\alpha_0 - \frac{1}{2}\sigma^{3/2}\tau^{1/2} \right)
-9 \sigma^2,
\]
as required.
\qed

\begin{figure}
    \centering
    \includegraphics[width=0.8\textwidth]{Figures/Manifold_Good_disc4.pdf}
    \caption{Assisting illustration for the proof of Claim \ref{thm:claim_ri_tyk_bound_clean}: the sphere section \emph{(}in purple\emph{)} represents the manifold; The $xy$-plane represents $T_p\MM$; the plane $\LL_{\tilde j}$ is spanned by $y_{\tilde j}\in T_p\MM^\perp$ and $u_{\tilde j}\in T_p\MM$; the blue disc on $T_p\MM$ is $B_{T_p\MM}(p, \sqrt{\sigma\tau} - \sigma)$; the manifold is considered to be the graph of the function $\phi_p:B_{T_p\MM}(p, \sqrt{\sigma\tau}- \sigma)\to T_p\MM^\perp$. }
    \label{fig:ManifoldGoodDisc_clean}
\end{figure}
\paragraph*{Proof of Claim \ref{thm:claim_ri_tyk_bound_clean}:}
As a result of Corollary \ref{cor:GraphOfFunctionTau_ROI}, $r_i \in U_\textrm{ROI}$ can be written as 
$$
r_i = \underbrace{p+(x_i, \phi_p(x_i))}_{p_i}+\eps_i ,
$$
where $x_i=P_{T_p\MM}(r_i - p)\in T_{p}\MM$ and $\|\eps_i\| \leq \sigma$.
Since 
$$ 
B_D(p, \sqrt{\sigma\tau}-\sigma ) \subset B_D(r, \sqrt{\sigma\tau}) = U_\textrm{ROI}
$$
(see the blue disc on $T_p\MM$ in Figure \ref{fig:ManifoldGoodDisc_clean}), we look for a point $x_0\in T_p\MM\cap \LL_{\tilde j}$ and a radius $\rho >0$ such that $\Gamma_{\phi_p, B_{T_p\MM}(x_0, \rho)}^\sigma\subset B_D(p, \sqrt{\sigma\tau} - \sigma)$, where
\begin{equation}
    \Gamma_{\phi_p, B_{T_p\MM}(x_0, \rho)}^\sigma \defeq \{x ~|~ dist(x, \Gamma_{\phi_p, B_{T_p\MM}(x_0, \rho)}) < \sigma\}
,\end{equation}
where we remind that
\[
\Gamma_{\phi_p, B_{T_p\MM}(x_0, \rho)} = \{p + (x, \phi_p(x))_{T_p\MM} ~|~ x\in B_{T_p\MM}(x_0, \rho)\}
.\]
Furthermore, we choose $x_0$ and $\rho$, such that $\lrangle{r_i - p, \tilde y_{\tilde j}}^2$ is large for any point $r_i\in\Gamma_{\phi_p, B_{T_p\MM}(x_0, \rho)}^\sigma$ (see Figure \ref{fig:ManifoldGoodDisc_clean} for an illustration).

For convenience, we denote by $q^y, q^x, q^{\vec x}$ the projections of $q\in \Gamma_{\phi_p, B_{T_p\MM}(x_0, \rho)}^\sigma$ onto $y_{\tilde j}$ (i.e., $T_p\MM^\perp\cap \LL_{\tilde j}$), $u_{\tilde j}$ (i.e., $T_p\MM\cap \LL_{\tilde j}$) and $T_p\MM$ respectively.
Let $a_1, a_2<1$ and we define 
$$\rho = a_2 \cdot\sqrt{\sigma\tau},$$
and 
$$x_0 = (a_1\sqrt{\sigma\tau})\cdot u_{\tilde j}\in T_p\MM \cap \LL_{\tilde j}.$$
That is, 
\[
\norm{p - x_0} = a_1 \cdot \sqrt{\sigma\tau}
.\]
In order to make sure that $\Gamma^\sigma_{\phi_p, B_{T_p\MM}(x_0, \rho)}\subset U_{\textrm{ROI}}$ we restrict the choice of $a_1, a_2$ such that
\begin{equation}\label{eq:a1a2_restriction_clean}
    \forall x\in B_{T_p\MM}(x_0, a_2\sqrt{\sigma\tau}): \norm{p+(x,\phi_p(x))_{T_p\MM} ~-~ p}
    + \sigma < \sqrt{\sigma\tau} - \sigma
,\end{equation}
We wish to reiterate that $U_\textrm{ROI} = \{r_i ~|~ \norm{r - r_i}\leq \sqrt{\sigma \tau}\}$ and so $U_\textrm{ROI}\subset B_D(r, \sqrt{\sigma\tau})$.
Furthermore, since $\norm{r - p}\leq \sigma$, we have $B_D(p, \sqrt{\sigma\tau} - \sigma)\subset  B_D(r, \sqrt{\sigma\tau})$.
Accordingly, all points $p+(x, \phi_p(x))_{T_p\MM}$ for  $x\in B_{T_p\MM}(x_0, a_2 \sqrt{\sigma\tau})$ are within our region of interest even when moved $\sigma$ away from the manifold in some direction into $\MM_\sigma$ (in Fig. \ref{fig:ManifoldGoodDisc_clean}  $B_{T_p\MM}(x_0, a_2\sqrt{\sigma\tau})$ is the orange disc and the projections of $p+(x, \phi_p(x))_{T_p\MM} + \eps(x)$ onto $T_p\MM$ are limited by the green disc containing the orange disc).
Using the calculations in Appendix \ref{subsec:simplification_app} we can use a simpler demand using $a_1$ and $a_2$, which ensures that the inequality \eqref{eq:a1a2_restriction_clean} is satisfied. The simplified requirement is
\begin{equation}\label{eq:a1a2_simplification_clean}
    (a_1+a_2)<\frac{1}{\sqrt{2}} -\frac{\sqrt{2\sigma}}{\sqrt{\tau}}
.\end{equation}

Let us now bound the value of $\lrangle{\tilde q - p, \tilde y_{\tilde j}}^2$ from below, for any $\tilde  q\in \Gamma^\sigma_{\phi_p,B_{T_p\MM}(x_0, a_2\sqrt{\sigma\tau})}$.
Every such $\tilde  q$ can be written as
\[
\tilde q = q + \eps
,\]
where $q\in \Gamma^0_{\phi_p,B_{T_p\MM}(x_0, a_2\sqrt{\sigma\tau})}$, and so 
 \begin{equation}\label{eq:qxBound_clean}
 (a_1 - a_2)\sqrt{\sigma\tau}  \leq \abs{q^x} \leq \|q^{\vec{x}}\| \leq \|p-x_0\|+a_2\sqrt{\sigma\tau} = (a_1 + a_2)\sqrt{\sigma\tau}   
.\end{equation}
From Lemma \ref{lem:f_bound_circle_clean} we get that
\[
\abs{\lrangle{q, y_{\tilde j}}} \leq \tau-\sqrt{\tau^2-\norm{q^{\vec{x}}}^2} 
.\]
Thus, 
\[
\abs{\langle \tilde{q}, y_{\tilde j} \rangle} = |\tilde{q}^y| \leq \tau+\sigma-\sqrt{\tau^2-\|q^{\vec{x}}\|^2}
,\]
and by plugging the right hand side of \eqref{eq:qxBound_clean} we get
\begin{align*}
\langle \tilde{q}, y_{\tilde j} \rangle^2 & \leq \left(\tau+ \sigma -\sqrt{\tau^2-(a_1+a_2)^2\sigma\tau}\right)^2
\\& =\tau^2 \left(1+ \sigma/\tau -\sqrt{1-(a_1+a_2)^2\sigma/\tau}\right)^2
\\&\leq \tau^2(1+\sigma/\tau -(1-(a_1+a_2)^2\sigma/\tau))^2
\\&= \tau^2 (\sigma/\tau +  (a_1+a_2)^2\sigma/\tau)^2
\\&=  \sigma^2(1 +  (a_1+a_2)^2)^2
\\&\leq 2\sigma^2
,\end{align*}
where the last inequality comes from \eqref{eq:a1a2_simplification_clean}.
Since $\angle(y_{\tilde j}, \tilde y_{\tilde j}) =\beta_{\tilde j} $ we can use the Euclidean geometry on $\LL_{\tilde j}$ (Figure \ref{fig:Lk_clean}) to get
\[
\lrangle{q, \tilde y_{\tilde j}} = \lrangle{q, \textrm{Rot}(\beta_{\tilde j}) y_{\tilde j}} = \lrangle{\textrm{Rot}(-\beta_{\tilde j}) q,  y_{\tilde j}}
,\]
where $\textrm{Rot}(\theta)$ denotes the rotation matrix in $\RR^D$ with respect to the angle $\theta$ in $\LL_{\tilde j}$ .
Therefore,
$$|\langle q, \tilde{y}_{\tilde j} \rangle| \geq \abs{-|q^x|\sin \beta_{\tilde j}  + |q^y|\cos \beta_{\tilde j}} = \abs{|q^x|\sin \beta_{\tilde j}  - |q^y|\cos \beta_{\tilde j}}.$$
Using Lemma \ref{lem:f_bound_circle_clean} and \eqref{eq:qxBound_clean} as before we get
\begin{align*}
 |\langle q, \tilde{y}_{\tilde j} \rangle| &=\abs{  |q^x|\sin \beta_{\tilde j}  -  \left(\tau-\sqrt{\tau^2-\|q^{\vec{x}}\|^2} \right)\cos \beta_{\tilde j}}
 \\
 &\geq\left((a_1-a_2)\sqrt{\sigma\tau}\sin \beta_{\tilde j}  -  \tau\left(1-\sqrt{1- (a_1+a_2)^2\sigma/\tau} \right)\cos \beta_{\tilde j} \right)
 \\
 &\geq\tau\left((a_1-a_2)\sqrt{\sigma/\tau}\sin \beta_{\tilde j}-\cos \beta_{\tilde j} + \cos \beta_{\tilde j}(1- (a_1+a_2)^2\sigma /\tau) \right)
 \\
 &= \tau \left((a_1-a_2)\sqrt{\sigma/\tau}\sin \beta_{\tilde j} -(a_1+a_2)^2\sigma/\tau \cos\beta_{\tilde j}\right)
 \\
 &\geq \tau \left((a_1-a_2)\sqrt{\sigma/\tau}\sin \beta_{\tilde j} -(a_1+a_2)^2\sigma/\tau \right) 
,\end{align*}
and
\begin{equation}\label{eq:q_yk_prod_clean}
    |\langle q, \tilde{y}_{\tilde j} \rangle| \geq \tau \left((a_1-a_2)\sqrt{\sigma/\tau}\sin \alpha_0 -(a_1+a_2)^2\sigma/\tau \right) 
.\end{equation}
Since 
$\|\tilde{q} - q\| \leq \sigma $
we get 
$$\abs{\lrangle{\tilde q - q, \tilde y_{\tilde j}}} = |\langle \tilde{q}, \tilde{y}_{\tilde j}\rangle - \langle q, \tilde{y}_{\tilde j}\rangle| \leq\sigma,$$
and so,
\begin{align*}
|\langle q, \tilde{y}_{\tilde j}\rangle| - |\langle \tilde{q}, \tilde{y}_{\tilde j}\rangle|&\leq |\langle q, \tilde{y}_{\tilde j}\rangle - \langle \tilde{q}, \tilde{y}_{\tilde j}\rangle| \leq  \sigma\\
|\langle \tilde{q}, \tilde{y}_{\tilde j}\rangle|&\geq |\langle q, \tilde{y}_{\tilde j}\rangle| - \sigma\\
\langle \tilde{q}, \tilde{y}_{\tilde j}\rangle^2&\geq \langle q, \tilde{y}_{\tilde j}\rangle^2 - 2\sigma|\langle q, \tilde{y}_{\tilde j}\rangle| + \sigma^2\\
\langle \tilde{q}, \tilde{y}_{\tilde j}\rangle^2&\geq \langle q, \tilde{y}_{\tilde j}\rangle^2 - 2\sigma|\langle q, \tilde{y}_{\tilde j}\rangle|
.\end{align*}
 Substituting $\abs{\lrangle{q, \tilde y_{\tilde j}}}$ with the bound from \eqref{eq:q_yk_prod_clean} we get
\begin{align*}
    \lrangle{\tilde q, \tilde y_{\tilde j}}^2 & \geq \tau^2 \left((a_1-a_2)\sqrt{\sigma/\tau}\sin \alpha_0 -  (a_1+a_2)^2{\sigma/\tau}\right)^2 - 2\sigma\tau \left((a_1-a_2)\sqrt{\sigma/\tau}\sin \alpha_0 -  (a_1+a_2)^2{\sigma/\tau}\right) \\
    &\geq \tau^2\left((a_1-a_2)\sqrt{\sigma/\tau}\sin \alpha_0 -  {\sigma/\tau}\right)^2 - 2\sigma\tau \left((a_1-a_2)\sqrt{\sigma/\tau}\sin \alpha_0\right)\\
    &= \tau^2 \left((a_1-a_2)^2{\sigma/\tau}\sin^2 \alpha_0 - 2(a_1-a_2)(\sigma/\tau)^{3/2}\sin\alpha_0 + (\sigma/\tau)^2\right) -2(a_1-a_2)\sigma^{3/2}\sqrt{\tau} \sin \alpha_0 \\
    &\geq (a_1-a_2)^2{\sigma\tau}\sin^2 \alpha_0 - 2(a_1-a_2)\sigma^{3/2}\tau^{1/2}\sin\alpha_0 + \sigma^2 -2(a_1-a_2)\sigma^{3/2}{\tau}^{1/2} \\
    &= (a_1-a_2)^2{\sigma\tau}\sin^2 \alpha_0 - 2(a_1-a_2)\sigma^{3/2}\tau^{1/2} + \sigma^2 -2(a_1-a_2)\sigma^{3/2}{\tau}^{1/2} \\
    &= (a_1-a_2)^2{\sigma\tau}\sin^2 \alpha_0 - 4(a_1-a_2)\sigma^{3/2}\tau^{1/2} + \sigma^2 
\end{align*}
Thus, by choosing for example $a_1 = \frac{1}{4}$ and $a_2 = \frac{1}{8}$ we get
\begin{equation}
    \lrangle{\tilde q, \tilde y_{\tilde j}}^2 \geq \frac{1}{64}\sigma \tau sin\alpha_0 - \frac{1}{2}\sigma^{3/2}\tau^{1/2}
.\end{equation}

\end{proof}


\begin{Lemma}\label{lem:J_1qH_clean}
Let the sampling assumption of Section \ref{sec:SamplingAssumptions} hold. Let ${p_r} = P_\MM(r)$ be the projection of $r$ onto $\MM$, and $ T_{p_r}\MM $ be the tangent to $ \MM $ at $ {p_r} $.
For $\alpha = \sqrt{C_M/M}$, where $M = \frac{\tau}{\sigma}$, and $C_M$ is a constant  the following holds:
For any $\delta >0$ there is $N_\delta$ (independent of $\alpha$) such that $\forall n > N_\delta$ all linear sub-spaces $H \in Gr(d,D)$ with $\angle_{max}(H,T_{{p_r}}\MM) > \alpha$, and all $q$ in the search space of \eqref{eq:argmin1_clean} yield a score 
\[
J_1(r; {q}, H) \geq 100\sigma^2 
\]
with probability of at least $1 - \delta$. 
\end{Lemma}
\begin{proof}
From Lemma \ref{thm:J1pH_clean}, we have that there is a constant $C_M$ such that for any $ {\alpha} < \frac{\pi}{2} $ and $M = C_M/\alpha^2$, and for any $\tau$ and $\sigma$ maintaining $\frac{\tau}{\sigma}>M$
    the following hold:
For any $\delta >0$ arbitrarily small there is $ N_\delta $ sufficiently large such that for all $ n>N_\delta $, \textbf{all} linear spaces $H$ with $\angle_{max}(H,T_{p_r}\MM) > \alpha$, yield a score $J_1(r; p_r, H) \geq 109 \cdot \sigma^2 $, with probability of at least $1 - \delta$. 
We now wish to show that $J_1(r; q, H) \geq 100 \sigma^2$ with high probability as well.
For convenience we wish to reiterate \eqref{eq:J1_clean}
\[
J_1(r; q, H) = \frac{1}{n}\sum_{r_i\in U_{\textrm{ROI}}} \dist^2 (r_i - q,H)
.\]
By Constraint 3 of \eqref{eq:J1_clean} we achieve that $dist(q, \MM) \leq 3\sigma$ and so for $p=P_\MM(q)$ we have
\[
\norm{q - p} \leq 3\sigma
.\]
Thus,
\[
\dist^2 (r_i - q,H) = \norm{p - q}^2 + \dist^2 (r_i - p,H) \geq \dist^2 (r_i - p,H) - 9\sigma^2
,\]
and we achieve
\[
J_1(r; q, H) \geq J_1(r; p, H) - 9\sigma^2
.\]
By Lemma \ref{thm:J1pH_clean} we conclude the proof of the current lemma.
\end{proof}
% --- end Section_ProofStep1.tex ---


\subsection{Proof of Theorem \ref{thm:Step2}}
\label{sec:Step2Analysis}
% --- begin Step2_lemma_map.tex ---



\tikzstyle{decision} = [diamond, aspect=2, draw, fill=yellow!0,
text width=7em, text badly centered, node distance=3cm, inner sep=0pt]
\tikzstyle{block} = [rectangle, draw, fill=blue!20,
text width=5em, text centered, rounded corners, minimum height=2em]
\tikzstyle{wblock} = [rectangle, draw, fill=blue!20,
text width=7em, text centered, rounded corners, minimum height=2em]
\tikzstyle{line} = [draw, -latex']
\tikzstyle{bbblock} = [rectangle, draw, fill=blue!20,
text width=10em, text centered, rounded corners, minimum height=2em]
\tikzstyle{xbblock} = [rectangle, draw, fill=blue!20,
text centered, rounded corners, minimum height=2em]
\tikzstyle{wwblock} = [rectangle, draw, fill=blue!20,
text width=9.5em, text centered, rounded corners, minimum height=2em]
\tikzstyle{Iwwblock} = [rectangle, draw, fill=blue!50,
text width=9.5em, text centered, rounded corners, minimum height=2em]
\tikzstyle{line} = [draw, -latex']


\tikzstyle{Iblock} = [rectangle, draw, fill=blue!40,
text width=5em, text centered, rounded corners, minimum height=2em]


\tikzstyle{Iwblock} = [rectangle, draw, fill=blue!40,
text width=7em, text centered, rounded corners, minimum height=2em]

\tikzstyle{Rblock} = [rectangle, draw, fill=red!20,
text width=5em, text centered, rounded corners, minimum height=2em]

\tikzstyle{Rbbblock} = [rectangle, draw, fill=red!20,
text width=10em, text centered, rounded corners, minimum height=2em]






\newcommand{\dd}{2cm}
\newcommand{\ddd}{4cm}
\newcommand{\thd}{4cm}
\newcommand{\qut}{3cm}
\newcommand{\oct}{2cm}
\newcommand{\sm}{0.8cm}
\newcommand{\wit}{20pt}


\tikzstyle{innerWhite} = [-, white,line width=1mm, shorten >= 4.5pt]
\tikzstyle{innerWhiteEq} = [-, white,line width=1mm, shorten >= 0pt]
\tikzstyle{imp} = [-{implies}, double, double distance=1.1mm, line width=0.3mm]
\tikzstyle{Rimp} = [-{implies}, red,double, double distance=1.1mm, line width=0.3mm]
\tikzstyle{Gimp} = [-{implies}, green,double, double distance=1.1mm, line width=0.3mm]
\tikzstyle{Req} = [=, red,double, double distance=1.1mm, line width=0.3mm]
\tikzstyle{eq} = [=,  double distance=1.1mm, line width=0.4mm]
\tikzstyle{Geq} = [=, green, double distance=1.1mm, line width=0.4mm]
\usetikzlibrary{arrows, shapes.gates.logic.US, calc}

\begin{figure}[omitted: a diagram of this paper's results, not reproduced here]\label{tikz:thm33_lemmas}\end{figure}






% \begin{figure}
%     \centering
%     \includegraphics[width=\textwidth]{Figures/Step2_lemma_map.pdf}
%     \caption{Road-map for proof of Theorem \ref{thm:Step2}}
%     \label{tikz:thm33_lemmas}
% \end{figure}
% --- end Step2_lemma_map.tex ---
% --- begin Section_ProofStep2.tex ---

Before delving into the details of the proof, we wish to reiterate the steps of Algorithm \ref{alg:step2_clean} while introducing some useful notations.
According to the assumptions of Theorem \ref{thm:Step2} we have a local coordinate system $(q, H)\in \RR^D\times Gr(d, D)$, such that $ \norm{q-p} \leq 3 \sigma $ and $\angle_\textrm{max}(T_p\MM, H) \leq \alpha_0$, where $p\in \MM$ .
As Algorithm \ref{alg:step2_clean} involves an iterative process, we denote the initial values as $q_{-1} := q$ and $H_0 := H$.
Then, at each iteration we update the origin $q_\ell$ and  the local coordinates' directions $H_\ell$ (for $\ell=1, \ldots, \kappa$). 

Similar to \eqref{eq:FunctionGraph_init} and using the result of Lemma \ref{lem:M_is_locally_a_fuinction_clean}, we begin by looking at the manifold patch $\MM \cap \textrm{Cyl}_{H_0}(p,c_{\pi/4}\tau,\tau/2)$ as the graph of a function $f_{-1/2}:(q_{-1},H_0) \simeq\RR^d\to\RR^{D-d}$; i.e., 
\begin{equation}
 \Gamma_{f_{-1/2}, q_{-1},H_0} \defeq \{q_{-1} + (x, f_{-1}(x))_{H_0} ~|~ x\in B_{H_0}(0, c_{\pi/4}\tau)\}  
,\end{equation}
where $\textrm{Cyl}_{H_0}(p,c_{\pi/4}\tau, \tau/2)$ is the $D$-dimensional cylinder with the base $B_H(p, c_{\pi/4}\tau)\subset H_0$ and height $\tau/2$ in any direction on $H_0^\perp$. 
For the remainder of this section we assume that $ \alpha $ is small enough.
Then, at the first step we estimate $f_{-1/2}(0)$ in order to update the origin from $q_{-1}$ to $q_0$ and make sure it is close to the manifold. 
This is done through the weighted least-squares minimization of \eqref{eq:argmin2_clean} and by evaluating the local polynomial estimate $\pi^*_{q_{-1}, H_0}$ at $ 0 $; i.e., we set $q_0 = q_{-1} + (0,\pi^*_{q_{-1}, H_0}(0))_{H_0}$.

Subsequently, we look at the function $f_0:(q_0,H_0) \simeq \RR^d\to\RR^{D-d}$ taking us from the shifted coordinate system $(q_0, H_0)$ to $\MM$; i.e., $\MM$ is now locally expressed by the graph 
\begin{equation}
 \Gamma_{f_0, q_0, H_0} \defeq \{q_{0} + (x, f_0(x))_{H_0} ~|~ x\in B_{H_0}(0,c_{\pi/4}\tau)\}
.\end{equation}

Using similar notation, each iteration of Algorithm \ref{alg:step2_clean} comprises two steps. First, we update $H_\ell$ to $H_{\ell + 1}$ by taking the linear space coinciding with the directions of the estimated tangent to $f_{\ell }$ at zero to get $f_{\ell+ 1/2}:(q_\ell, H_{\ell+1}) \simeq \RR^d \to \RR^{D-d}$. Second, we update $q_\ell$ to $q_{\ell+1} $ by taking the estimated value of $f_{\ell+1/2}$ at zero to get $f_{\ell+1}:(q_{\ell+1},H_{\ell+1}) \simeq \RR^d \to \RR^{D-d}$.

Explicitly, given $(q_{\ell}, H_{\ell})$, we look at the manifold as $\Gamma_{f_{\ell}, q_{\ell}, H_{\ell}}$ the local graph of a function $f_{\ell}$ and estimate $T_0 f_{\ell} \in Gr(d, D)$ the tangent to the graph of $f_{\ell}$ at $0$ through taking the image of $\DD_{\pi^*_{q_{\ell}, H_{\ell}}}[0]$, the first order differential of $\pi^*_{q_{\ell }, H_{\ell }}$ at zero.
In other words, we ``rotate" the coordinate system to the point where $H_{\ell+1}$ aligns with the former tangent estimation to get $f_{\ell+1/2}$. 
Then, we define 
 $ q_{\ell+1}  =  q_\ell + ( 0 , \pi^*_{q_\ell, H_{\ell+1}} )_{H_\ell} $ and get $f_{\ell +1}$ by shifting the coordinate system from $(q_\ell, H_{\ell+1})$ to $(q_{\ell+1}, H_{\ell+1})$.

However, if we want to use the well-known convergence rates of local polynomial regression (i.e., the minimization of \eqref{eq:argmin2_clean}), a key assumption in the analysis is that the noise is of zero mean \cite{stone1980optimal}. 
However, in our case, for any $x\in H$, the samples above $x$ are uniformly distributed~in 
\begin{equation}\label{eq:Omega_def}
\Omega(x) = (x + H^\perp) \cap  \MM_\sigma
,\end{equation}
where $x + H^\perp = \{x + y ~|~ y\in H^\perp\}$.
That is, $\Omega(x):H \simeq \RR^d \to 2^{\RR^{D-d}}$.
Thus, denoting $\eta(x)\sim \textrm{Unif}(\Omega(x))$ and defining 
\begin{equation}\label{eq:ftilde_def_general}
\wtilde f(x) = \mathbb{E}[\eta(y | x)]    
,\end{equation}
the result of the regression will estimate $\wtilde f$ rather than $f$ itself.

In our case, we are estimating $f_{\ell}$ and $f_{\ell + 1/2}$, and the corresponding $\eta_{\ell}(y | x), \eta_{\ell+1/2}(y | x)$ are compactly supported since the sampling is uniform on $\MM_\sigma$.
However, that same fact implies that $f_\ell \neq \mathbb{E}[\eta_\ell(y | x)]$ and $f_{\ell+1/2} \neq \mathbb{E}[\eta_{\ell+1/2}(y | x)]$ (see Figure \ref{fig:TiltedBias}).
Thus, when we try estimating $f_\ell$ or its derivatives, we in fact estimate a different function, for which the noise has zero mean (see Figure \ref{fig:bias_ex_clean}); that is, we are estimating the function of conditioned expected value 
\begin{equation}\label{eq:ftilde_def}
\wtilde{f}_\ell(x) = \mathbb{E}[\eta_\ell(y |x)]
,\end{equation}
where $\eta_\ell(y |x)$ is the distribution of values $y\in H_{\ell}^\perp$ conditioned by the values $ x\in H_{\ell} $.

\begin{figure}
	\centering
	\includegraphics[width=0.4\textwidth]{Figures/Tilted3.pdf}
	\caption{Illustration of $ \MM $ as a graph of a function $ f_\ell $ \emphbrackets{marked by the red solid line} above the coordinate system $ f_\ell $. The boundary of  $ \MM_\sigma $ is delineated by the pink lines and $ \wtilde f_\ell $ is the conditioned expectancy $ \EE[\eta_\ell(y ~|~ x)] $ of this domain with respect to the presented $ y $-axis. }
	\label{fig:bias_ex_clean}
\end{figure}



Since for any arbitrarily fixed $ x $ the density of the random variable $ \eta_\ell(y | x) $ is constant, $\wtilde{f}_\ell(x)$ can be computed as the mean of the set $\Omega(x)$ defined in \eqref{eq:Omega_def}.
Furthermore, the farthest point in each such direction is exactly $\sigma$ away from the graph of $f_\ell$ (see the red line in Figure \ref{fig:draw_step_2_bias_clean}).

\begin{figure}
	\centering
	\includegraphics[width=0.45\textwidth]{Figures/BoundingBias.pdf}
	\caption{Illustration of $g(0,\vec{y})$ and $\Omega(0)$ in the two dimensional case. Let $H$ be some local coordinate system and consider $\MM$ as a local graph of some function $f:H\to H^\perp$. The upper bound for the values of the sample distribution above $0$ in some direction $\theta$ of $H^\perp$ is $g(0,\theta)$. This value is $\sigma$-away from some point on $\MM$. We denote this point by $f(\wtilde x(0))\in \MM$.}
	\label{fig:draw_step_2_bias_clean}
\end{figure}
Below we show that Algorithm \ref{alg:step2_clean} starts with $(q, H)$ a rough estimate of the origin and  tangent (which plays the role of the coordinate system), and as $n\to \infty$, $\hat p_n$ approaches $\mathbf p\in \MM$ and $\widehat{T_{\hat p_n}\MM}$ approaches $T_{\mathbf p}\MM$ as well.  
We first note that by Lemma \ref{lem:Tf_Tftilde_Init} since $ \maxangle(H_1, T_p\MM) \leq \alpha $ is sufficiently small we get 
\[
\maxangle(H_0, T_0 f_{-1/2}) \leq \frac{3\alpha}{2}
,\]
and we denote $ \alpha_0 = \maxangle(H_0, T_0 f_{-1/2}) = \maxangle(H_0, T_0 f_{0}) $ with which we initiate the iterates.

A key lemma in the proof of Theorem \ref{thm:Step2} is the following

\begin{Lemma} \label{lem:main_support_theorem_step2}
    Let $H_\ell\in Gr(d, D)$, and let $f_\ell:H_\ell\simeq\RR^d \rightarrow \RR^{D-d}$, defined as in \eqref{eq:fl_def}. Define $H_{\ell+1} = \Ima(\DD_{\pi^*_{q_\ell,H_{\ell}}}[0])$, as in Algorithm \ref{alg:step2_clean}. Let $f_{\ell+1/2}$ be as defined in \eqref{eq:fl12_def} and let $r_1 = \frac{k-1}{2k +d}$. Assume that $M = \frac{\tau}{\sigma} \geq C_\tau \sqrt{D\log D}$ (where $C_\tau$ is a constant from Lemma \ref{lem:Tf_TtildeF_D}).
    For any $\delta>0$, there is $N$ such that for any number of samples $n>N$, and any $\alpha \leq\frac{3}{2} \sqrt{C_M/M}$ smaller than some constant (see Theorem \ref{thm:Step1} and Section \ref{sec:SamplingAssumptions} for the definition of $M$ and $C_M$).  If
    \begin{enumerate}
        \item $\maxangle (H_\ell, T_0 f_\ell) = \alpha_\ell\leq \alpha$
        \item $48\frac{\|f_\ell(0)\|}{\tau} \leq \alpha$
        \item $ 12\sqrt{d} \frac{C_0 \ln(1/\delta)}{n^{r_1}}\leq \alpha $ 
    \end{enumerate}
    hold, where $C_0$ is some constant. Then, we have  
    \[
    \alpha_{\ell+1} = \maxangle (H_{\ell+1}, T_0 f_{\ell+1})  = \maxangle (H_{\ell+1}, T_0 f_{\ell+1/2}) \leq \alpha/2
    \]
    with probability at least $1-\delta$.
\end{Lemma}


Lemma \ref{lem:main_support_theorem_step2} is proved in Section \ref{subsubsec:proof of main support lemma}.
Now we prove Theorem~\ref{thm:Step2}.
\begin{proof}[proof of theorem \ref{thm:Step2}]
We divide Algorithm \ref{alg:step2_clean} into three steps:
\begin{enumerate}
    \item[(i)] Initialize $q_0$ and $H_0$ - corresponds to rows 3 and 4 in Algorithm \ref{alg:step2_clean}.
    \item[(ii)] Estimate $H_{\ell+1}$ from $H_{\ell}$ and $q_\ell$ - corresponds to rows 6 to 8 in Algorithm \ref{alg:step2_clean}.
    \item[(iii)] Estimate $q_{\ell+1}$ from $H_{\ell+1}$ - corresponds to row 10 in Algorithm \ref{alg:step2_clean}.
\end{enumerate}
Where steps (ii) and (iii) are repeated $\kappa$ times.
We will treat the three steps one by one, and prove that the "output" of each step will fit our requirements for the input of the next step. Although we will prove this point later, we start with assuming that $\kappa$ and $\delta_1$ are such that 
\begin{equation}\label{eq:n_bound1}
    12\sqrt{d}\frac{C_0 \ln(1/\delta_1)}{n^{r_1}} \leq \alpha_0 2^{-\kappa+1},
\end{equation}
where $C_0$ is the constant from Lemma \ref{lem:main_support_theorem_step2}. 

We start with \textbf{step (i)}: From Lemma \ref{lem:Tf_Tftilde_Init} there is $N_{1,\delta_1}$ such that for all $n>N_{1,\delta_1}$ we have that:
\[
\Pr\left(\maxangle (T_0f_{0}, H_{0}) \leq \alpha_1 \right) \geq 1-\delta_1,
\]
where $\alpha_1 = \frac{3}{2} \alpha_0$.
From Lemma \ref{lem:dist_to_f_l0_weak} Since $\alpha <1/D^{1/4}$ there is $N_{2,\delta_1}$ such that for all $n>N_{2,\delta_1}$ we have that:
\[
\Pr\left(\|f_0(0)\| \leq \tau \alpha_1 /48\right) \geq 1-\delta_1
\]
hold.

Denote by $A_\ell$ the event that $\maxangle (T_0f_{\ell}, H_{\ell}) \leq \alpha_1 2^{-\ell}$ and by $B_\ell$ the event that $\|f_\ell(0)\| \leq \tau \alpha_1 2^{-\ell} /48$, and then, by the union bound, we have that $\Pr(A_0 \mbox{ and } B_0) \geq 1-2\delta_1$.

Next, considering \textbf{step(ii)}. Since $M>C_\tau \sqrt{D\log D}$, and $\alpha_1 \leq \frac{3}{2} \sqrt{ C_M/M}$ if events $A_\ell$ and $B_\ell$ hold, then, the requirements of Lemma \ref{lem:main_support_theorem_step2} are met for $\alpha = \alpha_1 2^{-\ell}$. Since $\maxangle (T_0f_{\ell+1/2}, H_{\ell+1}) =\maxangle (T_0f_{\ell+1}, H_{\ell+1}) $, we have that for $n> N_{3,\delta_1}$, with probability of at least $1-\delta_1$, the event $A_{\ell+1}$ holds, i.e 
\[
\Pr(A_{\ell+1}|A_\ell, B_\ell) \geq 1-\delta_1.
\]

For \textbf{step(iii)}, we show that given that the event $A_{\ell+1}$ holds, then, with probability of at least $1-\delta_1$, the event $B_{\ell+1}$ holds. From Lemma \ref{lem:dist_to_f_l0_weak} with  $\alpha = \alpha_1 2^{-\ell}$ there is $N_{4,\delta_1}$ such that for all $n>N_{4,\delta_1}$ we have that:
\[
\Pr\left(\|f_{\ell+1}(0)\| \leq \tau \alpha_1 2^{-(\ell+1)}/48\right) \geq 1-\delta_1
\]
hold. 

To conclude our arguments so far, since $\kappa$ satisfies \eqref{eq:n_bound1}, using the union bound on the events $A_0,\ldots, A_\kappa$ and  $B_0,\ldots, B_\kappa$  we have that 
\begin{equation}\label{eq:final_prob_bound1}
\Pr\left(\maxangle (T_0 f_{\kappa}, H_{\kappa}) \leq \alpha_1 2^{-\kappa}\right) \geq 1 -  2\kappa\delta_1.        
\end{equation}
Choosing $\delta_1 = \frac{\delta}{2\kappa}$, we have from Lemma \ref{lem:comute_kappa}, that $\kappa$ from \eqref{eq:kappa_res_lem} satisfies \eqref{eq:n_bound1_kappa_lem}, and thus $\kappa$ satisfies \eqref{eq:n_bound1} as well. We also have from Lemma \ref{lem:comute_kappa} that 
 \[
    \alpha_1 2^{-\kappa} \leq C_{d} \ln\left(\frac{1}{\delta}\right) n^{-r_1}\left(\ln\left(\ln(n) \right)\right)^{2r_1}.
    \]
Thus, we have that 
\begin{equation}
\Pr\left(\maxangle (Tf_{\kappa}(0), H_{\kappa}) \leq C_{d}\ln\left(\frac{1}{\delta}\right) n^{-r_1}\left(\ln\left(\ln(n) \right)\right)^{2r_1}\right) \geq 1-\delta.    
\end{equation}
This concludes the proof of Eq. \eqref{eq:thm3.3_tangent_ang}. 


Furthermore, assuming event $A_\kappa$ holds, there is $N_5$ such that for $n>N_{5}$ we have that $\alpha = \alpha_1 2^{-\kappa}<1/D$. Thus, from Lemma \ref{lem:dist_to_f_l0_srtong_small_alpha} we have that  
\[
\|f_\kappa(0)\| \leq DC_1 \alpha_1^2 2^{-2\kappa} + \frac{C_2\ln\left(\frac{1}{\delta}\right)}{n^{r_0}}
\]
holds with probability at least $1-\delta_1$. 
Denote this event by $\bar{B}_\kappa$.
Substituting $\kappa$ from \eqref{eq:kappa_res_lem} we have from Lemma \ref{lem:comute_kappa} that 
\[
\| f_\kappa(0)\| \leq D C_1 C_{d}^2 \ln\left(\frac{1}{\delta}\right)^2 n^{-2r_1}\left(\ln\left(\ln(n) \right)\right)^{4r_1} + C_2\ln\left(\frac{1}{\delta}\right)n^{-r_0}
\]
or, for $n$ large enough, and some constant $C$,
\[
\|f_\kappa(0)\| \leq C \ln\left(\frac{1}{\delta}\right) n^{-r_0}.
\]
%for some different constant $C$.

Thus, using the union bound on $A_1,\ldots, A_\kappa, B_1,\ldots, B_{\kappa-1}$, and $\bar{B}_\kappa$, we have that 
\begin{equation}
\Pr\left(  \|f_\kappa(0)\| \leq C\ln\left(\frac{1}{\delta}\right) n^{-r_0}   \right) \geq 1-\delta.    
\end{equation}
This concludes the proof of Eq. \eqref{eq:thm3.3_p_bound} and of Theorem \ref{thm:Step2}. 



\end{proof}





% --- end Section_ProofStep2.tex ---

%\subsection{Proof of Lemma \ref{lem:main_support_theorem_step2}}
%\label{sec:ProofofLemAlphaHalf}

% \subsubsection{Bounding the error between $T_0 {f_\ell}$ and $T_0\wtilde{f}_\ell$}\label{sec:Tf_Tf_tilde_err}
% \input{Step2_BoundTf_Tf_tilde_lemma_map}
% \input{Section_ProofStep2_BoundTf_Tf_tilde.tex}

% \subsubsection{Bounding the finite sample error}\label{sec:FiniteSample_err}
% \input{Section_ProofStep2_BoundTf_tilde_TF_hat.tex}

% \subsubsection{Bounding the distance of $q_\ell$ from $f_\ell(0)$}\label{sec:f0_est}
% \input{Section_ProofStep2_f0_est.tex}


% \subsubsection{Bounding the error induced by the shifted origin}\label{sec:ShiftedCenter_err}
% \input{Section_ProofStep2_BoundTf_move_origin.tex}


\subsection{Proof of Lemma \ref{lem:main_support_theorem_step2}}\label{subsubsec:proof of main support lemma}
 
% --- begin Section_ProofStep2_proofMainLemma.tex ---

 For road-map of the proof of Theorem \ref{thm:Step2} see Figure \ref{tikz:thm33_lemmas}.
 
\begin{proof}
From Lemma \ref{lem:Tf_TtildeF_D} we have that 
\begin{equation}\label{eq:bound_Tfl_Tftile_l}
\maxangle(T_0 f_\ell, T_0 {\wtilde{f}_\ell}) 
\leq  
\frac{\alpha}{6}.    
\end{equation}

From Lemma \ref{lem:AngleImprovement} we have that for any $\delta$ and $n>N_{\delta}$, 
\[
\Pr(\maxangle(T_0\wtilde f_\ell, T_0\pi^*_{q_0, H_\ell}) \geq 2\sqrt{d}\frac{C_0 \ln(1/\delta)}{n^{r_1}}) < \delta.
\]
Thus, as $ 12\sqrt{d} \frac{C_0 \ln(1/\delta)}{n^{r_1}}\leq \alpha $  (see assumption 3 in the lemma), we have 
\begin{equation}\label{eq:pr_alpha_6}
\Pr(\maxangle(T_0\wtilde f_\ell, T_0\pi^*_{q_0, H_\ell}) \geq \frac{\alpha}{6}) < \delta.
\end{equation}
And, from
\[
\maxangle (T_0f_{\ell}, H_{\ell+1}) \leq \maxangle (T_0f_\ell, T_0\wtilde{f}_\ell)  + \maxangle (T_0\wtilde{f}_\ell, H_{\ell+1}). 
\]
along with \eqref{eq:bound_Tfl_Tftile_l} and \eqref{eq:pr_alpha_6}, we have that 
\begin{equation}\label{eq:pr_alpha_3}
\Pr\left(\maxangle (T_0f_{\ell}, H_{\ell+1}) \geq \alpha/3\right) < \delta.    
\end{equation}

In order to apply Lemma \ref{lem:shift_ang_general} we denote $g_0 = f_\ell$, $g_1 = f_{\ell+1/2}$, $G_0 = H_\ell$, and $G_1 = H_{\ell+1}$.
Thus, we have that $\maxangle (T_0g_0, G_0) = \maxangle (T_0f_\ell, H_\ell)\leq \alpha$ and  $\maxangle(G_0,G_1) = \maxangle(H_\ell,H_{\ell+1})\leq \maxangle(H_\ell,T_0f_\ell) + \maxangle(T_0f_\ell,H_{\ell+1})$.
Under the assumption that the event described in \eqref{eq:pr_alpha_3} holds we also know that $\maxangle(G_0,G_1)  \leq \alpha + \alpha/3 = \beta$.
Lastly, as $\alpha \leq \pi/16$, $\beta = \frac{4\alpha}{3} \leq \alpha_c$, and  $\|f_\ell(0)\| \leq \frac{3\tau}{4 \cdot 16} $ (follows from assumption 2 of the lemma), we can use the result of Lemma  \ref{lem:shift_ang_general} and obtain 
    \[
    \alpha_{\ell+1} = \maxangle (T_0f_{\ell+1}, H_{\ell+1}) = \maxangle (T_0f_{\ell+1/2}, H_{\ell+1}) \leq \maxangle(T_0f_\ell,H_{\ell+1})  + \frac{ 8\|f_\ell(0)\|}{\tau}  .\]
Furthermore, since $\maxangle(T_0f_\ell,H_{\ell+1}) \leq \alpha/3$ and $\frac{ 8\|f_\ell(0)\|}{\tau} \leq \alpha/6$, we have 
\[
\maxangle (T_0f_{\ell+1}, H_\ell)  = \maxangle (T_0f_{\ell+1/2}, H_\ell) \leq \alpha/2
.\]
\end{proof}

% --- end Section_ProofStep2_proofMainLemma.tex ---


\section{A Possible Application}\label{sec:applications}
% --- begin Section_Applications.tex ---
While there are numerous applications for this method, we present here one example that demonstrate the potential of the presented approach.
In this example we show how the presented method can be used to follow the trajectory of a geodesic line on a manifold. 
We assume that a point $x_0$ is chosen on the manifold and some direction $\vec v_0$ on the tangent $T_{x_0} \MM$ (In practice, the process can initialized with a point near the manifold $\MM$ and then project it to the manifold). 

The process of tracking a geodesic line is iterative. 
At each step, we compute $\wtilde x_{i+1} =  x_i + \eps \vec{v}$, then ``project'' the new point back to the manifold $x_{i+1} \approx \mbox{Proj}_\MM \wtilde x_{i+1}$ through Algorithms \ref{alg:step1Inpractice_clean} and \ref{alg:step2InPractice_clean}, and parallel transport $\vec v_i$ to $T_{x_{i+1}}$ to get $\vec v_{i+1}$.

In the first toy case, the manifold $\MM$ is a circle of radius 10 in $\RR^2$. 
The dataset consists of 5000 points. 
We start with some sample (illustrated in red in Figure \ref{fig:geoWalk_circle}), project it onto the circle (in Figure \ref{fig:geoWalk_circle}, the circle is marked in blue, and the projected point in green), and than move the point in some direction, project it again (shown in another green point in Figure \ref{fig:geoWalk_circle}), etc.  
\begin{figure}
	\centering
	\includegraphics[width=0.35\textwidth]{Figures/geoWalk_circle.pdf}
	\caption{Geodesic ``walk" on a circle. The red point is the initial point. The yellow points are the data set. The red point is than projected onto the estimation of the blue circle. Then at each step, a new point is generated along the circle (the black arrows connect the points. The plot illustrates 30 steps.}
	\label{fig:geoWalk_circle}
\end{figure}


In the second example, we took a 3d model of an airplane\footnote{\textcolor{blue}{\href{http://3dmag.org/en/market/download/item/4740/}{http://3dmag.org/en/market/download/item/4740/}}}, rotated it in the z-axis, and took 2000 snapshots. Each snapshot is an image of $290 \times 209$ gray-scale pixels. The input data set consist of the unsorted images, sampled from a one dimensional manifold embedded in $\RR^{60,610}$.
Several such images appear in Figure \ref{fig:geoWalk_airplane}. 
Starting from some image, we create a movie of the rotating airplane. 
The movie can be found in \textcolor{blue}{\href{https://youtu.be/aHYyUvu1Q-8}{https://youtu.be/aHYyUvu1Q-8}}, and the code for generating it can be found in \textcolor{blue}{\href{https://github.com/aizeny/manapprox}{https://github.com/aizeny/manapprox}}

\begin{figure}[H]
	\centering
	\includegraphics[width=0.16\textwidth]{Figures/airplane_ang018.pdf}
	\includegraphics[width=0.16\textwidth]{Figures/airplane_ang3008.pdf}
	\includegraphics[width=0.16\textwidth]{Figures/airplane_ang6015.pdf}
	\includegraphics[width=0.16\textwidth]{Figures/airplane_ang9005.pdf}
	\includegraphics[width=0.16\textwidth]{Figures/airplane_ang12066.pdf}
	\includegraphics[width=0.16\textwidth]{Figures/airplane_ang15002.pdf}
	\\
	\includegraphics[width=0.16\textwidth]{Figures/airplane_ang17991.pdf}
	\includegraphics[width=0.16\textwidth]{Figures/airplane_ang20998.pdf}
	\includegraphics[width=0.16\textwidth]{Figures/airplane_ang23970.pdf}
	\includegraphics[width=0.16\textwidth]{Figures/airplane_ang26995.pdf}
	\includegraphics[width=0.16\textwidth]{Figures/airplane_ang29985.pdf}
	\includegraphics[width=0.16\textwidth]{Figures/airplane_ang32992.pdf}
	\caption{Sample images from a 3d model of an airplane.}
	\label{fig:geoWalk_airplane}
\end{figure}% --- end Section_Applications.tex ---
% --- end main_text.tex ---


\section{Acknowledgments}
We wish to thank Prof. Felix Abramovich for driving us to do this work and for the fruitful discussions.
We also thank Prof. Ingrid Daubechies and Prof. David Levin for various discussions along the road.
B. Sober is supported by Duke University, The Hebrew University of Jerusalem, and the Simons Foundation through Math+X grant 400837.

\bibliography{main}{}
\bibliographystyle{plain}

\begin{appendices}
% --- begin main_appendix.tex ---




\section{Preliminaries}
Before we delve into the proofs, we wish to introduce the concepts of Principal Angles between linear sub-spaces \cite{jordan1875essai,bjorck1973numerical} as well as develop some general results concerning the viewpoint of the manifold as being locally a graph of some function from a local coordinate system.
Both of these topics will play a key role in the proofs below.

In addition, two bounds resulting from the Taylor expansion will be used extensively in our proofs.
Thus, we note them here as the two following remarks:

\begin{remark}\label{rem:taylor_sqrt1-x2_clean}
For $x\in[0,\sqrt{3}/2]$
\begin{equation}\label{eq:Taylor2ndOrder}
    1-1/2x^2 \geq \sqrt{1 - x^2} \geq 1-x^2
\end{equation}
\end{remark}
\begin{remark}\label{rem:taylor2_sqrt1-x2_clean}
For $x\in[0,\sqrt{3}/2]$
\begin{equation}
1-1/2x^2-1/8x^4 \geq \sqrt{1 - x^2} \geq 1-\frac{1}{2}x^2 - x^4    
\end{equation}
\end{remark}

\subsection{Principal angles between linear sub-Spaces}
The concept of Principal Angles between flats were first introduced by Jordan in 1875 \cite{jordan1875essai}.
Below, we use the definition of Principal Angles between subspaces as described in \cite{bjorck1973numerical}.
\begin{definition}[Principal Angles]\label{def:principal_angles_clean}
	Let $V$ be an inner product space. Given two sub-spaces $\UU, \WW$ of dimensions $dim(\UU)=k, dim(\WW) = l$, where $k \leq l$ there exists a sequence of $k$ angles $0 \leq \beta_1 \leq \ldots \leq \beta_k \leq \pi/2$ called the principal angles and their corresponding principal pairs of vectors $(u_i, w_i)\in\UU\times\WW$ for $i=1,\ldots,k$ such that $\angle(u_i,w_i) = \beta_i$ are defined by: 
	\begin{equation*}
	\begin{array}{ll}
	    u_1, w_1 \defeq \argmin\limits_{\substack{u\in\UU, w\in\WW \\ \norm{u}=\norm{w} =1}}\arccos\left(\abs{\langle u, w \rangle}\right), &
	    \beta_1 \defeq \angle(u_1, w_1)
	\end{array}    
	,\end{equation*}
	and for $i > 1$
	\begin{equation*}
	\begin{array}{ll}
	    u_i, w_i \defeq \argmin\limits_{\substack{u\perp\UU_{i-1}, w\perp\WW_{i-1} \\ \norm{u}=\norm{w} =1}}\arccos\left(\abs{\langle u, w \rangle}\right), &
	    \beta_i \defeq \angle(u_i, w_i)
	\end{array}
	,\end{equation*}
    where $$\UU_i \defeq Span\{u_j\}_{j=1}^i ~,~ \WW_i \defeq Span\{w_j\}_{j=1}^i.$$
\end{definition}
In other words, given two linear subspaces of $\RR^D$ of the same dimension $d$ we can measure the distance between them based upon the principal angles. 
In our case, we measure the distance between two subspaces by taking the maximal principal angle (maximal angle) , and denote it as
\begin{equation}\label{eq:TheAngleLinear_clean}
    \angle_{\max}(\UU,\WW) \defeq \max_{1\leq i\leq d}\beta_i
.\end{equation}

Lemmas \ref{lem:angle_space_to_vec} and \ref{lem:angle_space_perp_space} are reformulation of results proven in \cite{knyazev2007majorization} and will be also used later on.
\begin{Lemma}\label{lem:angle_space_to_vec}
Let $F$ and $G$ be two linear spaces of dimension $d$ in $\RR^D$. Assume that $\maxangle(F,G)\leq \alpha$. Then for any vector $v \in F^\perp$
\[
\min\limits_{w\in G^\perp} \angle(v,w) \leq \alpha
\]
\end{Lemma}

\begin{Lemma}\label{lem:angle_space_perp_space}
Let $F$ and $G$ be two linear spaces of dimension $d$ in $\RR^D$. Assume that $\maxangle(F.G)\leq \alpha$. Then for any vector $v \in F^\perp$ and  $w\in G$,
\[
\angle(v,w) \geq \pi/2 - \alpha
\]
\end{Lemma}

\subsection{Viewing the manifold locally as a function graph}
It is well known that, locally, a sub-manifold of $\RR^D$ can be described as a graph of a function defined from the tangent space to its orthogonal complement.
In this section, we deal with expressing a manifold as a local function graph with respect to some tilted coordinate system and bounding the size of the neighborhood for which this definition still hold.
The results reported below are general and relates closely on the concept of the Reach (see Definition \ref{def:reach}) which was introduced by Federer \cite{federer1959curvature} and further studied by Boissonat, Lieutier and Wintraecken \cite{boissonnat2017reach}. 

\begin{Lemma}[Corollary 8 from \cite{boissonnat2017reach}] \label{lem:reach_ball_no_intersect_clean}
    Let $\MM$ be a sub-manifold of $\RR^D$ with reach $\tau$ and let $p\in\MM$.
    Then, any open $D$ dimensional ball of radius $\rho \leq \tau$ that is tangent to $\MM$ at $p$ does not intersect $\MM$.
\end{Lemma}




\begin{Lemma}[Bounding Ball]\label{lem:f_bound_circle_no_func_clean}
Let $\MM$ be a $d$-dimensional sub-manifold of $\RR^D$ with reach $\tau$. For any $p \in \MM$, let $T_{p}\MM$ be the tangent of $\MM$ at $p$.
    For any $x = p + x_T$, where $x_T\in T_p\MM$, and $y\in (T_p\MM)^\perp$  such that $\|x_T\|\leq \tau$, $\|y\| \leq \tau/2$ and $(x + y)\in \MM$, we have that 
    \[
    \norm{y} \leq \tau-\sqrt{\tau^2-\norm{x_T}^2}
    \]
\end{Lemma}
The proof follows directly from Lemma \ref{lem:reach_ball_no_intersect_clean}.


\begin{Lemma}[Bounding Ball with Noise]\label{lem:dist_ri_TpM_clean}
Under the sampling assumption of \ref{sec:SamplingAssumptions}. For $r_i\in U_{ROI}$ where $U_{ROI}$ is defined in \eqref{eq:ROI_clean}. Denote $p_r = P_\MM(r)$, and $x_i = P_{T_{p_r}\MM}(r_i - p_r)$. Then,
\[
\dist(r_i - p_r,T_{p_r}\MM) \leq \tau-\sqrt{\tau^2-\|x_i\|^2}  + \sigma
.\]    
\end{Lemma}
\begin{proof}
We first note that since $r$ is at distance at most $\sigma$ from $\MM$, we have that  $U_{\textrm{ROI}} \subset B_D(p_r, \sqrt{\sigma\tau}+\sigma )$. Denote $p_i = P_\MM(r_i)$.
By Lemma \ref{lem:f_bound_circle_no_func_clean} the distance between $r_i$ and $ T_{p_r}\MM $ is bounded by 
\[
\dist(r_i-p_r,T_{p_r}\MM) \leq \dist(p_i-p_r,T_{p_r}\MM) +\sigma \leq \tau-\sqrt{\tau^2-\|x_i\|^2}  + \sigma
.\]
\end{proof}

\begin{Lemma}[Bounding Ball From a Tilted Plane] \label{lem:f_bound_circle_H_no_func_clean}

    Let $\MM$ be a $d$-dimensional sub-manifold of $\RR^D$ with reach $\tau$. For any $p \in \MM$, let $T_{p}\MM$ be the tangent of $\MM$ at $p$ and let $ H\in Gr(d, D) $ such that $ \maxangle(T_p\MM, H) = \alpha \leq \pi/4 $.
    For any $x = p + x_H$, where $x_H\in H$, $\|x_H\|\leq c_{\pi/4}\tau$ for some constant $c_{\pi/4}$ , and $y\in H^\perp$  such that $\|x-p\|\leq \tau \cos \alpha$, $\|y\| \leq \tau/2$ and $(x + y)\in \MM$, we have that 
	\[
	-\tau\cos\alpha+\sqrt{\tau^2-(\norm{x} - \tau\sin\alpha)^2} \leq \|y\| \leq \tau\cos\alpha-\sqrt{\tau^2-(\norm{x} + \tau\sin\alpha)^2}
	\]
\end{Lemma}
The proof of Lemma \ref{lem:f_bound_circle_H_no_func_clean} follows directly from applying Lemma \ref{lem:f_bound_circle_no_func_clean} and observing the illustration in Figure~\ref{fig:BoundingCircle}.

\begin{figure}
	\centering
	\includegraphics[width=0.6\linewidth]{./Figures/BoundingCircle.pdf}
	\caption{Illustration of bounding ball for $ d=1,~D=2 $. The manifold $ \MM $ is marked by the red solid line, $ T_\MM $ is marked by the blue dashed line and $ \alpha = \maxangle(H, T_p\MM) $ is the angle betweeen the $ x $-axis and $ T_p\MM $. The bounding balls defined by the reach $ \tau $ are marked in pink \emphbrackets{solid and dotted}  }
	\label{fig:BoundingCircle}
\end{figure}

\begin{Lemma}\label{lem:f_bound_circle_H_no_func_ver2_clean}
    Let $\MM$ be a $d$-dimensional sub-manifold of $\RR^D$ with reach $\tau$. For any $p \in \MM$, let $T_{p}\MM$ be the tangent of $\MM$ at $p$ and let $ H\in Gr(d, D) $ such that $ \maxangle(T_p\MM, H) = \alpha  \leq \pi/4$.
    For any $x = p + x_H$, where $x_H\in H$,$\|x_H\|\leq c_{\pi/4}\tau$ for some constant $c_{\pi/4}$, and $y\in H^\perp$  such that $\|x-p\|\leq \tau \cos \alpha$, $\|y\| \leq \tau/2$ and $(x + y)\in \MM$, we have that 
	\[
	\|y\| \leq \norm{x_H}(\tan \alpha  +\OO(\|x_H\|/\tau))
	\]
\end{Lemma}
\begin{proof}
 Recalling Lemma \ref{lem:f_bound_circle_H_no_func_clean}, we have that 
	\begin{equation}
	\|y\| \leq \tau \cos \alpha - \sqrt{\tau^2-(\norm{x_H}+\tau \sin \alpha)^2}
	.\end{equation}
	Therefore,
	\begin{align*}
	\|y\| &\leq \tau (\cos \alpha - \sqrt{1-(\norm{x_H}/\tau + \sin \alpha)^2})\\
	      &=  \tau (\cos \alpha - \sqrt{1-\norm{x_H}^2/\tau^2  + 2\norm{x_H}/\tau\sin \alpha - \sin^2\alpha})\\
	      &=  \tau (\cos \alpha - \sqrt{\cos^2\alpha-\norm{x_H}^2/\tau^2 - 2\norm{x_H}/\tau\sin \alpha })\\
	      &=  \tau \cos \alpha(1 - \sqrt{1-\norm{x_H}^2/(\tau^2\cos^2\alpha)  - 2\norm{x_H}/\tau\tan \alpha / \cos\alpha })\\
	      &\leq  \tau \cos \alpha (1 - (1-\frac{\norm{x_H}^2}{2\tau^2\cos^2\alpha}  - \frac{\norm{x_H}\tan \alpha}{\tau \cos\alpha}  - \left(\frac{\norm{x_H}^2}{\tau^2\cos^2\alpha}  + 2\frac{\norm{x_H}\tan \alpha}{\tau\cos\alpha} \right)^2))\\
	      &= \tau   \cos\alpha(\frac{\norm{x_H}^2}{2\tau^2\cos^2\alpha}  + \frac{\norm{x_H}\tan \alpha}{\tau\cos\alpha}  + \left(\frac{\norm{x_H}^2}{\tau^2\cos^2\alpha}  + 2\frac{\norm{x_H}\tan \alpha}{\tau\cos\alpha} \right)^2)\\
	      &= \norm{x_H}\tan \alpha  +\frac{\norm{x_H}^2}{2\tau\cos\alpha}  +  \left(\frac{\norm{x_H}^2}{\tau^{3/2}\cos^{3/2}\alpha}  + 2\frac{\norm{x_H}\tan \alpha}{\tau^{1/2}\cos^{1/2}\alpha} \right)^2\\
	      &\leq \norm{x_H} (\tan \alpha + c\|x_H\|/\tau)
	\end{align*}
	for some constant $c$.
\end{proof}

\begin{Lemma}\label{lem:alpha-beta_x_no_func}
	Let $\MM$ be a $d$-dimensional sub-manifold of $\RR^D$ with reach $\tau$. For any $p \in \MM$, let $T_{p}\MM$ be the tangent of $\MM$ at $p$ and let $ H\in Gr(d, D) $ such that $ \maxangle(T_p\MM, H) = \alpha\leq \pi/4 $.
	For any $x = p + x_H$, where $x_H\in H$,$\|x_H\|\leq c_{\pi/4}\tau$ for some constant $c_{\pi/4}$, and $y\in H^\perp$  such that $\|x-p\|\leq \tau \cos \alpha$, $\|y\| \leq \tau/2$ and $(x + y)\in \MM$, we denote $ \gamma = \maxangle(T_{x+y}\MM, T_{p}\MM) $.
	Then, we get that
	\[
	\sin(\maxangle (T_{x+y}\MM, T_{p}\MM)) \leq \frac{\norm{x_H}}{\tau}(1 + \tan^2 \alpha)  +\OO(\|x_H\|^2/\tau^2)
	\]
\end{Lemma}
\begin{proof}
From  Corollary 3 in \cite{boissonnat2017reach} bounds the maximal angle between the tangent spaces at two points on $\MM$ through their Euclidean distance and $\tau$.
	Namely, let $p_1, p_2 \in \MM$
	\[
	sin(\maxangle(T_{p_1}\MM, T_{p_2}\MM)/2) \leq\frac{\norm{p_1 - p_2}}{2\tau}
	.\]
	Therefore, in our case we obtain,
	\begin{equation}\label{eq:sin_alpha_1}
	\sin(\maxangle (T_{x+y}\MM, T_{p}\MM)/2) \leq \frac{\sqrt{\norm{x_H}^2 + \|y\|^2}}{2\tau}
	.\end{equation}
		
	Recalling Lemma \ref{lem:f_bound_circle_H_no_func_ver2_clean}, we have that 
	\begin{equation*}
	\|y\| \leq \norm{x_H}\tan \alpha  +\OO(\|x_H\|^2/\tau)
	\end{equation*}
	
	\begin{align*}
	\sin(\maxangle (T_{x+y}\MM, T_{p}\MM)/2) &\leq \frac{\sqrt{\norm{x_H}^2 + \|y\|^2}}{2\tau}    \\
	&= \frac{\sqrt{\norm{x_H}^2 +  \norm{x_H}^2\tan^2 \alpha  +\OO(\|x_H\|^2/\tau)}}{2\tau}\\
	&=\frac{\norm{x_H}}{2\tau}\sqrt{1 + \tan^2 \alpha  +\OO(\|x_H\|/\tau)}\\
	&\leq\frac{\norm{x_H}}{2\tau}(1 + \tan^2 \alpha)  +\OO(\|x_H\|^2/\tau^2)
	\end{align*}
	Then, as for sufficiently small $\gamma$ $sin(\gamma) < 2\sin(\gamma/2)$, for $\tau$ large enough we have
	\[
	\sin(\maxangle (T_{x+y}\MM, T_{p}\MM)) \leq \frac{\norm{x_H}}{\tau}(1 + \tan^2 \alpha)  +\OO(\|x_H\|^2/\tau^2)
	\]
\end{proof}

\begin{Lemma}\label{lem:alpha_beta_x_no_func}
	Let $\MM$ be a $d$-dimensional sub-manifold of $\RR^D$ with reach $\tau$. For any $p \in \MM$, let $T_{p}\MM$ be the tangent of $\MM$ at $p$ and let $ H\in Gr(d, D) $ such that $ \maxangle(T_p\MM, H) = \alpha\leq \pi/4 $.
	For any $x = p + x_H$, where $x_H\in H$,$\|x_H\|\leq c_{\pi/4}\tau$ for some constant $c_{\pi/4}$, and $y\in H^\perp$  such that $\|x-p\|\leq \tau \cos \alpha$, $\|y\| \leq \tau/2$ and $(x + y)\in \MM$, we denote $ \beta = \maxangle(T_{x+y}\MM, H) $.
	Then, we get that
	\[
	\alpha - 2 \sqrt{ \frac{\norm{x_H}}{\tau}\lrbrackets{2 \alpha + \frac{\norm{x_H}}{\tau}} } 
	\leq \beta \leq
	\alpha + 2 \sqrt{ \frac{\norm{x_H}}{\tau}\lrbrackets{2 \alpha + \frac{\norm{x_H}}{\tau}} }
	\]
\end{Lemma}
\begin{proof}
	Using the triangle inequality for maximal angles we have that 
	\[
	\maxangle(T_p \MM, H) \leq \maxangle(T_p \MM,T_{x+y} \MM)+\maxangle(T_{x+y} \MM,H),\]
	which can be written as 
	\[
	|\alpha - \beta| \leq \maxangle(T_{x+y} \MM, T_p \MM)
	.\]
	From Lemma~\ref{lem:alpha-beta_x_no_func} we have that 
	\[
	\sin(\maxangle (T_{x+y}\MM, T_{p}\MM)) \leq \frac{\norm{x_H}}{\tau}(1 + \tan^2 \alpha)  +\OO(\|x_H\|^2/\tau^2)
	,\]
	and thus, 
	\[
	\sin(|\alpha - \beta|) \leq \frac{\norm{x_H}}{\tau}(1 + \tan^2 \alpha)  +\OO(\|x_H\|^2/\tau^2)
	\]
	Since for $x\leq \pi/2$, we have that $x/2<\sin(x)$, we have 
	
	\[
	|\alpha - \beta| \leq 2\frac{\norm{x_H}}{\tau}(1 + \tan^2 \alpha)  +\OO(\|x_H\|^2/\tau^2)
	.\]
	Since for $\alpha \leq \pi/4$, we have that $\tan \alpha \leq 1$, and $\tan^2\alpha \leq \alpha$, we have that 
	\[
	|\alpha - \beta| \leq
	2 \sqrt{ \frac{\norm{x_H}}{\tau}\lrbrackets{2 \alpha + \frac{\norm{x_H}}{\tau}} }
	\]
	or 
	\[
	\alpha - 2 \sqrt{ \frac{\norm{x_H}}{\tau}\lrbrackets{2 \alpha + \frac{\norm{x_H}}{\tau}} } 
	\leq \beta \leq
	\alpha + 2 \sqrt{ \frac{\norm{x_H}}{\tau}\lrbrackets{2 \alpha + \frac{\norm{x_H}}{\tau}} }
	\]
	as required.
\end{proof}



\begin{Lemma} [$\MM$ is locally a function graph over a tilted plane]\label{lem:M_is_locally_a_fuinction_clean}
Let $\MM$ be a $d$-dimensional sub-manifold of $\RR^D$ with reach $\tau$. For any $p \in \MM$, let $T_{p}\MM$ be the tangent of $\MM$ at $p$. Let $H \in Gr(d,D)$, such that $\maxangle(H,T_p\MM) = \alpha \leq \pi/4$.  
Then $\MM \cap \emph{Cyl}_H(p,\rho, \tau/2)$ is locally a function over $H$, where $\emph{Cyl}_H(p,\rho, \tau/2)$ is the $D$-dimensional cylinder with the base $B_H(p, \rho)\subset H$ and height $\tau/2$ in any direction on $H^\perp$. 
Furthermore, $\rho = c_{\pi/4}\tau$ for some constant $c_{\pi/4}$. Explicitly, there exists a function 
$$
f:B_H(p, \rho) \to H^\perp
$$
such that the graph of $f$ defined as
$$
\Gamma_f = \{p + (x,f(x)) | x\in B_H(p, \rho) \} 
$$
identifies with $\MM \cap \emph{Cyl}_H(p,\rho, \tau/2)$.
\end{Lemma}
\begin{proof}
We split our arguments to two separate parts.
First, we show that for $c < 0.02$ there exists a function $f$ such that $\Gamma_f \subset M\cap \textrm{Cyl}_H(p,  c \tau, \tau/2)$.
Then, in the second part of the proof, we show that there is a constant $c_{\pi/4} < c$ such that $f$ is defined uniquely and $\Gamma_f = M\cap \textrm{Cyl}_H(p, \rho, \tau/2)$.

By definition, there is an open ball $U_T\subset T_p\MM$ of $p$ such that there is a neighborhood $W_\MM \subset \MM$ that can be pronounced as a graph of a function from $U_T\simeq \RR^d$ to $T_p\MM^\perp\simeq \RR^{D-d}$.
Accordingly, for any $H$ such that $\maxangle(H, T_p\MM) < \pi/2$ there is an open ball $B_H(p,\eps) \subset H$ such that $W_H \subset \MM$ can be pronounced as a graph of a function $f$  from $B_H(p,\eps)$ to $H^\perp$.
We wish to show that $f$ can be extended to a ball $B_H(p, 0.02 \tau)\subset H$ such that the graph of $f$ is a subset of $\MM$ (note that $f(0) = 0$).

By contradiction, let us assume that $\frak{r}$ the maximal radius of an open ball such that the $\Gamma_f \subset \MM$, is strictly smaller than $0.02 \tau$.
We claim that the graph $\Gamma_f$ is defined on the closed ball $\bar B_H(p, \frak{r})$ and is also subset of $\MM$.
This is true, from the following argument:
Take a sequence of points $\{x_n\}$ converging to $ x\in \partial \bar B_H(p, \frak{r})$, a point on the boundary of $\bar B_H(p, \rho)$, and consider $\{y_n = p + (x_n, f(x_n))\in \MM\}$. From the compactness of $\MM$ the sequence $y_n$ has a converging subsequence $y_{n_k}$ and we denote its limit as $y$.
Since $x_n \to x$, we define $f(x)\defeq\lim f(x_{n_k})$ and  $y = p + (x, f(x))$.

We now wish to show that there is $\eps>0$ such that $f$ can be extended to $B_H(p, \frak{r} + \eps)$.
Using a similar argument to the one used in the beginning of the proof, by showing that for any $x\in\partial B_H(p, \frak{r})$ the angle $\maxangle(T_x f, H) < \pi/2$, we get that there is $W_H\subset\MM$, a neighborhood of $y\in \MM$ that is the image of some function from $B_H(x,\eps_x)$ to $H^\perp$.
Therefore, $f$ can be extended into this neighborhood.
Taking $\eps$ to be the minimum over all $\eps_x$, which exists since $x$ is in $H(p, \frak{r})$, which is compact, we get that $f$ can be extended to $B_H(p, \frak{r}+ \eps)$.

The remaining piece of the existence puzzle is showing that for all $x\in\partial B_H(p, \frak{r})$ we have $\maxangle(T_x f, H) < \pi/2$.
From Lemma \ref{lem:alpha_beta_x_no_func} we have that for any $x$ such that 
\[
\frac{\pi}{2} > \frac{\pi}{4} + 2\sqrt{\frac{\norm{x}}{\tau}\left(2\frac{\pi}{4} + \frac{\norm{x}}{\tau}\right)},
\]
$\maxangle(T_xf,H) <\pi/2$ holds.
Rewriting the inequality we get
\[
\begin{array}{cc}
      \\
    \frac{\pi}{4} > & 2\sqrt{\frac{\norm{x}}{\tau}\left(\frac{\pi}{2} + \frac{\norm{x}}{\tau}\right)} \\
    \frac{\pi^2}{8^2} > & \frac{\norm{x}}{\tau}\left(\frac{\pi}{2} + \frac{\norm{x}}{\tau}\right)\\
    0 > & \frac{\norm{x}^2}{\tau^2} + \frac{\pi}{2} \frac{\norm{x}}{\tau} - \frac{\pi^2}{8^2},\\
\end{array}
\]
and thus we require
\[
\frac{\norm{x}}{\tau} < \frac{-\frac{\pi}{2} + \sqrt{\frac{\pi^2}{4} + 4 \frac{\pi^2}{8^2} }}{2} = 
\pi \left(-\frac{1}{4} + \frac{1}{2}\sqrt{\frac{1}{4} + \frac{4}{8^2} }\right)
.\]
Finally, as 
$\pi \left(-\frac{1}{4} + \frac{1}{2}\sqrt{\frac{1}{4} + \frac{4}{8^2} }\right) >0.02$ we have that for $x<0.02\tau$,
$\maxangle(T_xf,H) <\pi/2$ holds.

We now turn to show that there is a constant $c_{\pi/4}$ for which $f$ is uniquely defined in $B_H(p, c_{\pi/4}\tau)$.
From Lemma \ref{lem:f_bound_circle_H_no_func_clean} we know that for any $x \in H$ with $\|x\|\leq \tau/2$ all the $y\in H^\perp$ such that $(x,y)\in\MM$ and $\|y\| \leq \tau/2$ must satisfy:
\begin{equation}
    \|y\| \leq \tau\cos\alpha-\sqrt{\tau^2-(\norm{x} + \tau\sin\alpha)^2} = \tau\left(\cos\alpha-\sqrt{1-(\frac{\norm{x}}{\tau} + \sin\alpha)^2}\right)
\end{equation}
Let $y_1, y_2$ be such that $(x,y_1),(x,y_2)\in\MM$ where $\norm{x}=\bar x\tau$.
Then,
\begin{equation}
    \|y_j\| \leq \tau\left(\cos\alpha-\sqrt{1-(\bar x+ \sin\alpha)^2}\right),\quad (j=1,2)
.\end{equation}
In other words, $y_1$ and $y_2$ cannot be too far from one another, and note that as $\bar x \to 0$
\begin{equation}\label{eq:y2_y1_diff}
\norm{y_2 - y_1} \to 0
\end{equation}

On the other hand, taking the point $(x, y_1)\in \MM$, denoting $\beta = \maxangle (T_{(x, y_1)}\MM, H)$, and applying Lemma \ref{lem:alpha_beta_x_no_func} we have that
\[
\beta \leq
	\alpha + 2 \sqrt{ \bar x\lrbrackets{2 \alpha + \bar x} }
,\]
which tends to $\alpha$ when $\bar x\to 0$.
From Lemma \ref{lem:reach_ball_no_intersect_clean} we know that $(x,y_2)$ cannot be in any ball tangent to $\MM$ at $(x, y_1)$ of radius $\tau$.
We denote by $v$ the direction $(0, y_2 - y_1)\in H^\perp$. 
From Lemma \ref{lem:angle_space_to_vec}  we know that there is $w \in \lrbrackets{ T_{(x, y_1)} \MM} ^\perp$ such that $\angle (v, w) \leq \beta$.
Therefore, we can now limit our discussion to $L_{y_1}$ the affine space spanned by $v$ and $w$ from $(x, y_1)$, and note that it contains $(x, y_2)$ as well.
Taking the two balls $B_D((x, y_1) \pm \tau\cdot w, \tau)$ and intersecting them with $L_{y_1}$ we get two 2-dimensional disks of radius $\tau$ (see Figure \ref{fig:bouding_balls_for lem_locally_func}).
Thus, $(x, y_2)$ cannot be within either disks.
From basic trigonometry we achieve that either $y_2 = y_1$ or
\[
\norm{y_2 - y_1} \geq 2\tau\cos(\beta) \geq 2\tau\cos(\alpha + 2\sqrt{\bar x (2\alpha + \bar x)})
,\]
which tends to $2\tau \cos \alpha$ as $\bar x \to 0$.
Combining this with \eqref{eq:y2_y1_diff} we get there is $c_{\pi/4}$  such that for all $\bar x \leq c_{\pi/4}$ we have $y_1 = y_2$.
\begin{figure}
    \centering
    \includegraphics
 [width=0.5\textwidth]{Figures/BoundingCircle_Twicev2.pdf}
    \caption{Illustration of bounding balls. The reach $\tau$ of the manifold $\MM$ (marked by the red line) bounds its sectional curvatures. Namely, the manifold cannot intersect a tangent ball of radius $\tau$. In this illustration we have some coordinate system $H$ and the manifold can be described locally as a graph of some function $f:H\to H^\perp$. The coordinate system $H$ is not aligned with $T_0f$, the tangent at zero; i.e., $\maxangle(H, T_0f) = \alpha$. Then the value of $f$ at $x$ is bounded to the markered interval in $H^\perp$ above $x$. Furthermore, in order to have two different points in $\MM$ above $x$, the manifold cannot curve too fast as it cannot enter neither the dotted balls nor the solid ones.}
    \label{fig:bouding_balls_for lem_locally_func}
\end{figure}

\end{proof}


\begin{corollary} \label{cor:GraphOfFunctionTau2_clean}
Under the requirements of Lemma \ref{lem:M_is_locally_a_fuinction_clean} we get that $\MM$ is a function graph over $T_{p}\MM$ in a ${\tau}/2$ neighborhood of $p$. 
Reiterating \eqref{eq:phi_def_clean}, we have a function 
\[
\phi_p:B_{T_p\MM}(0,\tau/2) \rightarrow T_p\MM^\perp
\]
such that the graph of $\phi$ shifted to $p$ coincides with $\MM \cap \mathrm{Cyl}_{T_p\MM}(p,\tau/2,\tau/2)$.

\end{corollary}

\begin{corollary}\label{cor:GraphOfFunctionTau_ROI}
Let the requirements of Lemma \ref{lem:M_is_locally_a_fuinction_clean}, and the sampling assumptions of Section \ref{sec:SamplingAssumptions} hold. Denote the projection of $r$ onto $\MM$ by $p_r = P_\MM(r)$, let $U_{\textrm{ROI}}$ be as defined in \eqref{eq:ROI_clean}. 
Then, any $r_i \in U_{\textrm{ROI}}$ can be written as
\begin{equation}\label{eq:FunctionGraph_pi_ei_clean}
r_i = \underbrace{p_r + (x_i, \phi_{p_r}(x_i))_{T_{p_r}\MM}}_{p_i}+\eps_i
,\end{equation}
where $x_i= P_{T_{p_r}\MM}(r_i - p_r)$ and $\|\eps_i\| \leq \sigma$.
\end{corollary}
\begin{proof}
From the assumptions of Section \ref{sec:SamplingAssumptions} we know that  $\tau/\sigma > M $.
Therefore, there is $M$ such that $\sqrt{\sigma\tau}+\sigma < \tau/2$ and thus, from Corollary \ref{cor:GraphOfFunctionTau2_clean}, the intersection of $\MM$ with $\textrm{Cyl}_{T_{p_r}\MM}(p_r,\sqrt{\sigma \tau}+\sigma, \tau/2)$, a cylinder with base $B_{T_{p_r}\MM}(p_r, \sqrt{\sigma\tau}+\sigma)\subset T_{p_r}\MM$ and heights $\tau/2$ in $T_{p_r}\MM^\perp$ can be written as $\Gamma_{\phi_{p_r}, B_{T_{p_r}\MM}(p_r, \sqrt{\sigma\tau} + \sigma)}$, the graph of $\phi_{p_r}:T_{p_r}\MM\to T_{p_r}\MM^\perp$. Since $r_i$ are in a tubular neighborhood of $\MM$, the proof is concluded. 
\end{proof}


\begin{Lemma}[Function version of Lemma \ref{lem:f_bound_circle_no_func_clean}]\label{lem:f_bound_circle_clean}
    Let $\MM$ be a $d$-dimensional sub-manifold of $\RR^D$ with reach $\tau$. For any $p \in \MM$, let $T_{p}\MM$ be the tangent of $\MM$ at $p$.
    Let $\phi_p:B_{T_p\MM}(0, \tau/2)\to T_p\MM^\perp$ be defined as in Corollary \ref{cor:GraphOfFunctionTau2_clean}; that is, 
    \[\Gamma_{\phi_p,B_{T_p\MM}(0, \tau/2)}\subset\MM,\]
    where
    \[\Gamma_{\phi_p,B_{T_p\MM}(0, \tau/2)} = \{p + (x, \phi_p(x)) | x\in B_{T_p\MM}(0, \tau/2)\}.\]
    Then, for any $v \in T_p\MM^\perp$, such that $\norm{v} = 1$
    \[
    \lrangle{v, \phi_p(x)} \leq \tau-\sqrt{\tau^2-\norm{x}^2}
    \]
\end{Lemma}
\begin{proof}
    This follows immediately from Lemma \ref{lem:f_bound_circle_no_func_clean} and Lemma \ref{cor:GraphOfFunctionTau2_clean}.
\end{proof}

\begin{corollary}
It follows immediately from Lemma \ref{lem:f_bound_circle_clean} 
\begin{equation}\label{eq:eq:bound_phix_clean}
    \norm{\phi_p(x)}_{\RR^{D-d}} \leq \tau - \sqrt{\tau^2 - \norm{x}_{\RR^d}^2}
,\end{equation}
and using the triangle inequality we can say that
\begin{equation}\label{eq:bound_x_phix_clean}
    \norm{(x,\phi_p(x))}_{\RR^{D}} \leq \norm{x}_{\RR^d} + \tau - \sqrt{\tau^2 - \norm{x}_{\RR^d}^2}
.\end{equation}
\end{corollary}


\begin{Lemma}[Function version of Lemma \ref{lem:f_bound_circle_H_no_func_clean}]\label{lem:f_bound_circle_H_clean}
    Let $\MM$ be a $d$-dimensional sub-manifold of $\RR^D$ with reach $\tau$. For any $p \in \MM$, let $T_{p}\MM$ be the tangent of $\MM$ at $p$ and let $ H\in Gr(d, D) $ such that $ \maxangle(T_p\MM, H) = \alpha \leq \pi/4$.
	Let $f_p:\RR^d\to \RR^{D-d}$ be such that the neighborhood $W_p\subset\MM$ can be descried as the the graph of $f_p$
	\[\Gamma_{f_p,W_p} = \{p + (x, f_p(x)) | x\in P_{H}(W_p)\}\]
	Then, for any $v \in H^\perp$, such that $\norm{v} = 1$
	\[
	-\tau\cos\alpha+\sqrt{\tau^2-(\norm{x} - \tau\sin\alpha)^2} \leq \lrangle{v, f_p(x)} \leq \tau\cos\alpha-\sqrt{\tau^2-(\norm{x} + \tau\sin\alpha)^2}
	\]
\end{Lemma}
\begin{proof}
    This follows directly from Lemma \ref{lem:f_bound_circle_H_no_func_clean} and Lemma \ref{lem:M_is_locally_a_fuinction_clean}.
\end{proof}


\begin{Lemma}[Function Version of Lemma \ref{lem:alpha_beta_x_no_func}]\label{lem:alpha_beta_x}
    Let $\MM$ be a $d$-dimensional sub-manifold of $\RR^D$ with reach $\tau$. For any $p \in \MM$, let $T_{p}\MM$ be the tangent of $\MM$ at $p$ and let $ H\in Gr(d, D) $ such that $ \maxangle(T_p\MM, H) = \alpha \leq \pi/4$.
	Let $f_p:\RR^d\to \RR^{D-d}$ be such that the neighborhood $W_p\subset\MM$ can be descried as the the graph of $f_p$
	\[\Gamma_{f_p,W_p} = \{p + (x, f_p(x)) | x\in P_{H}(W_p)\}\]
    Let $ x_0\in H $, $ \beta(x_0) = \maxangle(T_{x_0}f_p, H) $, where $ T_\xi f $ is the tangent to the graph of $ f_p $ at $ f(\xi) $.
	Then, we get
	\[
	\alpha - 2 \sqrt{ \frac{\norm{x_0}}{\tau}\lrbrackets{2 \alpha + \frac{\norm{x_0}}{\tau}} } 
	\leq \beta(x_0) \leq
	\alpha + 2 \sqrt{ \frac{\norm{x_0}}{\tau}\lrbrackets{2 \alpha + \frac{\norm{x_0}}{\tau}} }
	\]
\end{Lemma}

\begin{proof}
 This follows immediately from Lemma \ref{lem:alpha_beta_x_no_func} and Lemma \ref{lem:M_is_locally_a_fuinction_clean}.
\end{proof}
\section{Supporting lemmas for Step 1}
\subsection{Proof of Lemma \ref{lem:J1pTp_clean}} \label{subsec:proof_lem_J1pTp_clean}
\begin{proof}
We recall that $ U_\textrm{ROI} $ as defined in \eqref{eq:ROI_clean} is contained in $ B_D({p_r}, \sqrt{\sigma\tau} + \sigma) $.
Using Lemma \ref{lem:dist_ri_TpM_clean}, for any $r_i\in U_\textrm{ROI}$ we know that the distance between $r_i - p_r$ and $ T_{{p_r}}\MM $ is bounded by 
\begin{equation}\label{eq:d_ri_Tp_clean}
\dist(r_i - p_r,T_{{p_r}}\MM) \leq \tau-\sqrt{\tau^2-\|x_i\|^2}  + \sigma
,\end{equation}
where $ x_i = P_{T_{p_r}\MM}(r_i - p_r) $.
Since $ x_i \in B_D(0, \sqrt{\sigma\tau}+\sigma )$ we have that $\norm{x_i} \leq \sqrt{\sigma\tau}+\sigma $, and so
\[
J_1(r; {p_r}, T_{{p_r}}\MM) \leq \left(\tau-\sqrt{\tau^2-(\sqrt{\sigma\tau}+\sigma)^2} + \sigma\right)^2 
.\]
By simplifying and bounding this expression using Remark \ref{rem:taylor2_sqrt1-x2_clean} we get
	\begin{align*}
		\left((\sigma+\tau)-\sqrt{\tau^2-(\sqrt{\sigma\tau}+\sigma)^2}\right)^2 = & (\sigma+\tau)^2 + \tau^2-(\sqrt{\sigma\tau}+\sigma)^2 - 2(\sigma+\tau) \sqrt{\tau^2-(\sqrt{\sigma\tau}+\sigma)^2}\\
		= &
		\sigma^2 + 2\sigma\tau + 2 \tau^2 -  (\sqrt{\sigma\tau}+\sigma)^2 \\
		&- 2\tau(\sigma+\tau) \sqrt{1-(\sqrt{\sigma/\tau}+\sigma/\tau)^2}\\
		\leq &
		\sigma^2 + 2\sigma\tau + 2 \tau^2 - (\sqrt{\sigma\tau}+\sigma)^2\\
		& - 2\tau^2 (1-1/2 (\sqrt{\sigma/\tau}+\sigma/\tau)^2- (\sqrt{\sigma/\tau}+\sigma/\tau)^4 )\\& - 2\tau\sigma (1-1/2 (\sqrt{\sigma/\tau}+\sigma/\tau)^2- (\sqrt{\sigma/\tau}+\sigma/\tau)^4)\\
		= &
		\sigma^2 -  \tau^2(\sqrt{\sigma/\tau}+\sigma/\tau)^2 \\
		& + \tau^2  ((\sqrt{\sigma/\tau}+\sigma/\tau)^2  +2(\sqrt{\sigma/\tau}+\sigma/\tau)^4) \\& + \tau\sigma  ((\sqrt{\sigma/\tau}+\sigma/\tau)^2+2(\sqrt{\sigma/\tau}+\sigma/\tau)^4) \\
		= &
		\sigma^2 + 2\tau^2 (\sqrt{\sigma/\tau}+\sigma/\tau)^4\\& + \tau\sigma ( (\sqrt{\sigma/\tau}+\sigma/\tau)^2 +2(\sqrt{\sigma/\tau}+\sigma/\tau)^4)\\
		\leq &
		\sigma^2 + 2 \left(\sqrt{\sigma} + \frac{\sigma}{\sqrt{\tau}}\right)^4
		+ 3\sigma
		\left(\sqrt{\sigma} + \frac{\sigma}{\sqrt{\tau}}\right)^2 %+		1/4 \left(\frac{\tau^{0.25}\sigma^{0.25}}{k} + \frac{\sigma^{1.25}}{\tau^{0.75}}\right)^4
	.\end{align*}
Using the fact that $ \sigma \leq \tau $ we get
\begin{align*}
J_1(r; {p_r}, T_{{p_r}}\MM) &\leq
\sigma^2 + 2(2\sqrt{\sigma})^4 + 3\sigma (2\sqrt{\sigma})^2
\end{align*}
Thus, we obtain
$$J_1(r; {p_r}, T_{{p_r}}\MM) \leq 49 \sigma^2.$$
\end{proof}
\subsection{Technical parts of Lemma \ref{thm:J1pH_clean}}
\subsubsection{Proof of Claim \ref{thm:claim_R1pp_bound_clean} of  Lemma \ref{thm:J1pH_clean}}\label{sec:proof_claim_R1pp_bound_clean}
\begin{proof} 
It is clear from \eqref{eq:R1pp_def_clean} that
\[
R_1''(r; p, H) \geq
- \sum\limits_{j\in\KK''} \frac{1}{\#|U_\textrm{ROI}|}\sum\limits_{r_i\in U_\textrm{ROI}}   \langle r_i - p, {y}_j \rangle^2 =  - \frac{1}{\#|U_\textrm{ROI}|}\sum\limits_{r_i\in U_\textrm{ROI}}\sum\limits_{j\in\KK''}\langle r_i - p, {y}_j \rangle^2
.\]
Furthermore, since $y_j\in T_p\MM^\perp$ and by Lemma \ref{lem:dist_ri_TpM_clean} we get 
\[
\sum\limits_{j\in\KK''}\abs{\lrangle{r_i - p, y_j}}^2 \leq dist^2(r_i - p, T_p\MM) \leq \left(\tau - \sqrt{\tau^2 - \norm{P_{T_p\MM}(r_i- p) }^2} + \sigma\right)^2 
,\]
and since $r_i\in U_\textrm{ROI} \subset B_D(p, \sqrt{\sigma\tau} + \sigma)$ we get
\begin{align*}
     \sum\limits_{j\in\KK''}\langle r_i - p,y_j \rangle^2 
     & \leq 
	\left(\tau+\sigma-\sqrt{\tau^2-(\sqrt{\sigma\tau}+\sigma)^2}\right)^2 \\
	& = \lrbrackets{\tau + \sigma - \tau\sqrt{1 - \lrbrackets{\sqrt{\frac{\sigma}{\tau}} + \frac{\sigma}{\tau} }^2 }}^2
\end{align*}
by Taylor expansion (Remark \ref{rem:taylor_sqrt1-x2_clean})
\begin{align*}
     \sum\limits_{j\in\KK''}\langle r_i - p,y_j \rangle^2 
	& \leq 
	\lrbrackets{\tau + \sigma - \tau\lrbrackets{1 - \frac{1}{2}\lrbrackets{\sqrt{\frac{\sigma}{\tau}} + \frac{\sigma}{\tau} }^2}}^2 \\
	&=
	\lrbrackets{\sigma + \tau \frac{1}{2}\lrbrackets{\sqrt{\frac{\sigma}{\tau}} + \frac{\sigma}{\tau} }^2}^2 \\
\end{align*} 
since $ \sigma < \tau $ we get $\sqrt{\frac{\sigma}{\tau}} > \frac{\sigma}{\tau} $ and
\begin{align}\label{eq:ri_yk_bound_s_square_clean}
     \sum\limits_{j\in\KK''}\langle r_i - p,y_j \rangle^2 
	& \leq 
	\lrbrackets{\sigma + 2\tau \frac{\sigma}{\tau}}^2 = 9 \sigma^2
\end{align}
and
\begin{equation*}
R_1''(r; p, H) \geq
-  \frac{1}{\#|U_\textrm{ROI}|}\sum\limits_{r_i\in U_\textrm{ROI}}   9 \sigma^2 =  -9 \sigma^2
\end{equation*}

\end{proof}


\subsubsection{Simplification of \eqref{eq:a1a2_restriction_clean} to achieve \eqref{eq:a1a2_simplification_clean}}\label{subsec:simplification_app}
We wish to rewrite the requirement of \eqref{eq:a1a2_restriction_clean} in terms of $a_1$ and $a_2$.
By Lemma \ref{lem:f_bound_circle_clean}, for all $x\in B_{T_p\MM}(x_0, a_2 \sqrt{\sigma\tau})$:
\begin{equation*}
    \|\phi_p(x)\| \leq \tau - \sqrt{\tau^2 - (a_1+a_2)^2 {\sigma\tau}} 
;\end{equation*}
see Figure \ref{fig:ManifoldGoodDisc_clean}.
Thus,
\begin{align*}
\norm{(x,\phi_p(x))_{T_p\MM}}
& \leq \sqrt{(a_1+a_2)^2 {\sigma\tau} + \left(\tau - \sqrt{\tau^2 - (a_1+a_2)^2 {\sigma\tau}}\right)^2} \\
&= \sqrt{(a_1+a_2)^2{\sigma\tau} + \tau^2 - 2\tau\sqrt{\tau^2 - (a_1+a_2)^2 {\sigma\tau}} + \tau^2 - (a_1+a_2)^2 {\sigma\tau}} \\
&= \sqrt{2\tau^2 - 2\tau^2\sqrt{1 - (a_1+a_2)^2 {\sigma/\tau}} } \\
&= \tau\sqrt{2} \sqrt{1 - \sqrt{1 - (a_1+a_2)^2 {\sigma/\tau}} } \\
&\leq \tau\sqrt{2} \sqrt{1 - (1 - (a_1+a_2)^2 {\sigma/\tau}) } \\
&= \sqrt{2}(a_1+a_2) \tau\sqrt{ {\sigma/\tau} } \\
&= \sqrt{2}(a_1+a_2) \sqrt{ {\sigma\tau} },\end{align*}
Where the second inequality results from applying \eqref{eq:Taylor2ndOrder}.
Therefore, the requirement of \eqref{eq:a1a2_restriction_clean} translates to
$$
\norm{(x,\phi_p(x))_{T_p\MM}}_{T_p\MM,y_{\tilde j}}\leq \sqrt{2}(a_1+a_2)\sqrt{\sigma\tau}<\sqrt{\sigma\tau}-2\sigma
,$$
which can be simplified to 
$$
(a_1+a_2)<\frac{1}{\sqrt{2}} - \frac{\sqrt{2\sigma}}{\sqrt{\tau}} 
.$$


\subsection{Supporting Lemmas on Sample size in a given volume}
% --- begin Section_Lemmas_Step1.tex ---
In this section we concentrated all assisting lemmas that are used in the proofs of Step 1.

\begin{Lemma}[Number of samples in a ball]\label{lem:sampling_density_almost_uniform_clean}
Suppose $\nu$ is a distribution on $\Omega\subset B_d(0,R) \subset \RR^d$ which is close to the uniform distribution $\mu$. That is, there exists $\mu_{max}$, $\mu_{min}$  such that for any $x\in \Omega$ we have  $\mu_{\min}\mu(x) \leq \nu(x) \leq \mu_{\max}\mu(x)$. Suppose $X=\{x_j\}_{j=1}^n$ is a set of $n$ i.i.d. sample from $\nu$, and denote the volume of a $ d $-dimensional unit ball by  $ \mathrm{V}_d = \frac{\pi^{d/2}}{\Gamma(d/2+1)} $. For any $\eps$, $\delta$ and radius $\rho$, there is $N$, such that if $n>N$ the following holds: For any $x_0\in \Omega$ such that $B_d(x_0,\rho)\subset \Omega$, we have 
\[
n(\frac{\mu_{max}}{2}\cdot \mathrm{V}_d \cdot\rho^d-\eps) < \#|X \cap B_d(x_0, \rho)| < n(2\cdot \mu_{min}\cdot \mathrm{V}_d \cdot\rho^d+\eps)
\]
with probability of at least $1-\delta$, where $\#|A|$ denotes the number of elements in the set $A$.
\end{Lemma}
\begin{proof}
Since $\Omega \subset B_d(0,R)\subset\RR^d$, there exists an $\tilde\eps$-net (denoted by $\Xi$) such that 
\[
\#\abs{\Xi} = \ceil{ \frac{3R}{\tilde\eps}}^d
,\]
where $ \ceil{x} $ is the ceiling value of $ x $ \cite{vershynin2018high}.
Around each point $p$ in $\Xi$, we consider a ball $B_d(p, (1 - \tilde\eps)\rho)$. 
Note, that this $\tilde\eps$-net along with these balls are independent of the choice of a specific ball $B_d(x_0,\rho)$.

For each of the $B_d(p, (1-\tilde\eps)\rho)$ , we  consider our sample set as $n$ i.i.d random variables   $Z^p_j$ which return the value $1$ if the sample lies within $B_d(p, (1-\tilde\eps)\rho)$ and $0$ if not.
Naturally, we get for all $j$ that
\[
\Pr[z_j^p = 1] = \int\limits_{B_d(p, (1-\tilde\eps)\rho)} d\nu 
\] 
Applying Hoeffding's inequality for each of the $B_d(p, (1 - \tilde\eps)\rho)$ we arrive at
\begin{equation}\label{eq:HoeffingtonInequality1_clean}
\Pr[\overline{Z}^p - \EE[\overline{Z}^p] \leq -\eps] \leq e^{-2 n \eps^2}
,    
\end{equation}
and
\begin{equation}\label{eq:HoeffingtonInequality2_clean}
\Pr[\overline{Z}^p - \EE[\overline{Z}^p] \geq \eps] \leq e^{-2 n \eps^2}
,    
\end{equation}
where
\[
\overline{Z}^p = \frac{1}{n}\sum_{j=1}^n Z_j^p.
\]
As a result
\[
\EE[\overline{Z}^p] = \frac{1}{n}\sum_{j=1}^n\Pr[Z_j^p = 1] =
\frac{1}{n}\sum_{j=1}^n\int\limits_{B_d(p, (1 - \tilde\eps)\rho)}d\nu
,\]
and
\begin{equation}\label{eq:EzBounds_clean}
\mu_{min} \cdot \vol(B_d(p, (1 - \tilde\eps)\rho)) \leq \EE[\overline{Z}^p] \leq \mu_{max} \cdot \vol(B_d(p, (1 - \tilde\eps)\rho))
.\end{equation}
Plugging this into \eqref{eq:HoeffingtonInequality1_clean} we get
\begin{equation}\label{eq:PrBound1_clean}
\Pr\left[\overline{Z}^p - \mu_{max}\vol\left(B_d(p, (1- \tilde\eps)\rho)\right)\leq -\eps\right] \leq e^{-2n \eps^2}
,\end{equation}
or, alternatively, since $\#|X \cap B_d(p,(1-\tilde\eps)\rho)| = n \cdot \overline{Z}^p$ we get
\[
\Pr\left[\#|X \cap B_d(p,(1-\tilde\eps)\rho)| \leq n(\mu_{max}\vol(B_d(p,(1-\tilde\eps)\rho))-\eps)\right]\leq e^{-2n \eps^2}
.\]

Denoting by $A_p$ the event $\#|X \cap B_d(p,(1-\tilde\eps)\rho)| \leq n(\mu_{max}\vol(B_d(p,(1-\tilde\eps)\rho))-\eps)$ we use the union bound to achieve 
\begin{equation}\label{eq:UnionBoundAp_clean}
\Pr\left(\bigcup\limits_{p\in\Xi} A_p\right) \leq \sum\limits_{p\in\Xi} \Pr\left(A_p \right) \leq \ceil{ 3R/\tilde\eps}^d \cdot  e^{-2n \eps^2}
.\end{equation}
Explicitly, the chances that there exists $B_d(p,(1-\tilde\eps)\rho)$ containing \textbf{less} sampled points than $n(\mu_{max}\vol(B_d(p,(1-\tilde\eps)\rho))-\eps)$ are less than $c \cdot e^{-2n \eps^2}$, where $c = \ceil{ 3R/\tilde\eps}^d$.
Going back to $B_d(x_0, \rho)$, we know that there exists a point $\tilde p\in \Xi$ such that 
\[
B_d(\tilde p,(1-\tilde\eps)\rho)\subset B_d(x_0, \rho)
.\]
As a result, for any $\delta, \eps, \rho$ there exists $N$ such that for all $n>N$
\[
\#|X \cap B_d(x_0, \rho)| > n(\mu_{max}\vol(B_d((1-\tilde\eps)\rho))-\eps)
= n(\mu_{max}\cdot \mathrm{V}_d \cdot(1-\tilde\eps)^d\rho^d-\eps)
,\]
with probability larger than $1-\delta$, where
\[
\mathrm{V}_d = \frac{\pi^{d/2}}{\Gamma(d/2+1)}
.\]

Similarly, instead of considering $B_d(p, (1-\tilde\eps)\rho)$ we look at $B_d(p, (1+\tilde\eps)\rho)$ for $p\in\Xi$ and alter the definitions of $Z^p$ accordingly.
Then, by plugging the left inequality of \eqref{eq:EzBounds_clean} into \eqref{eq:HoeffingtonInequality2_clean} we get
\[
\Pr[\overline{Z}^p - \mu_{min}\vol(B_d((1+\tilde\eps)\rho))\geq \eps] \leq e^{-2n\eps^2} 
\]
Utilizing the union bound once more we get the same bound as in \eqref{eq:UnionBoundAp_clean}. 
In other words, the chances that there exists a $B_d(p, (1+\tilde\eps)\rho)$ containing \textbf{more} sampled points than $n(\mu_{min} \cdot \vol(B_d((1+\tilde\eps)\rho)) + \eps )$ are less than $\ceil{3R/\tilde\eps}^d\cdot e^{-2n\eps^2}$.

Going back to $B_d(x_0, \rho)$, we know that there exists a point $\tilde p\in \Xi$ such that 
\[
B_d(x_0, \rho) \subset B_d(\tilde p,(1+\tilde\eps)\rho)
,\]
and for $\tilde\eps$ small enough
\[
B_d(\tilde p,(1+\tilde\eps)\rho) \subset \Omega
.\]
As a result, for any $\delta, \eps, \rho$ there exists $N$ such that for all $n>N$
\[
\#|X \cap B_d(x_0, \rho)| < n(\mu_{min}\vol(B_d((1+\tilde\eps)\rho))+\eps)
= n(\mu_{min}\cdot \mathrm{V}_d \cdot(1+\tilde\eps)^d\rho^d+\eps)
,\]
with probability larger than $1-\delta$.

Finally, we get that for any $\delta, \eps, \rho$ there exists $N$ large enough such that if $n>N$ we get
\[
n(\mu_{max}\cdot \mathrm{V}_d \cdot(1-\tilde\eps)^d\rho^d-\eps) < \#|X \cap B_d(x_0, \rho)| < n(\mu_{min}\cdot \mathrm{V}_d \cdot(1+\tilde\eps)^d\rho^d+\eps)
\]
Since this is true for any $\tilde\eps$, $(1+1/a)^a \leq 3$ and $(1-1/a)^a \geq 0.25$ we can choose $\tilde\eps = \frac{1}{c_1 d}$ such that
\[(1 - \tilde\eps)^d \geq 0.25^{1/c_1} \geq 0.5,\]
and
\[(1 + \tilde\eps)^d \leq 3^{1/c_1} \leq 2,\]
and achieve
\[
n(\frac{\mu_{max}}{2}\cdot \mathrm{V}_d \cdot\rho^d-\eps) < \#|X \cap B_d(x_0, \rho)| < n(2\cdot \mu_{min}\cdot \mathrm{V}_d \cdot\rho^d+\eps)
\]
\end{proof}

\begin{Lemma}[The projection of the Lebesgue measure onto $T_p\MM$ is almost uniform] \label{lem:measure_projecting_to_tangent_clean}
    Let $\MM$ be a $d$-dimensional sub-manifold of $\RR^D$ with bounded reach $\tau$ and a Riemannian metric $G$ pronounced through the chart $\varphi_p$ around a point $ p\in \MM $ \eqref{eq:LocalChart_clean}. 
    Let $r\in\MM_\sigma$, and let $\mu, \mu_\MM, \mu_{T_p\MM}$ denote the uniform distribution on $\MM_\sigma\subset\RR^D$, $\MM$, $T_p\MM$ correspondingly.
    Denote $P_\MM, P_{T_p\MM}$ the projection operators onto $\MM$ and $T_p\MM$.
    Then $\left(P_{T_p\MM}P_{\MM}\right)_*\mu$ is a measure on $T_p\MM$, and upon restricting this measure to $B_d(0,\rho)$ for some $\rho < \tau/2$ we get
    \[
    (P_{T_p\MM}P_\MM)_*\mu = \mathrm{V}_{D-d}\sigma^{D-d} \sqrt{det(G)}\mu_{T_p\MM}
    ,\]
    where $\mathrm{V}_d$ denotes the volume of a $d$-dimensional ball.
\end{Lemma}

\begin{proof}
We first note that since $\mu$ is the Lebesgue measure on $\MM_\sigma$ we have
\[
\int_{\MM_\sigma} d\mu = \int_{\MM}\mathrm{V}_{D-d}\sigma^{D-d}\mu_\MM
,\]
where $\mu_\MM$ is the uniform distribution on $\MM$. Thus,
\[
P_\MM \mu = {\mathrm{V}_{D-d}\sigma^{D-d}}\mu_\MM
.\]
Now, $P_\MM\mu$ is a measure defined on $\MM$, which can be pulled back to the tangent domain $T_p\MM\simeq\RR^d$ in the neighborhood $B_d(p, \tau/2)$ according to Corollary \ref{cor:GraphOfFunctionTau2_clean}.
If we denote the chart from $T_p\MM\simeq\RR^d$ to $\MM$ by $\varphi_p$ we get
\[
(P_{T_p\MM}P_\MM)_*\mu = \mathrm{V}_{D-d}\sigma^{D-d} \sqrt{det(G)}\mu_{T_p\MM}
,\]
where $G$ is the Riemannian metric expressed in this chart, and $\mu_{T_p\MM}$ is the Lebesgue measure on $T_p\MM$.

\end{proof}

\begin{corollary}\label{cor:measure_projecting_bounds_clean}
	From the fact that $\MM\in\C^k$ is compact and the restriction to a ball of radius $\tau/2$ we get that $\sqrt{det(G)}$ is bounded and
	\[
	\mu_{min} \mu_{T_p\MM} \leq (P_{T_p\MM}P_\MM)_*\mu \leq \mu_{max} \mu_{T_p\MM}
	\]
	where $\mu_{min}, \mu_{max}$ are constans that depend on $\tau$, and $\mu_{T_p\MM}$ is the Lebesgue measure on $T_p\MM$.
	The constants $ \mu_{min}, \mu_{max} $ can be described explicitly to show their exact relationship to $ \tau $.
\end{corollary}

Combining Lemma \ref{lem:sampling_density_almost_uniform_clean} and Corollary \ref{cor:measure_projecting_bounds_clean}, we have the following result

\begin{Lemma}\label{lem:num_of_samples_in_a_ball_clean}
	Let $\MM$ be a compact $d$-dimensional sub-manifold of $\RR^D$ with reach $\tau$  bounded away from zero, and a Riemannian metric $G_p$ pronounced through the chart $\varphi_p$ around a point $ p\in \MM $ \emphbrackets{see \eqref{eq:LocalChart_clean}}. 
    Let $\MM_\sigma$ be a tubular neighborhood around $ \MM $ of radius $ \sigma $ \emphbrackets{see \eqref{eq:Msigma}}, and assume $\sqrt{\frac{\sigma}{\tau}} < \frac{1}{2}$. 
    Suppose that $\mu$ is the uniform distribution on $ \MM_\sigma $. 
    Let $X = \{r_1, \ldots,r_n\}$ be $n$ points sampled i.i.d from $\mu$, and denote the volume of a $d$-dimensional unit ball by  $ \mathrm{V}_d = \frac{\pi^{d/2}}{\Gamma(d/2+1)} $.
    Denote,
    \begin{align}
    \label{eq:mu_min_max_def}
    \begin{split}
    \mu_{\min} &= \mathrm{V}_{D-d}\sigma^{D-d} \min_{\substack{p\in \MM\\ x\in B_{T_p\MM}(0,\sqrt{\sigma\tau}- \sigma)}}  \sqrt{det(G_p(x))}\\
    \mu_{\max} &= \mathrm{V}_{D-d}\sigma^{D-d}\max_{\substack{p\in \MM \\ x \in B_{T_p\MM}(0,\sqrt{\sigma\tau}- \sigma)}}  \sqrt{det(G_p(x))}
    .\end{split}
    \end{align}
    
    Then for any $\eps$ and $\delta$, there is $N$, such that for all $n>N$ the following holds: For any $x_0\in T_p\MM$ and $\rho \in \RR^+$ such that $B_{T_p\MM}(x_0,\rho + \sigma)\subset  B_{T_p\MM}(0,\sqrt{\sigma\tau}- \sigma)$ \emphbrackets{see the red, green, and blue discs in Figure \ref{fig:ManifoldGoodDisc_clean}}, we have 
\begin{align*}
\begin{split}
&\#\{r_i|P_{T_p\MM}(r_i) \in B_{T_p\MM}(x_0,\rho) \}
\leq
n(2\cdot \mu_{min}\cdot \mathrm{V}_d \cdot\rho^d+\eps) \\
&\#\{r_i|P_{T_p\MM}(r_i) \in B_{T_p\MM}(x_0,\rho+\sigma) \}
\geq
n \left(\frac{\mu_{max}}{2}\cdot V \cdot \rho^d - \eps\right)
\end{split}
\end{align*}
with probability of at least $1-\delta$.
\end{Lemma}


\begin{proof}
    Note that both the minimum and maximum in \eqref{eq:mu_min_max_def} exist and finite since $ \MM $ is compact and the determinant is continuous.
	We mention that since $ \sqrt{\frac{\sigma}{\tau}} < \frac{1}{2} $ we get that $ \sqrt{\sigma \tau} < \frac{\tau}{2} $ and the conditions of Corollary \ref{cor:measure_projecting_bounds_clean} are met.
	Note that 
	\[
	\#\{r_i|P_{T_p\MM}(r_i) \in B_d(x_0,\rho) \} \leq \#\{r_i|P_{T_p\MM}\circ P_\MM(r_i) \in B_d(x_0,\rho) \}
	,\]
	and from Lemma \ref{lem:sampling_density_almost_uniform_clean} combined with Corollary \ref{cor:measure_projecting_bounds_clean} we get
	\[
	\#\{r_i|P_{T_p\MM}\circ P_\MM(r_i) \in B_d(x_0,\rho) \} \leq n(2\cdot \mu_{min}\cdot \mathrm{V}_d \cdot\rho^d+\eps)
	.\]
	Thus, the first inequality is achieved.
	
	On the other hand,
	\[
	\#\{r_i|P_{T_p\MM}(r_i) \in B_d(x_0,\rho+\sigma) \}
	\geq \#\{r_i|P_{T_p\MM}\circ P_\MM(r_i) \in B_d(x_0,\rho)\}
	.\]
	Using again Lemma \ref{lem:sampling_density_almost_uniform_clean} combined with Corollary \ref{cor:measure_projecting_bounds_clean} we get
	\[
	\#\{r_i|P_{T_p\MM}\circ P_\MM(r_i) \in B_d(x_0,\rho)\} \geq n \left(\frac{\mu_{max}}{2}\cdot \mathrm{V}_{d} \cdot \rho^d - \eps\right)
	,\]
	and the second inequality holds.
\end{proof}
% --- end Section_Lemmas_Step1.tex ---


\section{Supporting lemmas for Step 2}


\begin{Lemma}\label{lem:Tf_Tftilde_Init}
	Let the sampling assumptions of Section \ref{sec:SamplingAssumptions} hold.
	Let $(q_{-1}, H_0)\in \MM_\sigma\times Gr(d,D)$ be the initialization of Algorithm \ref{alg:step2_clean}. Assume $\maxangle (H_0, T_p\MM) \leq \alpha$, for $p = P_\MM(q_{-1})$.
	Let $ f_{-1/2}:\RR^d\to\RR^{D-d} $ be defined in \eqref{eq:fl12_def}; i.e., $f_{-1/2}$ is a function whose graph $ \Gamma_{f_{-1/2}, q_{-1}, H_0} $ coincides with $ \MM $ in the sense of Lemma \ref{lem:M_is_locally_a_fuinction_clean}; explicitly 
	\[
	\Gamma_{f_{-1/2}, q_{-1}, H_0} = \{q_{-1} + (x, f_{-1/2}(x))_{H_0} ~|~ x\in B_{H_0}(q_{-1}, c_{\pi/4}\tau)\}
	,\]
	where $c_{\pi/4}$ is a constant and $\Gamma_{f_{-1/2}, q_{-1}, H_0} \subset \MM$.
	Then, for any $\alpha \leq \alpha_c$ where $\alpha_c$ is a constant depending only on $c_{\pi/4}$ of Lemma \ref{lem:M_is_locally_a_fuinction_clean}, there is $M$ such that for $\tau/\sigma >M$ (see assumption \ref{assume:noise} in Section \ref{sec:SamplingAssumptions}) we have 
	\[
	\maxangle(H_0, T_0 f_{-1/2}) 
	\leq  \frac{3\alpha}{2}
	.\]
\end{Lemma}
\begin{proof}
    Let $\phi_p: T_p\MM \simeq \RR^d \rightarrow \RR^{D-d}$ be defined in \eqref{eq:phi_def_clean} and \eqref{eq:FunctionGraph_init}. Let $g_0:(q_{-1},T_p\MM) \simeq \RR^d \to \RR^{D-d}$ defined as $g_0(x) = \phi_p(x) + (p - q_{-1})$ (note that $ (p-q_{-1}) \in T_p \MM^\perp$ and thus this can be understood with some abuse of notation). From Corollary \ref{cor:GraphOfFunctionTau2_clean}, $\Gamma_{g_0, q_{-1}, T_p\MM}$ coincides with $\MM \cap \mathrm{Cyl}(p,\tau/2, \tau/2)$.
    We note that $\maxangle (T_0g_0, T_p\MM)=0$,  $\maxangle(H_0, T_p\MM)\leq \alpha$ where $\alpha \leq \alpha_c$, $\|p-q_{-1}\| \leq \sigma$ and from assumption \ref{assume:noise} in Section \ref{sec:SamplingAssumptions}, we have that $\sigma \leq \frac{3\tau}{4 \cdot 16}$.  
    Then, we can apply Lemma \ref{lem:shift_ang_general}, with the above defined $g_0$ and $g_1 = f_{-1/2}$, $G_0 = T_p\MM$, and $G_1 = H_0$ and get
    \[
    \maxangle (H_0, T_0f_{-1/2}) \leq \maxangle(H_0, T_0g_0)  + \frac{ 8\|g_0(0)\|}{\tau}
    .\]
Since $T_0g_0 = T_p\MM$ and $\frac{ 8\|g_0(0)\|}{\tau} \leq \frac{8\sigma}{\tau}$, we have 
\[
    \maxangle (H_0, T_0f_{-1/2}) \leq \alpha + \frac{8\sigma}{\tau}.
\]
Finally, from assumption \ref{assume:noise} in Section \ref{sec:SamplingAssumptions}, we can require that $8\frac{\sigma}{\tau}\leq \frac{\alpha}{2}$, and thus,
\[
    \maxangle (H_0, T_0f_{-1/2}) \leq \frac{3\alpha}{2}.
\]


\end{proof}
\begin{Lemma}\label{lem:comute_kappa}
For any $\delta$ and for any $n,\alpha_1,r_1$, Let $C_{0}$ be the constant from Theorem 3.2 of \cite{aizenbud2021VectorEstimation}. We have that 
    \begin{equation}\label{eq:kappa_res_lem}
        \kappa = r_1\log_2(n) + \bar{C}_{\alpha_1,d} - \log\left(\ln\left(\frac{2r_1\log_2(n) + 2\bar{C}_{ \alpha_1,d}}{\delta}\right)\right)
    ,\end{equation}
    where
  \[
    \bar{C}_{ \alpha_1,d} = 1 + \log_2\left(\frac{ \alpha_1}{12\sqrt{d}}\right) -\log_2(C_0 )
    ,
    \]
    satisfies 
    \begin{equation}\label{eq:n_bound1_kappa_lem}
        2^{\kappa-1}C_0 \ln(1/\delta_1) \leq \frac{ \alpha_1}{12\sqrt{d}}n^{r_1},
    \end{equation}
    for $\delta_1 = \frac{\delta}{2\kappa}$, and $C_0$ from Theorem 3.2 of \cite{aizenbud2021VectorEstimation}.
    Furthermore, for $\kappa$ as in \eqref{eq:kappa_res_lem} we have 
    \[
     \alpha_1 2^{-\kappa} \leq C_{d} \ln\left(\frac{1}{\delta}\right) n^{-r_1}\left(\ln\left(\ln(n) \right)\right)^{2r_1}
    \]
\end{Lemma}

\begin{proof}
We find $\kappa$ that will satisfy \eqref{eq:n_bound1_kappa_lem}. Recalling that $\delta_1 = \frac{\delta}{2\kappa}$, we have that
	\[
    C_0 \left(\ln \frac{1}{\delta_1}\right) = C_0 \left(\ln \frac{2\kappa}{\delta}\right).
	\]
Rewriting \eqref{eq:n_bound1_kappa_lem}, we need $\kappa$ to satisfy
    \begin{equation}
        2^{\kappa-1}C_0\left(\ln \frac{2\kappa}{\delta}\right) \leq \frac{ \alpha_1}{12\sqrt{d}}n^{r_1},
    \end{equation}
    or, taking $\log_2$ of both sides, we have
    \begin{equation}
        \kappa - 1 +\log_2(C_0 ) + \log_2\left(\ln \frac{2\kappa}{\delta}\right) \leq \log_2\left(\frac{ \alpha_1}{12\sqrt{d}}\right) + r_1\log_2(n),
    \end{equation}
    or, 
    \begin{equation}\label{eq:kappa_bound1}
        \kappa  + \log_2\left(\ln \frac{2\kappa}{\delta}\right) \leq 1 + \log_2\left(\frac{ \alpha_1}{12\sqrt{d}}\right) -\log_2(C_0 ) +  r_1\log_2(n),
    \end{equation}
    To simplify the expression we denote the RHS of \eqref{eq:kappa_bound1} by $x$. Then we are looking for $\kappa$ such that
    \begin{equation}\label{eq:kappa_bound_simplified_x}
    \kappa+\log_2\left(\ln\left(\frac{2\kappa}{\delta}\right)\right) < x     
    \end{equation}
    We note that 
    \[
    \kappa = x-\log_2\left(\ln\left(\frac{2x}{\delta}\right)\right)
    \]
    satisfies \eqref{eq:kappa_bound_simplified_x} since
    \[
    x-\log_2\left(\ln\left(\frac{2x}{\delta}\right)\right)   +   \log_2\left(\ln\left( \frac{2x-2\log_2\left(\ln\left(\frac{x}{\delta}\right)\right)}{\delta} \right) \right)<x
    \]
    Thus, the following $\kappa$ satisfies Eq. \eqref{eq:kappa_bound1}
    \begin{equation}\label{eq:kappa_res}
        \kappa = r_1\log_2(n) + \bar{C}_{ \alpha_1,d} - \log\left(\ln\left(\frac{2r_1\log_2(n) + 2\bar{C}_{ \alpha_1,d}}{\delta}\right)\right)
    \end{equation}
    where 
    \[
    \bar{C}_{ \alpha_1,d} = 1 + \log_2\left(\frac{ \alpha_1}{12\sqrt{d}}\right) -\log_2(C_0 ).
    \]
    
    We now bound $ \alpha_1 2^{-\kappa}$ by 
    
       \begin{equation}
    \begin{array}{ll}
     \alpha_1 2^{-r_1\log_2(n) - \bar{C}_{\alpha_1,d} + \log\left(\ln\left(\frac{2r_1\log_2(n) + 2 \bar{C}_{\alpha_1,d}}{\delta}\right)\right)} &=   \alpha_1 n^{-r_1}2^{-\bar{C}_{ \alpha_1,d}}\left(\ln\left(\frac{2r_1\log_2(n) + 2\bar{C}_{ \alpha_1,d}}{\delta}\right)\right)
     \\
         & =  \alpha_1 2^{-\bar{C}_{\alpha_1,d}} n^{-r_1}\left(\ln\left(2r_1\log_2(n) + 2\bar{C}_{ \alpha_1,d}\right)+\ln\left(\frac{1}{\delta}\right)\right)\\
         & \leq C_{ \alpha_1,d} \ln\left(\frac{1}{\delta}\right) n^{-r_1}\left(\ln\left(\ln(n) \right)\right)^{2r_1}
    \end{array}
    \end{equation}
    Since $\alpha_1$ is bounded from above, there is  $ C_{d}$ independent of $\alpha_1$ for which
    \begin{equation}
      \alpha_1 2^{-\kappa} = \alpha_1 2^{-r_1\log_2(n) - \bar{C}_{ \alpha_1,d} + \log\left(\ln\left(\frac{2r_1\log_2(n) + 2 \bar{C}_{ \alpha_1,d}}{\delta}\right)\right)} \leq C_{ d} \ln\left(\frac{1}{\delta}\right) n^{-r_1}\left(\ln\left(\ln(n) \right)\right)^{2r_1}
    \end{equation}
\end{proof}



\subsection{Supporting lemmas for Lemma \ref{lem:main_support_theorem_step2}}
\label{sec:ProofofLemAlphaHalf}

\subsubsection{Bounding the error between $T_0 {f_\ell}$ and $T_0\wtilde{f}_\ell$}\label{sec:Tf_Tf_tilde_err}
% --- begin Step2_BoundTf_Tf_tilde_lemma_map.tex ---
\usetikzlibrary{calc}

\begin{figure}[omitted: a diagram of this paper's results, not reproduced here]\label{tikz:lemTf-Tftilde_lemmas}\end{figure}





% \begin{figure}
%     \centering
%     %\resizebox{1.2\textwidth}{!}{%
%     \includegraphics[width = \textwidth]{Figures/Step2_BoundTf_Tf_tilde_lemma_map.pdf}
%     \caption{Road-map for proof of Lemma  \ref{lem:Tf_TtildeF_D}}
%     \label{tikz:lemTf-Tftilde_lemmas}
% \end{figure}% --- end Step2_BoundTf_Tf_tilde_lemma_map.tex ---
% --- begin Section_ProofStep2_BoundTf_Tf_tilde.tex ---

Back to Theorem \ref{thm:Step2} proof road-map see Figure \ref{tikz:thm33_lemmas}.


The main Lemma in the Section is Lemma \ref{lem:Tf_TtildeF_D}. A road-map for the proof appears in Figure~\ref{tikz:lemTf-Tftilde_lemmas}.

\begin{Lemma}\label{lem:Tf_TtildeF_D}
    Let the assumptions of Theorem \ref{thm:Step2} hold, and let $ f:H\simeq\RR^d\to\RR^{D-d} $ be a function such that its graph coincides with a neighborhood on the manifold $\MM$ (see Lemma \ref{lem:M_is_locally_a_fuinction_clean}). Let $\wtilde f:\RR^d\to\RR^{D-d} $ be the regression function defined in \eqref{eq:ftilde_def_general}. Denoting $\alpha = \maxangle(H, T_0 f)$, there is a constant $C_\tau$ large enough such that for $M = \frac{\tau}{\sigma}>C_\tau\sqrt{D \log D}$, and $\alpha \leq \frac{3}{2}\sqrt{C_M/M}$, where $C_M$ is from Theorem \ref{thm:Step1}, we have 
    \[
	\maxangle(T_0 f, T_0 {\wtilde{f}}) 
	\leq  
	 \frac{\alpha}{6} 
	\]
\end{Lemma}
\begin{proof}
From Lemma \ref{lem:bound_Df-Dfwtilde} we have 
\[
\|\DD_{\wtilde{f}}[0] - \DD_f[0]\| < \frac{\alpha}{6}
\]
From Lemma \ref{lem:Diff_Tf_bound} we have that 
    \[
	\maxangle(T_0 f, T_0 {\wtilde{f}}) 
	\leq   \frac{\alpha}{6}
	\]
	and the proof is concluded. 
\end{proof}


\begin{Lemma}\label{lem:bound_Df-Dfwtilde}
 Let the assumptions of Theorem \ref{thm:Step2} hold, and let $ f:H\simeq\RR^d\to\RR^{D-d} $ be a function such that its graph coincides with a neighborhood on the manifold $\MM$ (see Lemma \ref{lem:M_is_locally_a_fuinction_clean}). Let $\wtilde f:\RR^d\to\RR^{D-d} $ be the regression function   defined in \eqref{eq:ftilde_def_general}. 
 For any $\eps>0$, denoting $\alpha = \maxangle(H, T_0 f)$, there is a constant $C_\tau$ large enough such that for $M = \frac{\tau}{\sigma} >C_\tau\sqrt{D \log D}$, and $\alpha \leq \frac{3}{2} \sqrt{C_M/M}$, where $C_M$ is from Theorem \ref{thm:Step1}, we have 
    \[
    \|\DD\wtilde{f}(0) - \DD f(0)\| < \eps \alpha
    \]
\end{Lemma}
\begin{proof}
We reiterate the definition of $\Omega(x):H \simeq \RR^d \to 2^{\RR^{D-d}}$ from \eqref{eq:Omega_def}
\begin{equation*}
\Omega(x) = (x + H^\perp) \cap  \MM_\sigma
,    
\end{equation*}
where $x + H^\perp = \{x + y ~|~ y\in H^\perp\}$.
Next, denoting $S_{D-d}$ to be the $D-d$ dimensional unit sphere, we define $g(x, \theta):H\times S_{D-d}\to \RR$ the maximal length from $f(x)$ in the direction $\theta$ that is inside $\Omega(x)$.
Explicitly,
\begin{equation}\label{eq:g_def}
g(x, \theta) = \max\{y\in \RR ~|~ (x,f(x) + y\cdot \theta)_H  \in \Omega(x)\}
.    
\end{equation}
Note, that the farthest point from $f(x)$ in $\Omega(x)$ at each direction $\theta$, by which we define $g(x, \theta)$, belongs to $\partial\MM_\sigma$ (the boundary of the domain $\MM_\sigma$), and is therefore exactly $\sigma$ away from some point on the manifold itself.
Since we are viewing the manifold locally as the graph of the function $f:H\to H^\perp$ we denote this point by $(\wtilde x_\theta, f(\wtilde x_\theta(x)))_H$.
Explicitly, $\wtilde{x}_\theta:H \to H$ is such that
\begin{equation}\label{eq:x_tilde_def}
    (x ,f(x) + g(x, \theta) \cdot \theta)_H^T  = (\wtilde x_\theta(x), f(\wtilde x_\theta(x)))_H^T + \sigma \vec N_\theta (x)
,\end{equation}
where $\vec N_\theta \in \RR^D$ is perpendicular to $T_{\wtilde x_\theta}f$.
We introduced the definition of $\wtilde{x}_\theta$ here as it will be pivotal in the proofs of Lemmas \ref{lem:g_bond_weak} and \ref{lem:bound_nablag} upon which the the current proof relies.
Furthermore, we wish to stress here that by Lemma \ref{lem:bounding_stuff} $\wtilde x_\theta(0) \leq \sigma \sin (\alpha + 3 \sqrt{\sigma/\tau}\alpha)$.
Therefore, for sufficiently large $M = \frac{\tau}{\sigma}$, $\wtilde x_\theta(0)$ is within the domain of definition of the function $f$ which by Lemma \ref{lem:M_is_locally_a_fuinction_clean} is of radius of at least $c_{\pi/4}\tau$.

Next, by the definitions and Eq. \eqref{eq:g_def} and \eqref{eq:Omega_def} we have that 
\begin{equation*}
\wtilde{f}(x) - f(x) = \frac{\int\limits_{y\in\Omega(x)} y dy}{\int\limits_{y\in\Omega(x)} dy} 
.\end{equation*}
Since $\Omega \subset \RR^D$ is perpendicular to $H\simeq \RR^d$ we get that $\Omega\simeq \RR^{D-d}$.
Thus, by change of variables we can breakdown the integrals over $\Omega(x)$ to a radial component $r$ and directions on the $(D-d-1)$-dimensional sphere.
Explicitly,
\begin{equation*}
\wtilde{f}(x) - f(x) = 
\frac{\int\limits_{S_{D-d-1}}\int\limits_0^{g(x,\theta)} \theta r r^{D-d-1} drd\theta}{\int\limits_{S_{D-d-1}}\int\limits_0^{g(x,\theta)} r^{D-d-1} drd\theta}  = 
\frac{(D-d)\int\limits_{S_{D-d-1}}\theta g(x,\theta)^{D-d+1} d\theta}{(D-d+1)\int\limits_{S_{D-d-1}}g(x,\theta)^{D-d} drd\theta}     
,\end{equation*}
where $dr$ is the measure over the radial component, $r^{D-d-1}$ is the Jacobian introduced by the change of variables and $d\theta$ is the measure over the $(D-d-1)$-dimensional sphere.
For brevity we introduce the notation $\wtilde D = D - d$ and get
\begin{equation}\label{eq:f_tilde-f_integrals}
\wtilde{f}(x) - f(x) = 
\frac{\int\limits_{\mathbb{S}_{\wtilde D-1}}\int\limits_0^{g(x,\theta)} \theta r r^{\wtilde D-1} drd\theta}{\int\limits_{\mathbb{S}_{\wtilde D-1}}\int\limits_0^{g(x,\theta)} r^{\wtilde D-1} drd\theta}  = 
\frac{\wtilde D \int\limits_{\mathbb{S}_{\wtilde D-1}}\theta g(x,\theta)^{\wtilde D+1} d\theta}{(\wtilde D+1)\int\limits_{\mathbb{S}_{\wtilde D-1}}g(x,\theta)^{\wtilde D} drd\theta}     
.\end{equation}
 
Next, by taking the differential of that expression with respect to $x$ we have that 
\begin{multline*}
    \DD_{\wtilde{f}}[x] - \DD_f[x] =
\frac{\wtilde D}{(\wtilde D+1)}\left(
\frac{(\wtilde D+1)\int\limits_{\mathbb{S}_{\wtilde D-1}}\theta g(x,\theta)^{\wtilde D} \nabla_x g^T d\theta}
{\int\limits_{\mathbb{S}_{\wtilde D-1}}g(x,\theta)^{\wtilde D} d\theta} \right.\\
-\left.
\frac{\wtilde D\int\limits_{\mathbb{S}_{\wtilde D-1}}\theta g(x,\theta)^{\wtilde D+1} d\theta \int\limits_{\mathbb{S}_{\wtilde D-1}} g(x,\theta)^{\wtilde D-1}\nabla_x g^T d\theta}
{\left(\int\limits_{\mathbb{S}_{\wtilde D-1}}g(x,\theta)^{\wtilde D} d\theta\right)^2}     
\right)
,\end{multline*}
where $\nabla_x g$ stands for the gradient of $g(\theta, x)$ with respect to the $x$ variables only.
As can be seen in the above equation there is a multiplicative factor of size $\approx (D - d)$ for both summends.
In order to deal with this obstacle we wish to utilize the fact that in high dimensions most of the volume of a sphere is concentrated near an equator.
Thus, we split the domain into two different regions that we will deal with separately (we will use this trick in few of the other proofs). 
In case $D-d$ is small, then the following computations can be done without splitting the domain into two regions (one near the equator and the the second being the remaining cap), and include the $D-d$ factor in the constant that will be cancelled by $M$.
Therefore, we assume without losing the generality of our claim that $D-d = \wtilde D > 3$.
For any direction/unit vector $\vec{z}\in S_{\wtilde D}$ denote $\vec{z}^T\theta = z$ and 
\begin{equation}\label{eq:Omegas_def}
    \begin{aligned}
    \Omega_1 &= \{\theta~|~0\leq \vec{z}^T\theta \leq \xi\}\\
    \Omega_2 &= \{\theta~|~\vec{z}^T\theta > \xi\}
    \end{aligned}
\end{equation}
For some $\xi$ to be chosen later. Using the above notation, we have 
\begin{multline}\label{eq:zDftilde}
    \vec{z}^T\cdot  (\DD_{\wtilde{f}}[x] - \DD_f[x] )=
\frac{\wtilde D}{(\wtilde D+1)}\left(
\underbrace{\frac{(\wtilde D+1)\int\limits_{\mathbb{S}_{\wtilde D-1}}z g(x,\theta)^{\wtilde D} \nabla_x g^T d\theta}
{\int\limits_{\mathbb{S}_{\wtilde D-1}}g(x,\theta)^{\wtilde D} drd\theta}  }_{I} \right.  \\
-\left.
\underbrace{\frac{\wtilde D\int\limits_{\mathbb{S}_{\wtilde D-1}}z g(x,\theta)^{\wtilde D+1} d\theta \int\limits_{\mathbb{S}_{\wtilde D-1}} g(x,\theta)^{\wtilde D-1}\nabla_x g^T d\theta}
{\left(\int\limits_{\mathbb{S}_{\wtilde D-1}}g(x,\theta)^{\wtilde D} drd\theta\right)^2}}_{II}     
\right)
\end{multline}
First we treat part $(I)$ of Eq. \eqref{eq:zDftilde} by splitting the domain into $\Omega_1$ and $\Omega_2$ of Eq. \ref{eq:Omegas_def}.

\begin{align*}
    (I) &= \frac{(\wtilde D+1)\left(\int\limits_{\Omega_1}z g(x,\theta)^{\wtilde D} \nabla_x g^T d\theta + \int\limits_{\Omega_2}z g(x,\theta)^{\wtilde D} \nabla_x g^T d\theta \right) }
{\int\limits_{\mathbb{S}_{\wtilde D-1}}g(x,\theta)^{\wtilde D} d\theta} \\
&\leq\frac{(\wtilde D+1)\left(\xi\int\limits_{\Omega_1} g(x,\theta)^{\wtilde D} \nabla_x g^T d\theta + \int\limits_{\Omega_2} g(x,\theta)^{\wtilde D} \nabla_x g^T d\theta \right)}
{\int\limits_{\mathbb{S}_{\wtilde D-1}}g(x,\theta)^{\wtilde D} d\theta}
\end{align*}

From \ref{lem:g_bond_weak} we have that $0 < \sigma \leq g(0,\theta)\leq \sigma + 4\sigma\alpha^2$, and from Lemma \ref{lem:bound_nablag} we have $\|\nabla_x g(0,\theta)\|\leq c_1 \sigma \alpha/\tau$. Then,
\begin{align*}
\norm{(I)} &\leq \frac{(\wtilde D+1)\left(\xi\int\limits_{\Omega_1} g(x,\theta)^{\wtilde D} \|\nabla_x g^T\| d\theta + \int\limits_{\Omega_2} g(x,\theta)^{\wtilde D} \|\nabla_x g^T\| d\theta \right)}
{\int\limits_{\mathbb{S}_{\wtilde D-1}}g(x,\theta)^{\wtilde D} d\theta} \\
&\leq  \frac{(\wtilde D+1)c_1 \sigma \alpha \left(\xi\int\limits_{\Omega_1} g(x,\theta)^{\wtilde D}  d\theta + \int\limits_{\Omega_2} g(x,\theta)^{\wtilde D}  d\theta \right)}
{\tau\int\limits_{\mathbb{S}_{\wtilde D-1}}g(x,\theta)^{\wtilde D} d\theta}
\\
&\leq  \frac{\xi(\wtilde D+1)c_1 \sigma \alpha}{\tau} + \frac{(\wtilde D+1)c_1 \sigma \alpha \int\limits_{\Omega_2} g(x,\theta)^{\wtilde D}  d\theta }
{\tau\int\limits_{\mathbb{S}_{\wtilde D-1}}g(x,\theta)^{\wtilde D} d\theta}
\\
&\leq  \frac{\xi(\wtilde D+1)c_1 \sigma \alpha}{\tau} + \frac{(\wtilde D+1)c_1 \sigma \alpha ( \sigma + 4\sigma\alpha^2)^{\wtilde D}\int\limits_{\Omega_2}  d\theta }
{\tau \sigma^{\wtilde D}\int\limits_{\mathbb{S}_{\wtilde D-1}} d\theta}
\end{align*}

Furthermore, using the following concentration of measure inequality (see e.g. \cite{milman2009asymptotic,Venkatesan2012high}
\begin{equation}\label{eq:Omega3_to_ball_ratio}
\frac{\int\limits_{\Omega_2}  d\theta}{\int\limits_{\mathbb{S}_{\wtilde D-1}}  d\theta} \leq \frac{2}{\xi\sqrt{\wtilde D-2}}e^{-(\wtilde D - 2 )\xi^2/2}
,\end{equation}
we have that,
\begin{align*}
    \norm{(I)} 
&\leq  \frac{\xi(\wtilde D+1)c_1 \sigma\alpha}{\tau} + \frac{(\wtilde D+1)c_1 \sigma\alpha (1 + 4\alpha^2)^{\wtilde D} }
{\tau}\frac{2}{\xi\sqrt{\wtilde D-2}}e^{-(\wtilde D - 2 )\xi^2/2}
\\
&\leq  \frac{\xi(\wtilde D+1)c_1 \sigma\alpha}{\tau} + \frac{c_2 \sigma\alpha \sqrt{\wtilde D+1}( 1 + 4\alpha^2)^{\wtilde D} }
{\tau}\frac{2}{\xi}e^{-(\wtilde D - 2 )\xi^2/2}
,\end{align*}
where $c_2 = 4c_1 \geq \frac{\sqrt{\wtilde D +1}}{\sqrt{\wtilde D - 2}} c_1$.

Since $M = \frac{\tau}{\sigma}>C_\tau\sqrt{D \log D}$, we choose $\xi = 2\sqrt{\log D/D}$ and we have

\begin{align*}
    \norm{(I)} 
&\leq \frac{2c_1 \alpha(\wtilde D+1)}{C_\tau D} + \frac{c_2 \alpha \sqrt{\wtilde D+1}( 1 + 4\alpha^2)^{\wtilde D} }
{C_\tau\sqrt{D \log D}}\frac{2}{\xi}e^{-(\wtilde D - 2 )\xi^2/2}
\\
&\leq \alpha\left(\frac{2c_1 }{C_\tau} + \frac{c_1 \sqrt{\wtilde D+1}( 1 + 4\alpha^2)^{\wtilde D} } 
{C_\tau\sqrt{D\log D }}\frac{4\sqrt{D}}{\sqrt{\log D}}e^{-\log D}\right)
\\
&\leq \alpha\left(\frac{2c_1 }{C_\tau} + \frac{2c_1 \sqrt{D}( 1 + 4\alpha^2)^{\wtilde D} }
{C_\tau D\log D }\right)
,\end{align*}
where the second inequality is true since  $(\wtilde D-2)/D \geq 1/2$.
Since $\alpha \leq \frac{3}{2} \sqrt{C_M/M} $, and we also have that $M>C_\tau\sqrt{D \log D}$ we have that 
$( 1 + 4\alpha^2)^{\wtilde D}$ is bounded for any $D$. Thus, we have that for $C_\tau$ large enough, 
\begin{align}
\norm{(I)} 
&\leq \frac{\eps}{2}\alpha
\end{align}
for any $D$.


Next we bound $\|(II)\|$ from \eqref{eq:zDftilde}:
\begin{align*}
\frac{\wtilde D\int\limits_{\mathbb{S}_{\wtilde D-1}}z g(x,\theta)^{\wtilde D+1} d\theta \int\limits_{\mathbb{S}_{\wtilde D-1}} g(x,\theta)^{\wtilde D-1}\nabla_x g^T d\theta}
{\left(\int\limits_{\mathbb{S}_{\wtilde D-1}}g(x,\theta)^{\wtilde D} d\theta\right)^2}
\end{align*}
First, we note that 
\begin{align*}
\norm{\frac{\int\limits_{\mathbb{S}_{\wtilde D-1}} g(x,\theta)^{\wtilde D-1}\nabla_x g^T d\theta}
{\int\limits_{\mathbb{S}_{\wtilde D-1}}g(x,\theta)^{\wtilde D} d\theta}} & \leq \frac{c_1 \sigma\alpha}{\tau} \frac{\int\limits_{\mathbb{S}_{\wtilde D-1}} g(x,\theta)^{\wtilde D-1} d\theta}
{\int\limits_{\mathbb{S}_{\wtilde D-1}}g(x,\theta)^{\wtilde D-1} d\theta} = \frac{c_1 \sigma \alpha}{\tau}.
\end{align*}
Thus, similarly to the way we bounded $(I)$, we split the domain to $\Omega_1$ and $\Omega_2$ of Eq. \eqref{eq:Omega_def} and achieve
\begin{align*}
\norm{(II)} &\leq \frac{\wtilde Dc_1 \sigma\alpha}{\tau}\frac{\int\limits_{\mathbb{S}_{\wtilde D-1}}z g(x,\theta)^{\wtilde D+1} d\theta}
{\int\limits_{\mathbb{S}_{\wtilde D-1}}g(x,\theta)^{\wtilde D} d\theta}
\\
&\leq \frac{\wtilde Dc_1 \sigma\alpha}{\tau} \frac{ \left(\xi\int\limits_{\Omega_1} g(x,\theta)^{\wtilde D+1}  d\theta + \int\limits_{\Omega_2} g(x,\theta)^{\wtilde D+1}  d\theta \right)}
{\int\limits_{\mathbb{S}_{\wtilde D-1}}g(x,\theta)^{\wtilde D} d\theta}
\\
&\leq  \frac{\xi\wtilde Dc_1 \sigma\alpha}{\tau} + \frac{\wtilde Dc_1 \sigma\alpha ( \sigma + 4\sigma \alpha^2)^{\wtilde D+1}\int\limits_{\Omega_2}  d\theta }
{\tau \sigma^{\wtilde D}\int\limits_{\mathbb{S}_{\wtilde D-1}} d\theta}
\\
&\leq  \frac{\xi\wtilde Dc_1 \sigma\alpha}{\tau} + \frac{\wtilde Dc_1 \sigma^2\alpha ( 1 + 4\alpha^2)^{\wtilde D+1}\int\limits_{\Omega_2}  d\theta }
{\tau \int\limits_{\mathbb{S}_{\wtilde D-1}} d\theta}
\end{align*}
From \eqref{eq:Omega3_to_ball_ratio} we have that 
\begin{align*}
\norm{(II)} &\leq  \frac{\xi\wtilde Dc_1 \sigma\alpha}{\tau} + \frac{\wtilde Dc_1 \sigma^2\alpha ( 1 + 4\alpha^2)^{\wtilde D+1}}
{\tau} \frac{2}{\xi\sqrt{\wtilde D-2}}e^{-(\wtilde D - 2 )\xi^2/2}
\end{align*}
Since $M = \frac{\tau}{\sigma}>C_\tau\sqrt{D\log D}$, and choosing again $\xi = 2\sqrt{\log D/D}$ we have

\begin{align*}
\norm{(II)} &\leq  \frac{2\wtilde Dc_1 \alpha}{C_\tau D} + \frac{\wtilde Dc_1 \sigma\alpha ( 1 + 4\alpha^2)^{\wtilde D+1}}
{C_\tau \sqrt{D\log D}} \frac{\sqrt{D}}{\sqrt{\log D} \sqrt{\wtilde D-2}}e^{-\log D}
\\
&\leq \frac{2c_1 \alpha}{C_\tau } + \frac{\alpha c_1 \sigma ( 1 + 4\alpha^2)^{\wtilde D+1}}
{C_\tau D\log D} \frac{1}{\sqrt{\wtilde D-2}}
\end{align*}
Since From Theorem \ref{thm:Step1} we have that $\alpha \leq \frac{3}{2} \sqrt{ C_M/M}$, and thus 
$( 1 + \alpha)^{\wtilde D}$ is bounded for any $D$. Thus, we have that for $C_\tau$ large enough, 
\begin{align}
\norm{(II)}&\leq \frac{\eps}{2}\alpha
\end{align}
for any $D$.

Thus, we have from \eqref{eq:zDftilde} that for any unit vector $z$, $\|\vec{z}( \DD_{\wtilde{f}}[x] - \DD_f[x])\| \leq \eps\alpha$
or 
\[
\| \DD_{\wtilde{f}}[x] - \DD_{f}[x]\|_{op} \leq \eps\alpha
\]
\end{proof}


\begin{Lemma}\label{lem:Diff_Tf_bound}
	Let $ f, \tilde f $ be functions from $ \RR^d $ to $ \RR^{D-d} $, and denote their differentials by $ \DD_f, \DD_{\wtilde f} $ respectively.
	Denote by $ T_0 f, T_0 \wtilde f $ the tangent planes of the graphs of $ f $ and $ \wtilde f $ respectively.
	Assume that
	\[
	\norm{\DD_f[0] - \DD_{\tilde f}[0]}_{op}\leq \varepsilon 
	.\]
	Then, for sufficiently small $ \varepsilon $
	\[
	\sin(\maxangle(T_0 f, T_0 \wtilde{f})) \leq  \varepsilon 
	.\]
\end{Lemma}

\begin{proof}
	By definition we know that
	\[
	T_0 f = \{ \DD_{f}[0]v ~|~ v\in \RR^d \} \subset \RR^D
	,\]
	and
	\[
	T_0 \wtilde f_\ell = \{ (v, \DD_{\wtilde f}[0]v) ~|~ v\in \RR^d \} \subset \RR^D
	.\] 
	Let $ L_1, L_2 $ be two linear spaces, denote by $ Q_{L_1, L_2} :L_1\to L_2^\perp $  the following operator
	\[
	Q_{L_1, L_2} = {v - P_{L_2}(v)}
	.\]
	By Lemma \ref{lem:Q_angle_between_flats} we know that 
	\[
	\sin(\maxangle\lrbrackets{L_1, L_2}) = \norm{Q_{L_1, L_2}}_{op}
	.\]
	We now turn to look at the operator $ Q_{T_0 f, T_0 \wtilde f } $ operating on some vector $ v\in T_0 f $.
	\begin{align*}
	\norm{Q_{T_0 f, T_0 \wtilde f}(v,\DD_{f}[0]v )} 
	&\leq 
	\norm{(v,\DD_{f}[0]v ) - (v,\DD_{\wtilde f}[0]v )}\\
	&\leq 
	\norm{\DD_{f}[0] - \DD_{\wtilde f}[0]}_{op}\norm{(v,\DD_{f}[0]v)}
	,\end{align*}
	and thus,
	\[
	\sin (\angle_{max}(T_0 f, T_0 \wtilde f ))  = \|Q_{T_0 f, T_0 \wtilde f}\|_{op} \leq \|\DD_{f}[0] - \DD_{\wtilde f}[0]\|_{op} \leq  \varepsilon  
	\]
\end{proof}

\begin{Lemma}\label{lem:Q_angle_between_flats}
    Let $L_1$ and $L_2$ be two linear subspaces of $\RR^D$. Denote $Q_{L_1,L_2}:L_1 \rightarrow L_2^\perp$ defined as
    \[
    Q_{L_1,L_2}(v) = v - P_{L_2}(v)
    ,\]
    where $ P_{L_2} $ is the projection onto $ L_2 $.
    Then, 
    \[
    \sin \angle_{max}(L_1,L_2) = \|Q_{L_1,L_2}\|_{op}
    .\]
\end{Lemma}
\begin{proof}
	Recalling Definition \ref{def:principal_angles_clean}, the Principal Angles $ \beta_i $ between  $ L_1, L_2 $ and their corresponding pairs of vectors $ u_i\in L_1, w_i\in L_2 $ are defined as
	\begin{equation*}
	\begin{array}{ll}
	u_1, w_1 \defeq \argmin\limits_{\substack{u\in L_1, w\in L_2 \\ \norm{u}=\norm{w} =1}}\arccos\left(\abs{\langle u, w \rangle}\right), &
	\beta_1 \defeq \angle(u_1, w_1)
	\end{array}    
	,\end{equation*}
	and for $i > 1$
	\begin{equation*}
	\begin{array}{ll}
	u_i, w_i \defeq \argmin\limits_{\substack{u\perp\UU_{i-1}, w\perp\WW_{i-1} \\ \norm{u}=\norm{w} =1}}\arccos\left(\abs{\langle u, w \rangle}\right), &
	\beta_i \defeq \angle(u_i, w_i)
	\end{array}
	,\end{equation*}
	where $$\UU_i \defeq Span\{u_j\}_{j=1}^i ~,~ \WW_i \defeq Span\{w_j\}_{j=1}^i.$$
	We now wish to show that for all $ i $ we can choose
	\[
	w_i = \frac{P_{L_2}(u_i)}{\norm{P_{L_2}(u_i)}}
	.\]
	Since the definition is inductive so will be our proof.
	\paragraph*{Basis of the induction $ i=1 $:}
	We first denote
	\[
	v_1 = P_{L_2(u_1)}
	.\]
	Note that,
	\[
	\beta_1 = \angle(u_1, w_1) = \angle(u_1, \norm{v_1} w_1 )
	,\]
	and by the minimization problem defining $ \beta_1 $ we know that
	\[
	\beta_1 \leq \angle(u_1, v_1)
	.\]
	Then, since the projection onto a linear space minimizes the Least-Squares norm we get
	\begin{align*}
	\lrangle{u_1 - v_1, u_1 - v_1} 
	&\leq 
	\lrangle{u_1 - \norm{v_1}w_1, u_1 - \norm{v_1}w_1}\\
	\lrangle{u_1,u_1} - 2\lrangle{u_1, v_1} + \lrangle{v_1,v_1} 
	&\leq
	 \lrangle{u_1,u_1} - 2\lrangle{u_1, \norm{v_1}w_1} + \lrangle{\norm{v_1}w_1,\norm{v_1}w_1} \\
	 - 2\lrangle{u_1, v_1} + \norm{v_1}^2 
	 &\leq
	 - 2\lrangle{u_1, \norm{v_1}w_1} + \norm{v_1}^2 \\
	 \lrangle{u_1, \norm{v_1}w_1} 
	 &\leq
	 \lrangle{u_1, v_1}\\
	 \lrangle{u_1, w_1} 
	 &\leq
	 \lrangle{u_1, \frac{v_1}{\norm{v_1}}}
	 \\
	 \beta_1 =
	 \angle(u_1, w_1) =
	 \arccos(\lrangle{u_1, w_1}) 
	 &\geq
	 \arccos(\lrangle{u_1, \frac{v_1}{\norm{v_1}}}) = \angle(u_1, v_1)
	.\end{align*}
	Thus, 
	\[
	\angle(u_1, v_1) = \beta_1
	,\]
	and we can choose $ w_1 = v_1 $.
	\paragraph*{The induction step:}
	Now we assume that for all $ i\leq j $
	\[
	w_i = \frac{v_i}{\norm{v_i}}
	,\]
	where
	\[
	v_i = P_{L_2}(u_i)
	.\]
	And, we wish to show that 
	\[
	\angle(u_{j+1}, v_{j+1}) = \beta{j+1}
	,\]
	where the fact that $ \beta_{j+1} \leq  \angle(u_{j+1}, v_{j+1})$ results directly from the definition of $ \beta_{j+1} $.
	
	We first note that since
	\[
	u_{j+1}\perp\UU_j
	,\]
	we have
	\[	P_{L_2}(u_{j+1})\notin Span\{P_{L_2}(u_i)\}_{i=1}^j = \WW_j
	,\]
	thus, 
	\[
	v_{j+1} \in \WW_j^\perp  
	.\]
	From here on we can repeat the same argument as in the basis of the induction, just replacing $ L_1, L_2 $ with $ \UU_j^\perp, \WW_j^\perp $ respectively.
\end{proof}


\begin{Lemma}\label{lem:g_bond_weak}
Let the conditions of Lemma \ref{lem:bound_Df-Dfwtilde} hold.
Let $g(x, \theta)$ be as defined in Equation \eqref{eq:g_def}.
Then, for $\alpha$ smaller than some constant and $M$ larger then some constant, 
\[
\sigma \leq g(0,\theta) \leq \sigma +  4\sigma \alpha^2
\]
\end{Lemma}

\begin{proof}
    
    Since $(0,f(0))_H \in \MM$, any point $p\in \RR^D$ such that $\norm{p - (0,f(0))_H}\leq\sigma$ belongs to $\MM_\sigma$, the $\sigma$-tubular neighborhood of $\MM$. 
    In particular, this is also true for $p\in H^\perp\subset \RR^D$, and thus, we obtain the lower bound 
    \begin{equation*}
     \sigma \leq g(0,\theta)   
    ,\end{equation*}
    as by the definition of \eqref{eq:g_def} we have $g(0,\theta) = \max_{p\in (\textrm{Span}\{\theta\}\cap\MM_\sigma)} \norm{p - (0, f(0))_H}$.
    
    From Lemma \ref{lem:bounding_stuff} we have that 
    \[
    \|f(\wtilde{x}_\theta(0)) - f(0)\| \leq 2\sigma \sin (\alpha + 3\sqrt{\frac{\sigma}{\tau}}\alpha) \tan \alpha,
    \]
    where $\wtilde{x}_\theta$ defined in \eqref{eq:x_tilde_def}.
    
    Since $(0, f(0))_H + \theta g(0,\theta)$ is at distance  $\sigma$ from $(\wtilde{x}_\theta(0), f(\wtilde{x}_\theta(0)))_H$, and denoting $\theta\in H^\perp\subset \RR^D$ by $(0, \bar \theta)_H$ we have that
    \[
    \sigma  =
    \norm{(0, f(0) + \bar \theta g(0, \theta)  )_H -(\wtilde x_{\theta}(0), f(\wtilde x_\theta(0))_H}
    = \sqrt{\norm{\wtilde x_\theta}^2 + \norm{f(\wtilde x_\theta(0) - f(0)}^2 + g(0,\theta)^2 },\]
    and thus,
    \[
    g(0,\theta) \leq \sigma + 2\sigma \sin (\alpha + 3\sqrt{\frac{\sigma}{\tau}}\alpha) \tan \alpha,
    .\]
    Then, for $\alpha$ smaller than some constant and $M$ larger than some constant we have 
    \[
    g(0,\theta) \leq \sigma + 4\sigma \sin^2 \alpha,
    \] 
    or,
    \[
    g(0,\theta) \leq \sigma + 4\sigma \alpha^2,
    \] 
\end{proof}
\begin{Lemma}\label{lem:bound_nablag}
Let the conditions of Lemma \ref{lem:bound_Df-Dfwtilde} hold.
Let $g(x, \theta)$ be as defined in Equation \eqref{eq:g_def}. Then, 
\[
\|\nabla_x{g}(0,\theta)\| \leq C \cdot \frac{\sigma}{\tau}\alpha 
,\]
where $\nabla_x{g}$ denotes the gradient of $g(x, \theta)$ with respect to the $x$ variables only, and $C$ is some constant.
\end{Lemma}
\begin{proof}
Following the definition of $\wtilde{x}_\theta$ in \eqref{eq:x_tilde_def} and $g(x,\theta)$ in \eqref{eq:g_def}, we have the following equations that describe the connection between $x, \wtilde{x}_\theta$ and $g(x,\theta)$
\begin{equation*}
\left(\begin{array}{c}
     \wtilde{x}_\theta \\
     f(\wtilde{x}_\theta) 
\end{array}\right)
+
\sigma \vec N(x,\wtilde{x}_\theta,\theta) = 
\left(\begin{array}{c}
     x \\
     f(x) 
\end{array}\right)_H
+ 
\left(\begin{array}{c}
     0 \\
     \bar \theta
\end{array}\right)_H g(x,\theta)
,\end{equation*}
where $\theta\in \RR^D$ is written as $(0, \bar \theta)_H$, and $ \vec N(x,\wtilde{x}_\theta,\theta)\in \RR^D$ is some unit vector perpendicular to $T_{\wtilde{x}_\theta}f$.
Explicitly, 
\begin{equation}\label{eq:c_2}
 \vec N(x,\wtilde{x}_\theta,\theta) \perp T_{\wtilde{x}_\theta}f
,\end{equation}
and 
\begin{equation}\label{eq:c_3}
\|\vec N(x,\wtilde{x}_\theta,\theta) \| = 1
.\end{equation}
Alternatively, we can write,
\begin{equation}\label{eq:c_1}
\sigma \vec N(x,\wtilde{x}_\theta,\theta) = 
\left(\begin{array}{c}
     x \\
     f(x) 
\end{array}\right)_H
-
\left(\begin{array}{c}
     \wtilde{x}_\theta \\
     f(\wtilde{x}_\theta) 
\end{array}\right)_H
+ 
\left(\begin{array}{c}
     0 \\
     \bar \theta
\end{array}\right)_H g(x,\theta)
\end{equation}


Taking the norm of \eqref{eq:c_1} and using \eqref{eq:c_3} we have  

\begin{equation*}
\left\|\left(\begin{array}{c}
     x \\
     f(x) 
\end{array}\right)
-
\left(\begin{array}{c}
     \wtilde{x}_\theta \\
     f(\wtilde{x}_\theta) 
\end{array}\right)
+ 
\left(\begin{array}{c}
     0 \\
     \bar  \theta
\end{array}\right) g(x,\theta)\right\|^2 = \sigma^2
\end{equation*}
or,
\[
g^2 + 2g \left(\begin{array}{c}
     x - \wtilde{x}_\theta \\
     f(x) - f(\wtilde{x}_\theta) 
\end{array}\right) ^T
\left(\begin{array}{c}
     0 \\
     \bar \theta
\end{array}\right)
+
\left\|\left(\begin{array}{c}
     x - \wtilde{x}_\theta \\
     f(x) - f(\wtilde{x}_\theta) 
\end{array}\right)\right\|^2 
-
\sigma^2 =0
\]
or,
\[
g^2 + 2g (f(x) - f(\wtilde{x}_\theta))^T \bar \theta 
+
\|x - \wtilde{x}_\theta\|^2 
+
\| f(x) - f(\wtilde{x}_\theta) \|^2
-
\sigma^2 =0
\]
the two solutions are 
\begin{multline}\label{eq:sol_g_pm}
    g_{\pm}(x,\theta) =\\ -f(x)^T \bar \theta + f(\wtilde{x}_\theta)^T \bar \theta \pm
    \sqrt{\sigma^2 + (f(x)^T \bar \theta - f(\wtilde{x}_\theta)^T \bar \theta)^2 - \|x- \wtilde{x}_\theta\|^2 - \|f(x)- f(\wtilde{x}_\theta) \|^2}
\end{multline}

From Lemma \ref{lem:bounding_stuff}, for $\alpha$ smaller than some constant, we have that
\[ 
\|0- \wtilde{x}_\theta(0)\|^2 + \|f(0)- f(\wtilde{x}_\theta(0)) \|^2 \leq \sigma^2
\]
and thus, the solutions of Eq. \eqref{eq:sol_g_pm} at $x=0$ are $g_{-}(0, \theta) < 0$ and $g_{+}(0, \theta) > 0$.
Therefore, from continuity we get that there is a neighborhood of $x=0$ such that the only non-negative solution is
\begin{multline}\label{eq:sol_g}
    g(x,\theta) =\\ -f(x)^T \bar \theta + f(\wtilde{x}_\theta)^T \bar \theta + 
    \sqrt{\sigma^2 + (f(x)^T \bar \theta - f(\wtilde{x}_\theta)^T \bar \theta)^2 - \|x- \wtilde{x}_\theta\|^2 - \|f(x)- f(\wtilde{x}_\theta) \|^2}
.\end{multline}
In addition, from the definition of $g$ the only valid solution is the non-negative one which appears on Eq. \eqref{eq:sol_g}.
Thus, denoting 
\begin{equation}\label{eq:Delta}
\Delta = \sigma^2 + (f(x)^T \bar \theta - f(\wtilde{x}_\theta)^T \bar \theta)^2 - \|x- \wtilde{x}_\theta\|^2 - \|f(x)- f(\wtilde{x}_\theta) \|^2,
\end{equation}
we have that near $x=0$
\begin{multline}\label{eq:nabla_g}
    \nabla_x{g}(x,\theta) =\\ -\DD_f[x]^T\bar\theta + J_{\wtilde{x}_\theta} \DD_f[\wtilde{x}_\theta]^T \bar \theta 
    +
    \frac{1}{\sqrt{\Delta}}
    \left( 
     (f(\wtilde{x}_\theta)^T \bar \theta - f(x)^T \bar \theta)\left(J_{\wtilde{x}_\theta}\DD_f[\wtilde{x}_\theta]^T \bar \theta - \DD_f[x]^T \bar \theta\right) 
    \right. \\
    \left.- 
    (I_d-J_{\wtilde{x}_\theta})(x- \wtilde{x}_\theta) 
    - 
    (\DD_f[x]^T - J_{\wtilde{x}_\theta}\DD_f[\wtilde{x}_\theta]^T)(f(x)- f(\wtilde{x}_\theta))  \right)
,\end{multline}
where $J_{\wtilde x_\theta(x)} = \DD_{\wtilde x_\theta (x)}[x]$ is the Jacobi matrix of the function $\wtilde x_\theta (x)$, and $I_d$ is the $d$-dimensional identity matrix.
Alternatively, we can write
\begin{multline}\label{eq:nabla_g2}
    \nabla_x{g}(x,\theta) =\\ (J_{\wtilde{x}_\theta} \DD_f[\wtilde{x}_\theta]^T - \DD_f[x]^T ) \bar\theta 
    +
    \frac{1}{\sqrt{\Delta}}
    \left( 
     (f(\wtilde{x}_\theta)^T\bar\theta - f(x)^T\bar\theta ) \left(J_{\wtilde{x}_\theta}\DD_f[\wtilde{x}_\theta]^T - \DD_f[x]^T \right)\bar\theta 
    \right. \\
    \left.- 
    (I_d-J_{\wtilde{x}_\theta})(x- \wtilde{x}_\theta) 
    - 
    (\DD_f[x]^T - J_{\wtilde{x}_\theta}\DD_f[\wtilde{x}_\theta]^T)(f(x)- f(\wtilde{x}_\theta))  \right)
\end{multline}


Next, using Lemma \ref{lem:bound_J}, Lemma \ref{lem:Dfx-Dfxtilde} and Lemma \ref{lem:bounding_stuff}, we bound 
\begin{equation}\label{eq:Df-JDfwtilde}
\begin{aligned}
     \|\DD_f[0]^T - J_{\wtilde{x}_\theta} \DD_f[\wtilde{x}_\theta(0)]^T \| &= \|\DD_f[0]^T -\DD_f[\wtilde{x}_\theta(0)]^T+\DD_f[\wtilde{x}_\theta(0)]^T - J_{\wtilde{x}_\theta(0)} \DD_f[\wtilde{x}_\theta(0)]^T \| 
     \\
     &\leq \|\DD_f[0]^T -\DD_f[\wtilde{x}_\theta(0)]^T\|+\|\DD_f[\wtilde{x}_\theta(0)]^T - J_{\wtilde{x}_\theta(0)} \DD_f[\wtilde{x}_\theta(0)]^T \|
     \\
     &\leq \|\DD_f[0]^T -\DD_f[\wtilde{x}_\theta(0)]^T\|+\|I_d - J_{\wtilde{x}_\theta(0)} \|\|\DD_f[\wtilde{x}_\theta(0)]^T \| 
     \\
     &= \OO(\frac{\sigma}{\tau}\sin \alpha) + \OO(\frac{\sigma}{\tau}) \|\DD_f[\wtilde{x}_\theta(0)]^T \|
     \\
     &= \OO(\frac{\sigma}{\tau}\cdot\alpha) 
\end{aligned}
\end{equation}


Now, using \eqref{eq:nabla_g2} and the fact that $\norm{\bar \theta} = 1$ we get
\begin{align*}
    \|\nabla_x{g}(0,\theta)\|  &\leq\|\DD_f[0]^T - J_{\wtilde{x}_\theta}[0] \DD_f[\wtilde{x}_\theta(0)]^T\|  \\\
    &~~~
    +
    \frac{1}{\sqrt{\Delta}}
    \left( 
     \|f(\wtilde{x}_\theta(0)) - f(0)\| \|J_{\wtilde{x}_\theta}[0]\DD_f[\wtilde{x}_\theta(0)]^T - \DD_f[0]^T \| 
    \right. \\
    &~~~ \left.+ 
    \|I_d-J_{\wtilde{x}_\theta}[0]\|\|\wtilde{x}_\theta(0)\| 
    + 
    \|\DD_f[0]^T - J_{\wtilde{x}_\theta}[0]\DD_f[\wtilde{x}_\theta(0)]^T\|\|f(0)- f(\wtilde{x}_\theta(0))\|\right)
\end{align*}
From Lemmas \ref{lem:bound_J} and \ref{lem:bounding_stuff} we know that $\|I_d - J_{\wtilde x_\theta}[0]\| =\OO(\frac{\sigma}{\tau})$, $\|\wtilde x_\theta(0)\|\leq \sigma \sin(\alpha + 3\sqrt{\sigma/\tau} \alpha)$, and that $\|f(\wtilde x_\theta(0)) - f(0)\| \leq 2\sigma \sin(\alpha + 3\sqrt{\sigma/\tau}\alpha)\tan(\alpha)$.
In other words, for $\alpha$ smaller than some constant we can say that $\|I_d - J_{\wtilde x_\theta}[0]\| = \OO(\sigma/\tau)$, $\|\wtilde x_\theta(0)\|\leq \OO(\sigma \alpha)$, and that $\|f(\wtilde x_\theta(0)) - f(0)\| \leq \OO(\sigma \alpha^2)$.
Combining this with \eqref{eq:Df-JDfwtilde} as well, we have
\begin{align*}
    \|\nabla_x{g}(0,\theta)\| &\leq\OO(\frac{\sigma}{\tau}\cdot\alpha)
    +
    \frac{1}{\sqrt{\Delta}}
    \left( 
    \OO(\frac{\sigma^2}{\tau}\cdot\alpha^3)
    + 
    \OO(\frac{\sigma^2}{\tau}\cdot\alpha) 
    +
    \OO(\frac{\sigma^2}{\tau}\cdot\alpha^3)\right)
\end{align*}
Since Lemma \ref{lem:Dxtildew+Dxw} gives us 
\[
 \Delta(0,\wtilde x_\theta (0)) \geq \frac{1}{2} \sigma^2
,\]
we have that $\frac{1}{\sqrt{\Delta}}\leq \sqrt{2}/\sigma$, and 
\begin{equation*}
    \|\nabla_x{g}(x,\theta)\| =  \OO(\frac{\sigma}{\tau}\cdot\alpha)
.\end{equation*}
\end{proof}

\begin{Lemma}\label{lem:bound_J}
Let the conditions of Lemma \ref{lem:bound_nablag} hold.
Let $g(x, \theta), \wtilde x(x)$ be as defined in Equation \eqref{eq:g_def} and \eqref{eq:x_tilde_def} respectively, let $I_d:\RR^d\to\RR^d$ denote the identity matrix and let $J_{\wtilde x_\theta}$ denote the differential of $\wtilde x(x)$ with respect to $x$. Then, 
\begin{equation*}
\|I_d-J_{\wtilde{x}_\theta}[0]\| = \OO(\frac{\sigma}{\tau})
.\end{equation*}    
\end{Lemma}
\begin{proof}
We begin by reiterating equations \eqref{eq:c_1},\eqref{eq:c_2}, and \eqref{eq:c_3}.
Namely we have
\begin{equation*}
\sigma \vec N(x,\wtilde{x}_\theta,\theta) = 
\left(\begin{array}{c}
     x \\
     f(x) 
\end{array}\right)
-
\left(\begin{array}{c}
     \wtilde{x}_\theta \\
     f(\wtilde{x}_\theta) 
\end{array}\right)
+ 
\left(\begin{array}{c}
     0 \\
     \bar \theta
\end{array}\right) g(x,\theta)
,\end{equation*}
where
\begin{equation*}
 \vec N(x,\wtilde{x}_\theta,\theta) \perp T_{\wtilde{x}_\theta}f
,\end{equation*}
and 
\begin{equation*}
\|\vec N(x,\wtilde{x}_\theta,\theta) \| = 1
.\end{equation*}
Thus, there is a vector $v(x,\wtilde x_\theta)\in \RR^{D-d}$ with $\|v(x,\wtilde x_\theta)\| = 1$,
\[
\vec N(x,\wtilde{x}_\theta,\theta) = \frac{1}{\sqrt{\|\DD_f[\wtilde{x}_\theta] v(x,\wtilde x_\theta)\|^2 + 1}} 
\left(\begin{array}{c}
     -\DD_f[\wtilde{x}_\theta]^T v(x,\wtilde x_\theta) \\
     v(x,\wtilde x_\theta)
\end{array}\right)
\]
or, denoting $ w(x,\wtilde x_\theta) = \frac{\sigma}{\sqrt{\|\DD_f[\wtilde{x}_\theta]^T v(x,\wtilde x_\theta))\|^2 + 1}}  v$, we have 
\[
\sigma \vec N(x,\wtilde{x}_\theta) = 
\left(\begin{array}{c}
     -\DD_f[\wtilde{x}_\theta]^T  w(x,\wtilde x_\theta) \\
     w(x,\wtilde x_\theta)
\end{array}\right)
\]

Using this pronunciation of $\vec N$ we can rewrite the above equation as 
\begin{equation}\label{eq:c_1_w}
\left(\begin{array}{c}
     -\DD_f[\wtilde{x}_\theta]^T w(x,\wtilde x_\theta) \\
     w(x,\wtilde x_\theta)
\end{array}\right) = 
\left(\begin{array}{c}
     x \\
     f(x) 
\end{array}\right)
-
\left(\begin{array}{c}
     \wtilde{x}_\theta \\
     f(\wtilde{x}_\theta) 
\end{array}\right)
+ 
\left(\begin{array}{c}
     0 \\
     \bar \theta
\end{array}\right) g(x,\theta)
.\end{equation}
From Eq. \eqref{eq:sol_g} in the proof of Lemma \ref{lem:bound_nablag} we know that 
$g(x, \theta) = -f(x)^T\bar\theta + f(\wtilde x_\theta)^T\bar \theta + \sqrt{\Delta}$, near $x=0$, where $\Delta$ is defined in Eq. \eqref{eq:Delta}.
Combining this with the last $D-d$ equations  we get,
\begin{equation}\label{eq:def_w}
    w(x,\wtilde x_\theta) =\left(f(x) - f(\wtilde{x}_\theta) + \bar \theta \left(- f(x)^T \bar \theta + f(\wtilde{x}_\theta) ^T \bar \theta + 
    \sqrt{\Delta}\right)\right)
.\end{equation}
Looking at the first $d$ equations of \eqref{eq:c_1_w}, we have 
\[
-\DD_f[\wtilde{x}_\theta]^T w(x,\wtilde x_\theta) - x + \wtilde{x}_\theta = 0
.\]
Denoting the function 
\begin{equation} \label{eq:G_def}
G(x,\wtilde{x}_\theta) = -\DD_f[\wtilde{x}_\theta]^T w(x,\wtilde x_\theta) - x + \wtilde{x}_\theta    
,\end{equation}
we aim at using the Inverse Function Theorem (IFT) to compute $J_{\wtilde{x}_\theta}$.
First, we compute  $\DD^{x}_G $ and $\DD^{\wtilde{x}_\theta}_G$, the partial differentials of $G$ with respect to the variables $x$ and $\wtilde x_\theta$:
\begin{equation*}
    \DD^x_G[x, \wtilde x_\theta] = - \DD_f[\wtilde{x}_\theta]^T \DD^x_w[x, \wtilde x_\theta] - I_d
\end{equation*}
\begin{equation*}
    \DD^{\wtilde{x}_\theta}_G[x, \wtilde x_\theta] = - \HH f(\wtilde x_\theta)^w - \DD_f[\wtilde{x}_\theta]^T \DD^{\wtilde{x}_\theta}_w[x, \wtilde x_\theta]  + I_d
,\end{equation*}
where $\HH f(\wtilde x_\theta)^w \in \RR^{d\times d}$ is the tensor Hessian of $f(\wtilde x_\theta):\RR^d\to\RR^{D-d}$ projected onto the target direction $w\in \RR^{D-d}$; that is
\begin{equation}\label{eq:H_w_def}
\HH f(\wtilde x_\theta)^w =
\left(\frac{\partial\left(\DD_f[\wtilde{x}_\theta]^T\right)}{\partial \wtilde{x}_1} w \right|\left. \frac{\partial\left(\DD_f[\wtilde{x}_\theta]^T\right)}{\partial \wtilde{x}_2} w \right| \cdots \left|\frac{\partial\left(\DD_f[\wtilde{x}_\theta]^T\right)}{\partial \wtilde{x}_d} w \right)
.\end{equation}
Notice that $\DD_f[\wtilde x_\theta]^T\in \RR^{d\times D-d}$; therefore, $\partial_{\wtilde x_j}\DD_f[\wtilde x_\theta]^T\in \RR^{d\times D-d}$ and $\partial_{\wtilde x_j}\DD_f[\wtilde x_\theta]^T w\in \RR^{d}$.

Next, using the IFT we have that 
\[
J_{\wtilde{x}_\theta} = -(\DD^{\wtilde{x}_\theta}_{G})^{-1} \DD^x_G = \left( I_d - \DD_f[\wtilde{x}_\theta]^T \DD^{\wtilde{x}_\theta}_ w + \HH f(\wtilde x_\theta)^w  \right)^{-1} \left(I_d + \DD_f[\wtilde{x}_\theta]^T \DD^x_w \right)
,\]
and thus
\[
J_{\wtilde{x}_\theta}[0] = (A(I_d-A^{-1}B))^{-1}A = (I_d-A^{-1}B)^{-1}A^{-1}A =  (I_d-A^{-1}B)^{-1}
,\]
where
\begin{equation}\label{eq:A_def}
A = I_d +  \DD_f[\wtilde{x}_\theta(0)]^T \DD^x_w[0,\wtilde{x}_\theta(0)]    
,\end{equation}
\begin{equation}\label{eq:B_def}
B = \DD_f[\wtilde{x}_\theta(0)]^T (\DD^x_w[0,\wtilde{x}_\theta(0)] + \DD^{\wtilde{x}_\theta}_w[0,\wtilde{x}_\theta(0)] )+ \HH f(\wtilde x_\theta(0))^w
.\end{equation}
From Lemma \ref{lem:bound_A_inv_B} we have that $\|A^{-1}B\| \leq \OO(\frac{\sigma}{\tau}) \leq 1/2$ for $\sigma/\tau$ smaller than some constant, and thus, using the first order approximation of this term we get that there is a matrix
\[\mathcal{E} = \sum_{t=2}^{\infty}(A^{-1}B)^t\]
such that
\[
J_{\wtilde{x}_\theta} = (I_d - A^{-1}B)^{-1} = I_d +A^{-1}B + \mathcal{E}
,\]
and
with
\[
\|\mathcal{E}\| \leq \|A^{-1}B\| \sum_{t=1}^\infty \frac{1}{2^t} = \|A^{-1}B\|
.\]

Thus we have, 
\begin{equation}\label{eq:I-J}
\|I_d -J_{\wtilde{x}_\theta}\| \leq  2\|A^{-1}B\|  =  \OO(\frac{\sigma}{\tau})
\end{equation}

\end{proof}

\begin{Lemma}\label{lem:bound_A_inv_B}
Let the conditions of Lemma \ref{lem:bound_J} hold and let $A$ and $B$ be as defined in \eqref{eq:A_def} and \eqref{eq:B_def}.
Then,
    \[
    \|A^{-1}B\| \leq \OO(\frac{\sigma}{\tau})
    \]
\end{Lemma}
\begin{proof}
We begin by noting that 
\[
\|A^{-1}B\| \leq \|A^{-1}\|\|B\|
,\]
where
\begin{equation*}
A = I_d +  \DD_f[\wtilde{x}_\theta(0)]^T \DD^x_w[0,\wtilde{x}_\theta(0)]    
,\end{equation*}
\begin{equation*}
B = \DD_f[\wtilde{x}_\theta(0)]^T (\DD^x_w[0,\wtilde{x}_\theta(0)] + \DD^{\wtilde{x}_\theta}_w[0,\wtilde{x}_\theta(0)] )+ \HH f(\wtilde x_\theta(0))^w
.\end{equation*}
Moreover,
\[
A^{-1} = (I_d +  \DD_f[\wtilde{x}_\theta(0)]^T \DD^x_w[0, \wtilde x_\theta(0)])^{-1} = I_d + \sum_{t=1}^\infty(\DD_f[\wtilde{x}_\theta(0)]^T \DD^x_w[0, \wtilde x_\theta(0)])^t
.\]
From Lemma \ref{lem:Dxtildew+Dxw} we have that $\|\DD^x_w[0, \wtilde x_\theta(0)]\| = \OO(\sin \alpha)$, where we remind the reader that $\DD^x_w[x, \wtilde x_\theta]$ is the partial differential of $w(x, \wtilde x_\theta)$ with respect to the $x$ variables only.
And, from Lemma \ref{lem:bounding_stuff}we have that $\|\DD_f[\wtilde{x}_\theta(0)]\|_2 \leq \sin (\alpha + 3\sqrt{\frac{\sigma}{\tau}}\alpha)$ . Thus $\|\DD_f[\wtilde{x}_\theta]^T \DD^x_w[0, \wtilde x_\theta(0)]\| = \OO(\sin\alpha)$, and thus, for $\alpha$ smaller than some constant we have
\begin{equation}\label{eq:bound_A}
    \|A^{-1}\| = 1 + \OO(\sin^2\alpha)
\end{equation}

Furthermore, from Lemma \ref{lem:Dxtildew+Dxw} we also know that $(\DD^x_w[0, \wtilde x_\theta(0)] + \DD^{\wtilde{x}_\theta}_ w[0, \wtilde x_\theta(0)] ) = \OO(\frac{\sigma}{\tau} \sin\alpha)$ and so 
\[
\| \DD_f[\wtilde{x}_\theta(0)]^T (\DD^x_w[0, \wtilde x_\theta(0)] + \DD^{\wtilde{x}_\theta}_ w[0, \wtilde x_\theta(0)] )\| \leq \OO(\frac{\sigma}{\tau}\sin^2 \alpha)
.\]
Combining this bound with the fact that $\|\HH f(\wtilde x_\theta(0))^w\| \leq \OO(\frac{\sigma}{\tau}) $ shown in Lemma \ref{lem:bound_DDf} we have 
\begin{equation}\label{eq:bound_B}
\|B\| \leq  \|\DD_f[\wtilde{x}_\theta]^T (\DD^x_w + \DD^{\wtilde{x}_\theta}_ w )\|+\| \HH f(\wtilde x_\theta)^w\| \leq \OO(\frac{\sigma}{\tau}) +  \OO(\frac{\sigma}{\tau}\sin^2 \alpha)
.
\end{equation}

Finally, from  \eqref{eq:bound_A} and \eqref{eq:bound_B} we have that for $\alpha$ smaller than some constant
\[
\|A^{-1}B\| \leq \OO(\frac{\sigma}{\tau})
\]
\end{proof}

\begin{Lemma}\label{lem:bound_DDf}
Let the conditions of Lemma \ref{lem:bound_J} and let $\HH f(\wtilde x_\theta(0))^w$ be as defined in \eqref{eq:H_w_def}.
Then, we have
\[
\|\HH f(\wtilde x_\theta(0))^w\|_{op} = \frac{\sqrt{2}\sigma}{\tau} + \OO\left(\frac{\sigma\sin \alpha}{\tau}\right)
\]
\end{Lemma}
\begin{proof}
We denote the tensor Hessian of $f:\RR^d\to\RR^{D-d}$ at $\wtilde x_\theta(0)$ by $\HH f(\wtilde x_\theta(0)):\RR^d\times\RR^d \to \RR^{D-d}$. 
For brevity of notation, throughout this proof we will use $\HH$ instead of $\HH f(\wtilde x_\theta(0))$. For any chosen direction $u\in \RR^{D-d}$ (i.e., a unit vector), $\HH$ can be thought of as a function:
$\HH ^u:\RR^d\times\RR^d \to \RR$
defined as $\HH ^u(v_1,v_2) = \lrangle{\HH (v_1,v_2), u}$. 
We note that this definition is consistent with the definition of $\HH f(\wtilde x_\theta(0))^u$ in \eqref{eq:H_w_def}.
Given $v\in\RR^d$, $\norm{v} = 1$ we also define $\HH_v:\RR^d \to \RR^{D-d}$
as $\HH_v(\cdot) = \HH(v,\cdot)$.
Note, that for any $w\in \RR^{D-d}$ 
\[
{\|\HH^w\|_{op}} = {\sup\limits_{v_1, v_2\in \mathbb{S}^{d}}|\lrangle{\HH(v_1, v_2), w}|} 
\leq \sup\limits_{v_1, v_2\in \mathbb{S}^{d}}\|\HH(v_1, v_2)\|_2\norm{w}_2
= \norm{w}_2\sup_{v\in\mathbb{S}^d}\|\HH_v\|_{op}
,\]
where the right-most equality is true since $\HH$ is symmetric.

Thus, in essence, we need to bound $\|\HH_v\|_{op}$ for an arbitrary $v$. 
From the definition of $\HH_v$ we know that 
\[
\HH_v\ =  \lim_{t\to 0} \frac{\DD_f[\wtilde x_\theta(0)] - \DD_f[\wtilde x_\theta(0) + tv]}{t} 
.\]
Then, from Lemma \ref{lem:alpha-beta_x_no_func}  we have for small $t$
\[
\sin(\maxangle (T_{\wtilde x_\theta}f,T_{\wtilde x_\theta + tv}f))\leq \frac{t}{\tau}(1+\tan^2\beta)+\OO(t^2/\tau^2)
\]
where $\beta = \maxangle(T_{\wtilde x_\theta}f,H)$. 
Therefore, applying Lemma \ref{lem:norm_from_angle} we get 
\[
 \|\DD_f[\wtilde x_\theta(0)] - \DD_f[\wtilde x_\theta(0) + tv] \|_{op} \leq \frac{t}{\tau}(1+\tan^2\beta)(1+\sin \beta) + \OO(t^2/\tau^2)
,\]
and we get 
\[
\|\HH_v\| \leq  \frac{1}{\tau}(1+\tan^2\beta)(1+\sin \beta)
.\]
Furthermore, from Lemma \ref{lem:alpha_beta_D} we know 
\[
\beta \leq \alpha + 3\alpha\sqrt{\frac{\sigma}{\tau}}
,\]
and so,
\[
\|\HH_v\| \leq  \frac{1}{\tau}(1+\tan^2
\left(\alpha + 3\alpha\sqrt{\frac{\sigma}{\tau}}\right))(1+\sin \left(\alpha + 3\alpha\sqrt{\frac{\sigma}{\tau}}\right))
= \frac{1}{\tau} + \OO(\frac{\sin \alpha}{\tau})
.\]
Thus, we obtain
\begin{equation}\label{eq:norm_H^w_bound}
\|\HH^w\| \leq \|w\| \left(\frac{1}{\tau} + \OO(\frac{\sin \alpha}{\tau})\right)    
.\end{equation}


Hence, all we are left with is bounding  $\|w(0,\wtilde x_\theta (0))\|$. 
From Eq. \eqref{eq:def_w}, Lemma \ref{lem:bounding_stuff}, and Lemma \ref{lem:Dxtildew+Dxw}, we have 
\begin{align*}
    \|w(0,\wtilde x_\theta(0))\| &=\|f(0) - f(\wtilde{x}_\theta(0)) + \theta \left(- f(0)\cdot \theta + f(\wtilde{x}_\theta(0)) \cdot \theta + 
    \sqrt{\Delta(0,\wtilde x_\theta (0))}\right)\|
    \\
    &\leq \|f(0) - f(\wtilde{x}_\theta(0))\| + \left|(- f(0) + f(\wtilde{x}_\theta(0))) \cdot \theta\right| + 
    \sqrt{\Delta(0,\wtilde x_\theta (0))}
    \\
    &\leq 2\sigma \sin \alpha + 2\sigma \sin \alpha + \sqrt{2}\sigma 
    \\
    &\leq \sqrt{2}\sigma  + 4\sigma \sin \alpha 
\end{align*}

Since $\|w\| \leq \sqrt{2}\sigma + \OO(\sin \alpha)$ 
we have from Eq. \eqref{eq:norm_H^w_bound}
\[
\|\HH^w\| \leq \frac{\sqrt{2}\sigma}{\tau} + \OO(\frac{\sigma\sin \alpha}{\tau})
\]
\end{proof}


\begin{Lemma}\label{lem:Dxtildew+Dxw}
Let the conditions of Lemma \ref{lem:bound_J} hold, let $w(x, \wtilde x_\theta)$ be as defined in \eqref{eq:def_w}, $\wtilde x_\theta(x)$ be as defined in \eqref{eq:x_tilde_def}, and $\Delta$ as defined in \eqref{eq:Delta}. 
Denote by $\DD^x_w, \DD^{\wtilde{x}_\theta}_w$ the partial differentials of $w$ with respect to the variables $x$ and $\wtilde x_\theta$.
Then, for $\alpha$ smaller than some constant
\begin{equation}\label{eq:Delta_bound}
    \frac{1}{2} \sigma^2 \leq \Delta(0,\wtilde x_\theta (0)) \leq 2 \sigma^2
,\end{equation}
\begin{equation} \label{eq:Dws_bound}
\| \DD^{\wtilde{x}_\theta}_ w[0, \wtilde x_\theta(0)] \| \leq  \mathcal{O}(\sin \alpha) \qquad \|\DD^x_w[0, \wtilde x_\theta(0)]\| \leq \mathcal{O}(\sin \alpha)
,\end{equation}
and
\begin{equation}\label{eq:Dxtildew+Dxw_bound}
\| \DD^{\wtilde{x}_\theta}_w[0, \wtilde x_\theta(0)] + \DD^x_w[0, \wtilde x_\theta(0)]\| \leq \mathcal{O}(\frac{\sigma}{\tau}\sin \alpha)
.\end{equation}
\end{Lemma}
\begin{proof}
First we bound $\Delta$ from Eq. \eqref{eq:Delta} at $x=0,\wtilde x_\theta = x_\theta(0)$ using Lemma \ref{lem:bounding_stuff}, and assuming $\alpha$ is smaller than some constant.
\begin{align}
\Delta(0,\wtilde x_\theta (0)) &= \sigma^2 + (f(0)^T \theta - f(\wtilde{x}_\theta(0))^T \theta)^2 - \|0- \wtilde{x}_\theta(0)\|^2 - \|f(0)- f(\wtilde{x}_\theta(0)) \|^2 
\notag\\
&\geq \sigma^2 -\|\wtilde{x}_\theta(0)\|^2 - 2\|f(0)- f(\wtilde{x}_\theta(0)) \|^2
\notag\\
&\geq \sigma^2 -2\sigma^2 \sin^2 \alpha - 4\sigma^2 \sin^2 \alpha 
\notag\\
&\geq \sigma^2(1-6\sin^2 \alpha )
\notag\\
&\geq \frac{1}{2} \sigma^2
.\end{align}
Similarly, 
\begin{align}
\Delta(0,\wtilde x_\theta (0)) &= \sigma^2 + (f(0)^T \theta - f(\wtilde{x}_\theta(0))^T \theta)^2 - \|0- \wtilde{x}_\theta(0)\|^2 - \|f(0)- f(\wtilde{x}_\theta(0)) \|^2 
\notag\\
&\leq \sigma^2 +\|\wtilde{x}_\theta(0)\|^2 + 2\|f(0)- f(\wtilde{x}_\theta(0)) \|^2
\notag\\
&\leq \sigma^2 +2\sigma^2 \sin^2 \alpha + 4\sigma^2 \sin^2 \alpha 
\notag\\
&\leq \sigma^2(1+6\sin^2 \alpha )
\notag\\
&\geq 2 \sigma^2
.\end{align}
and thus we showed Eq. \eqref{eq:Delta_bound}.

Next we compute $\DD^x_w$ and $D^{\wtilde{x}}_ w$
\begin{align}\label{eq:D_xw}
    \begin{split}
     \DD^x_w =& \DD_f[x] 
    +
    \theta {(\theta^T \DD_f[x])} \\
    &~~+
    {\frac{1}{\sqrt{\Delta}}} \theta \left( 
    2 {(f(x)\cdot \theta - f(\wtilde{x}_\theta) \cdot \theta)}{\theta^T \DD_f[x]} 
    - 
    2(x- \wtilde{x}_\theta)^T 
    - 
    2(f(x)- f(\wtilde{x}_\theta))^T\DD_f[x]  \right)\\
    =&\left(I_{D-d} 
    +
    \theta \theta^T  
    +
    \frac{1}{\sqrt{\Delta}} \theta \left( 
    2 (f(x)^T- f(\wtilde{x}_\theta)^T ) (\theta \theta^T - I_{D-d} )
      \right)\right) \DD_f[x] -
    \frac{2}{\sqrt{\Delta}} \theta (x- \wtilde{x}_\theta)^T 
    \end{split}
,\end{align}

\begin{align}\label{eq:D_xwtilde}
    \begin{split}
    \DD^{\wtilde{x}_\theta}_ w = & -\DD_f[\wtilde{x}_\theta] 
    - 
    \theta {(\theta^T \DD_f[\wtilde{x}_\theta])} \\
    &~~+
    {\frac{1}{\sqrt{\Delta}}} \theta \left( 
    -2 {(f(x)\cdot \theta - f(\wtilde{x}_\theta) \cdot \theta)}{\theta^T \DD_f[\wtilde{x}_\theta]} 
    + 
    2(x- \wtilde{x}_\theta)^T 
    +
    2(f(x)- f(\wtilde{x}_\theta))^T\DD_f[\wtilde{x}_\theta]  \right)\\
    =&-\left(I_{D-d} 
    +
    \theta \theta^T  
    +
    \frac{1}{\sqrt{\Delta}} \theta \left( 
    2 (f(x)^T- f(\wtilde{x}_\theta)^T ) (\theta \theta^T - I_{D-d} )
      \right)\right) \DD_f[x] +
    \frac{2}{\sqrt{\Delta}} \theta (x- \wtilde{x}_\theta)^T 
    \end{split}
.\end{align}
From Eq. \eqref{eq:D_xw}, Lemma \ref{lem:bounding_stuff}, and Eq. \eqref{eq:Delta_bound},  we have that 
\begin{align*}
     \|\DD^x_w[0,\wtilde x_\theta(0)]\| \leq& 
    \|I_{D-d} 
    +
    \theta \theta^T  
    +
    \frac{1}{\sqrt{\Delta}} \theta \left( 
    2 (f(0)^T- f(\wtilde{x}_\theta(0))^T ) (\theta \theta^T - I_{D-d} )
      \right)\|_{op} \|\DD_f[0]\|_{op} \\
      &~~+
    \frac{2}{\sqrt{\Delta}} \|(0- \wtilde{x}_\theta(0))^T\|_2 
    \\
    \leq& \OO( \sin \alpha)
\end{align*}
Similarly, from Eq. \eqref{eq:D_xwtilde}, Lemma \ref{lem:bounding_stuff}, and Eq. \eqref{eq:Delta_bound},  we have that 
\begin{equation}
     \|\DD^x_w[0,\wtilde x_\theta(0)]\| 
    \leq \OO( \sin \alpha)
\end{equation}
Thus, we showed \eqref{eq:Dws_bound}.


Now we show \eqref{eq:Dxtildew+Dxw_bound}. From \eqref{eq:D_xw} and \eqref{eq:D_xwtilde} we have 
\begin{equation}\label{eq:D_xw-Dxwtilde}
\begin{aligned}
    \sigma \DD^x_w + \sigma \DD^{\wtilde{x}_\theta}_ w 
    &= \DD_f[x] -\DD_f[\wtilde{x}_\theta] +\theta \theta^T (\DD_f[x] - \DD_f[\wtilde{x}_\theta]) \\
    &~~+\frac{1}{\sqrt{\Delta}} \theta 
    \bigg( 2 (f(x)\cdot \theta - f(\wtilde{x}_\theta) \cdot \theta)\theta^T - 2(f(x)- f(\wtilde{x}_\theta))^T\bigg) (\DD_f[x] - \DD_f[\wtilde{x}_\theta])\\
    &= \left(I_{D-d} +\theta \theta^T +\frac{1}{\sqrt{\Delta}} \theta 
    \bigg( 2 (f(x)^T - f(\wtilde{x}_\theta)^T)\theta\theta^T - 2(f(x)- f(\wtilde{x}_\theta))^T\bigg)\right) (\DD_f[x] - \DD_f[\wtilde{x}_\theta])\\
    &=\left(I_{D-d} +\theta \theta^T +\frac{2}{\sqrt{\Delta}} \theta 
      (f(x)^T - f(\wtilde{x}_\theta)^T) (\theta\theta^T - I_{D-d})\right) (\DD_f[x] - \DD_f[\wtilde{x}_\theta])
\end{aligned}
\end{equation}
Since we are bounding for $x= 0 $, we have from Lemma \ref{lem:bounding_stuff}, that $\|f(0) - f(\wtilde{x}_\theta(0))\| \leq \OO(\sigma \sin^2 \alpha)$, for $\alpha$ smaller than some constant.
Taking the norm of \eqref{eq:D_xw-Dxwtilde} we have that 
\[
\|\sigma \DD^x_w + \sigma \DD^{\wtilde{x}_\theta}_ w \| \leq \left\|I_{D-d} +\theta \theta^T +\frac{2}{\sqrt{\Delta}} \theta 
      (f(x)^T - f(\wtilde{x}_\theta)^T) (\theta\theta^T - I_{D-d})\right\| \|\DD_f[x] - \DD_f[\wtilde{x}_\theta]\|
\]
From Eq. \eqref{eq:Delta_bound} and Lemma \ref{lem:Dfx-Dfxtilde} we have,
\[
\|\sigma \DD^x_w + \sigma \DD^{\wtilde{x}_\theta}_ w \| =  \OO(\frac{\sigma}{\tau}\sin \alpha)
\]
\end{proof}


\begin{Lemma}\label{lem:Dfx-Dfxtilde}
Let $f$ be a differentiable function from $H$, a $d$-dimensional subspace of $\RR^D$, to $\RR^{D-d}$.
Assume, $\maxangle(T_0f, H)\leq \alpha$ and that $\mathrm{rch}(\Gamma_f)$ the reach of $\Gamma_f$ (the graph of the function $f$), is bounded by $\tau$.
\[
\|\DD_f[\wtilde{x}_\theta(0)] - \DD_f[0]\|_{op} \leq 
\frac{\sigma \sin (\alpha + 3\sqrt{\frac{\sigma}{\tau}}\alpha)}{\tau}(1 + \tan^2 \alpha)(1 + \sin\alpha)  +\OO(\sigma^2 \sin^2 (\alpha)/\tau^2)
\]
or
\[
\|\DD_f[\wtilde{x}_\theta(0)] - \DD_f[0]\|_{op} \leq \OO(\frac{\sigma}{\tau}\sin \alpha)
,\]
for $\alpha$ smaller than some constant.
\end{Lemma}
\begin{proof}
From Lemma \ref{lem:alpha-beta_x_no_func} we have that 
\[
	\sin(\maxangle (T_{0}f, T_{\wtilde{x}_\theta}f)) \leq \frac{\norm{\wtilde{x}_\theta}}{\tau}(1 + \tan^2 \alpha)  +\OO(\norm{\wtilde{x}_\theta}^2/\tau^2)
\]
From Lemma \ref{lem:bounding_stuff}, we have $\|\wtilde{x}_\theta\| \leq \sigma \sin (\alpha + 3\sqrt{\frac{\sigma}{\tau}}\alpha)$, and thus
\[
	\sin(\maxangle (T_{0}f, T_{\wtilde{x}_\theta}f)) \leq \frac{\sigma \sin (\alpha + 3\sqrt{\frac{\sigma}{\tau}}\alpha)}{\tau}(1 + \tan^2 \alpha)  +\OO(\sigma^2 \sin^2 (\alpha)/\tau^2)
\]
Moreover, we have that 
\[
\sin(\maxangle (T_{0}f, H)) \leq \sin \alpha
\]
Using Lemma \ref{lem:norm_from_angle} we have 
\[
\|\DD_f[\wtilde{x}_\theta] - \DD_f[x]\| 
\leq 
\frac{\sigma \sin (\alpha + 3\sqrt{\frac{\sigma}{\tau}}\alpha)}{\tau}(1 + \tan^2 \alpha)(1 + \sin\alpha)  +\OO(\sigma^2 \sin^2 (\alpha)/\tau^2) 
\]

\end{proof}
\begin{Lemma}\label{lem:norm_from_angle}
Let $L_1, L_2$ be two linear operators from $H$ a $d$-dimensional subspace of $\RR^D$ to $\RR^{D-d}$.
Let, $\maxangle(H,(H,L_1(H))_H) \leq \alpha$, where $(H,L_1(H)_H)$ is the subspace spanned by $H$ and $L_1(H)$, the target space of $L_1$. Furthermore, let $\maxangle((H,L_1(H))_H,(H,L_2(H))_H) \leq \beta$.
Then, 
\[\|L_1 - L_2\|_{op} \leq \sin \beta(1 + \sin\alpha).\]
\end{Lemma}
\begin{proof}
For any $x\in H, \|x\| = 1$, from Lemma \ref{lem:Q_angle_between_flats}, there is $y\in H$ such that 
\[
\|(x,L_1(x)) - (y,L_2(y))\| \leq \sin \beta
.\]
Therefore, 
\[
\|x-y\|^2+\|L_1(x) - L_2(y)\|^2  = \|(x,L_1(x)) - (y,L_2(y))\|^2\leq \sin^2 \beta
,\]
and
\[
\|x-y\| \leq \sin \beta ~,~ \|L_1(x) - L_2(y)\| \leq \sin\beta
.\]
Note, that $\|L_1(y) - L_1(x)\| \leq \|L_1\|_{op}\|x-y\|$.
Since $\maxangle(H,(H,L_1(H))) \leq \alpha$ we have that $\|L_1\|_{op}\leq \sin \alpha$, and we get that 
\[
\|L_1(y) - L_1(x)\| \leq \sin \alpha \sin \beta
.\]
Furthermore,
\[
\|L_1(x) - L_2(x)\| = \|L_1(x) - L_1(y) + L_1(y) - L_2(x)\| \leq \|L_1(x) - L_1(y)\|+\|L_1(y) - L_2(x)\|
,\]
and so
\[
\|L_1(x) - L_2(x)\| \leq \sin \beta(1 + \sin\alpha)
.\]
\end{proof}

\begin{Lemma}\label{lem:bounding_stuff}
Let the conditions of Lemma \ref{lem:bound_Df-Dfwtilde} hold. Let $\wtilde x_\theta(x)$ be as defined in equation \eqref{eq:x_tilde_def} in the proof of Lemma \ref{lem:bound_Df-Dfwtilde}.
Then, for $\alpha$ smaller than some constant and $\frac{\tau}{\sigma}$ larger than some constant, we have
\begin{align} \label{eq:xtilde_bound_l}
& \|\wtilde{x}_\theta(0)\| &\leq \sigma \sin (\alpha + 3\sqrt{\frac{\sigma}{\tau}}\alpha)
\\
\label{eq:f_xtilde_bound_l}
&\|f(\wtilde{x}_\theta(0)) - f(0)\| &\leq 2\sigma \sin (\alpha + 3\sqrt{\frac{\sigma}{\tau}}\alpha) \tan \alpha 
\\ \label{eq:df_bound_l}
&\|\DD_f[0]\|_2 &\leq \sin \alpha    
\\ \label{eq:df_xtilde_bound_l}
&\|\DD_f[\wtilde{x}_\theta(0)]\|_2 &\leq \sin (\alpha + 3\sqrt{\frac{\sigma}{\tau}}\alpha)
\end{align}

\end{Lemma}
\begin{proof}
In essence, this lemma is a summary and rewriting of results from other lemmas which are meant to be used conveniently in the proof of Lemma \ref{lem:bound_Df-Dfwtilde}. 
Accordingly, \eqref{eq:xtilde_bound_l} is already achieved in Lemma \ref{lem:wtilde_x_bound_clean_D}.
Then, from Lemma \ref{lem:f_bound_circle_H_no_func_ver2_clean} we have
\[
    \|f(\wtilde{x}_\theta(0))\| \leq \norm{\wtilde{x}_\theta(0)}\tan \alpha  +\OO(\|\wtilde{x}_\theta(0)\|^2/\tau)
.\]
Thus, for $\alpha$ and $\frac{\sigma}{\tau}$ smaller than some constants, we achieve Eq. \eqref{eq:f_xtilde_bound_l}.
Next, since $\maxangle(H, T_0f) \leq \alpha$, by Lemma \ref{lem:Q_angle_between_flats} we have \eqref{eq:df_bound_l}.
Finally, denoting $ \beta(\wtilde x_\theta(0)) =  \maxangle(T_{\wtilde x_\theta(0)}f, H)$ by Lemma \ref{lem:Q_angle_between_flats} we have $\|\DD_f[\wtilde{x}_\theta(0)]\|_2 \leq \sin \beta(\wtilde{x}_\theta(0))$, and combining this with Lemma \ref{lem:alpha_beta_D} we obtain \eqref{eq:df_xtilde_bound_l}

\end{proof}

\begin{Lemma}\label{lem:wtilde_x_bound_clean_D}
Let the conditions of Lemma \ref{lem:bound_Df-Dfwtilde} hold. Let $\wtilde x_\theta(x)$ be as defined in equation \eqref{eq:x_tilde_def} in the proof of Lemma \ref{lem:bound_Df-Dfwtilde}.
Then, for any unit vector $\theta\in \RR^{D-d}$, if $\alpha, \frac{\sigma}{\tau}$ are smaller than some constants, we have
	\begin{equation*}
	\|\wtilde{x}_\theta(0)\| \leq \sigma \sin (\alpha + 3\sqrt{\frac{\sigma}{\tau}}\alpha)
	\end{equation*}
\end{Lemma}
\begin{proof}
	From Lemma \ref{lem:x_tilde_bound_beta_D}, we get
	\begin{equation}\label{eq:xtilde_alpha_bound}
	\norm{\wtilde{x}_\theta(0)} \leq \sigma \sin \beta (\wtilde x_\theta(0)) 
	\end{equation}
	where $ \beta (\wtilde x_\theta(0)) = \maxangle(T_{\wtilde x_{\theta}(0)}f, H) $.
	Then, using Lemma \ref{lem:alpha_beta_D} we get
	\[
	\alpha - 4\alpha \sqrt{\frac{\sigma}{\tau}}  \leq \beta(\wtilde x_\theta(0)) \leq \alpha + 3\alpha\sqrt{\frac{\sigma}{\tau}}
	.\]
	Thus, we obtain
	\[
	\norm{\wtilde x_\theta(0)} \leq \sigma \sin (\alpha + 3\sqrt{\frac{\sigma}{\tau}}\alpha)
	,\]
	as required.
\end{proof}

\begin{Lemma}\label{lem:x_tilde_bound_beta_D}
Let the conditions of Lemma \ref{lem:bound_Df-Dfwtilde} hold. Let $\wtilde x_\theta(x)$ be as defined in equation \eqref{eq:x_tilde_def} in the proof of Lemma \ref{lem:bound_Df-Dfwtilde}.
Let $ T_{\wtilde x_\theta(0)}f $ be the tangent to the graph of $ f $ at the point $ (\wtilde x_\theta(0), f(\wtilde x_\theta(0))) $, $ \beta(\wtilde x_\theta(0)) =  \maxangle(T_{\wtilde x_\theta(0)}f, H)$.
	Then,
	\[
	\norm{\wtilde x_\theta(0) } \leq \sigma \sin \beta(\wtilde x_\theta(0))
	\]
\end{Lemma}
\begin{proof}
From \eqref{eq:x_tilde_def} of the proof of Lemma \ref{lem:bound_Df-Dfwtilde} (or more conveniently \eqref{eq:c_1} from the proof of Lemma \ref{lem:bound_nablag}) we have that 
\[
\|x-\wtilde{x}_\theta\| = \sigma \|Proj_H(\vec N_\theta(x,\wtilde{x}_\theta,\theta))\|,
\]
Using Lemma \ref{lem:angle_space_perp_space}, since $\vec N_\theta \in T_{\wtilde x_\theta(0)}f^\perp$ we have that 
\[
\|Proj_H(\vec N_\theta(x,\wtilde{x}_\theta,\theta))\| \leq \cos (\frac{\pi}{2}  - \beta(\wtilde x_\theta(0))) = \sin(\beta(\wtilde x_\theta(0))),
\]
and thus
\[
\|\wtilde{x}_\theta(0)\| = \sigma \sin(\beta(\wtilde x_\theta(0)))
\]
\end{proof}


\begin{Lemma}\label{lem:alpha_beta_D}
Let the conditions of Lemma \ref{lem:bound_Df-Dfwtilde} hold. Let $\wtilde x_\theta(x)$ be as defined in equation \eqref{eq:x_tilde_def} in the proof of Lemma \ref{lem:bound_Df-Dfwtilde}.
Let $ x_0 \in H $ be such that $ \norm{x_0} \leq \norm{\wtilde x_\theta(0)} $.
	Denote $ \beta(x) = \maxangle(T_{x}f, H) $ and let $ \alpha = \beta(0) $.
	Then,
	\[
	\alpha - 4\alpha\frac{\sigma}{\tau}  \leq \beta(x_0) \leq \alpha + 3\alpha\sqrt{\frac{\sigma}{\tau}}
	.\]
\end{Lemma}
\begin{proof}
	For convenience of notations we denote in this proof 
	\[\beta \defeq \beta(x_0).\]
	Using the result of Lemma \ref{lem:alpha_beta_x} we achieve
	\[
	\alpha - 2 \sqrt{ \frac{\norm{x_0}}{\tau}\lrbrackets{2 \alpha + \frac{\norm{x_0}}{\tau}} } 
	\leq \beta \leq
	\alpha + 2 \sqrt{ \frac{\norm{x_0}}{\tau}\lrbrackets{2 \alpha + \frac{\norm{x_0} }{\tau}}}
	,\]
	and from the fact that $ \norm{x_0} \leq \norm{\wtilde x_\theta(0)} $ we get
	\[
	\alpha - 2 \sqrt{ \frac{\norm{\wtilde x_\theta(0)}}{\tau}\lrbrackets{2 \alpha + \frac{\norm{\wtilde x_\theta(0)}}{\tau}} } 
	\leq \beta \leq
	\alpha + 2 \sqrt{ \frac{\norm{\wtilde x_\theta(0)}}{\tau}\lrbrackets{2 \alpha + \frac{\norm{\wtilde x_\theta(0)}}{\tau}} }
	\]
	From Lemma \ref{lem:x_tilde_bound_beta_D} we know that
	\[
	\norm{\wtilde x_\theta(0)} \leq \sigma \sin\beta
	,\]
	and so we get
	\begin{align*}
	\lrbrackets{\beta - \alpha}^2 
	&\leq
	4\frac{\sigma}{\tau}sin\beta\lrbrackets{2\alpha   + \frac{\sigma}{\tau}sin\beta} \\
	\frac{\beta^2-2\alpha\beta + \alpha^2}{4} 
	&\leq 
	2 \frac{\sigma}{\tau}\alpha\beta + \frac{\sigma^2}{\tau^2}\beta^2
	,\end{align*}
	that can be written as the parabola 
	\begin{align*}
	\lrbrackets{\frac{1}{4} - \frac{\sigma^2}{\tau^2}}\beta^2 - \lrbrackets{\frac{1}{2} + 2 \frac{\sigma}{\tau} }\alpha \beta + \frac{1}{4}\alpha^2 & \leq
	0 \\
	\lrbrackets{1 - 4\frac{\sigma^2}{\tau^2}}\beta^2 - \lrbrackets{2 + 8 \frac{\sigma}{\tau} }\alpha \beta + \alpha^2 & \leq
	0 
	.\end{align*}
	The left hand side of this expression is a parabola with respect to $\beta$.
	Note, that for $\frac{\sigma}{\tau} = 0$ the roots are $\beta = \alpha$.
	Solving this parabola we get the roots
	\begin{align*}
	\beta_{1,2} 
	&=
	\frac{\lrbrackets{2 +8\frac{\sigma}{\tau}}\alpha \pm \sqrt{\lrbrackets{2 + 8 \frac{\sigma}{\tau} }^2\alpha^2 - 4 \lrbrackets{1 - 4 \frac{\sigma^2}{\tau^2} }\alpha^2}}{2 - 8\frac{\sigma^2}{\tau^2}} \\
	&=
	\frac{\lrbrackets{1 +4\frac{\sigma}{\tau}}\alpha \pm \lrbrackets{1 + 4 \frac{\sigma}{\tau} }\alpha\sqrt{1 - \frac{\lrbrackets{1 - 4 \frac{\sigma^2}{\tau^2} }}{\lrbrackets{1 + 4 \frac{\sigma}{\tau} }^2}}}{1 - 4\frac{\sigma^2}{\tau^2}} \\
	&=
	\frac{\lrbrackets{1 +4\frac{\sigma}{\tau}}\alpha \pm \lrbrackets{1 + 4 \frac{\sigma}{\tau} }\alpha\sqrt{1 - \frac{\lrbrackets{1 - 4 \frac{\sigma^2}{\tau^2} }}{\lrbrackets{1 + 4 \frac{\sigma}{\tau} }^2}}}{1 - 4\frac{\sigma^2}{\tau^2}}
	.\end{align*}
	Therefore, from Remark \ref{rem:taylor_sqrt1-x2_clean} the inequality holds for
	\begin{align*}
	\beta 
	&\geq 
	\frac{\lrbrackets{1 +4\frac{\sigma}{\tau}}\alpha - \lrbrackets{1 + 4 \frac{\sigma}{\tau} }\alpha\sqrt{1 - \frac{\lrbrackets{1 - 4 \frac{\sigma^2}{\tau^2} }}{\lrbrackets{1 + 4 \frac{\sigma}{\tau} }^2}}}{1 - 4\frac{\sigma^2}{\tau^2}}\\
	&\geq
	\frac{\lrbrackets{1 +4\frac{\sigma}{\tau}}\alpha - \lrbrackets{1 + 4 \frac{\sigma}{\tau} }\alpha\lrbrackets{1 - \frac{\lrbrackets{1 - 4 \frac{\sigma^2}{\tau^2} }}{\lrbrackets{1 + 4 \frac{\sigma}{\tau} }^2}}}{1 - 4\frac{\sigma^2}{\tau^2}}\\
	&=
	\alpha\frac{1}{{1 + 4 \frac{\sigma}{\tau} }} \\
	&\geq 
	\alpha - 4\alpha\frac{\sigma}{\tau}
	,\end{align*}
	where on the other hand
	\begin{align*}
	\beta 
	&\leq
	\frac{\lrbrackets{1 +4\frac{\sigma}{\tau}}\alpha + \lrbrackets{1 + 4 \frac{\sigma}{\tau} }\alpha\sqrt{1 - \frac{\lrbrackets{1 - 4 \frac{\sigma^2}{\tau^2} }}{\lrbrackets{1 + 4 \frac{\sigma}{\tau} }^2}}}{1 - 4\frac{\sigma^2}{\tau^2}} \\
	&=
	\frac{\lrbrackets{1 +4\frac{\sigma}{\tau}}\alpha + \lrbrackets{1 + 4 \frac{\sigma}{\tau} }\alpha\sqrt{1 - \frac{\lrbrackets{1 - 2 \frac{\sigma}{\tau} }\lrbrackets{1 + 2 \frac{\sigma}{\tau} }}{\lrbrackets{1 + 4 \frac{\sigma}{\tau} }^2}}}{1 - 4\frac{\sigma^2}{\tau^2}} \\
	&\leq
	\frac{\lrbrackets{1 +4\frac{\sigma}{\tau}}\alpha + \lrbrackets{1 + 4 \frac{\sigma}{\tau} }\alpha\sqrt{1 - \frac{\lrbrackets{1 - 2 \frac{\sigma}{\tau} }\lrbrackets{1 + 2 \frac{\sigma}{\tau} }}{\lrbrackets{1 + 2 \frac{\sigma}{\tau} }^2}}}{1 - 4\frac{\sigma^2}{\tau^2}} \\
	&=
	\frac{\lrbrackets{1 +4\frac{\sigma}{\tau}}\alpha + \lrbrackets{1 + 4 \frac{\sigma}{\tau} }\alpha\sqrt{1 - \frac{{1 - 2 \frac{\sigma}{\tau} }}{{1 + 2 \frac{\sigma}{\tau} }}}}{1 - 4\frac{\sigma^2}{\tau^2}} 
	,\end{align*}
	since we know that $ \frac{1}{1 + x} \geq 1 - x $ we get
	\begin{align*}
	\beta 
	&\leq
	\frac{\lrbrackets{1 +4\frac{\sigma}{\tau}}\alpha + \lrbrackets{1 + 4 \frac{\sigma}{\tau} }\alpha\sqrt{1 - \lrbrackets{1 - 2 \frac{\sigma}{\tau}}^2}}{1 - 4\frac{\sigma^2}{\tau^2}} \\
	&\leq
	\frac{\lrbrackets{1 +4\frac{\sigma}{\tau}}\alpha + \lrbrackets{1 + 4 \frac{\sigma}{\tau} }\alpha\sqrt{4\frac{\sigma}{\tau} - 4\frac{\sigma^2}{\tau^2}}}{1 - 4\frac{\sigma^2}{\tau^2}}\\
	&\leq
	\frac{\lrbrackets{1 +4\frac{\sigma}{\tau}}\alpha + 2\lrbrackets{1 + 4 \frac{\sigma}{\tau} }\alpha\sqrt{\frac{\sigma}{\tau}}}{1 - 4\frac{\sigma^2}{\tau^2}}
	,\end{align*}
	and by using the bound $ \frac{1}{1 - x^2} \leq 1 + 2 x^2 $ for $ 0 \leq x \leq 0.5 $	
	\begin{align*}
	\beta
	&\leq
	\lrbrackets{\lrbrackets{1 +4\frac{\sigma}{\tau}}\alpha + 2\lrbrackets{1 + 4 \frac{\sigma}{\tau} }\alpha\sqrt{\frac{\sigma}{\tau}}}\lrbrackets{1 + 8\frac{\sigma^2}{\tau^2}} \\
	&=
	\alpha + 2\alpha\sqrt{\frac{\sigma}{\tau}} + 4\alpha\frac{\sigma}{\tau} + 8 \alpha\frac{\sigma^{1.5}}{\tau^{1.5}} + 8\alpha\frac{\sigma^2}{\tau^2} + 16\alpha\frac{\sigma^{2.5}}{\tau^{2.5}} + 32\alpha\frac{\sigma^{3}}{\tau^{3}} + 32\alpha\frac{\sigma^{3.5}}{\tau^{3.5}}
	.\end{align*}
	Then, since $ \frac{\sigma}{\tau}\leq \frac{1}{36} $ is sufficiently small we obtain 
	\[
	\beta \leq \alpha + 3\alpha\sqrt{\frac{\sigma}{\tau}}
	,\]
	and 
	\[
	\alpha - 4\alpha\frac{\sigma}{\tau}  \leq \beta \leq \alpha + 3\alpha\sqrt{\frac{\sigma}{\tau}}
	\]
\end{proof}% --- end Section_ProofStep2_BoundTf_Tf_tilde.tex ---

\subsubsection{Bounding the finite sample error}\label{sec:FiniteSample_err}
% --- begin Section_ProofStep2_BoundTf_tilde_TF_hat.tex ---
Back to Theorem \ref{thm:Step2} proof road-map see Figure \ref{tikz:thm33_lemmas}.

\noindent In this section we show that  $ \maxangle(T_0 \wtilde f_\ell, H_{\ell + 1}) $ the angle between the tangent of $\wtilde{f}_\ell(0)$ and the tangent estimated using $n$ samples decays to zero as $n\to \infty$.
Namely, the main result of this subsection is pronounced in the lemma below.

\begin{Lemma}\label{lem:AngleImprovement}
	Let $(q_\ell, H_\ell)$ be defined in Algorithm \ref{alg:step2_clean} and $ \pi^*_{q_\ell, H_\ell}(x) $ be defined in \eqref{eq:argmin2_clean} and let $ H_{\ell + 1} = T_0 \pi^*_{q_\ell, H_\ell} $ the tangent to the graph of $ \pi^*_{q_\ell, H_\ell} $ at $ \pi^*_{q_\ell, H_\ell}(0) $.
	Then, for all $ \delta>0 $ there is $ N_\delta $ such that for all $ n > N_\delta $ we have with probability $ 1-\delta $
	\[
	\maxangle(T_0\wtilde f_\ell, H_{\ell + 1}) \leq 2\sqrt{d}\frac{C_0 \ln(1/\delta)}{n ^{r_1}}
	,\]
	where ${r_1} = \frac{k-1}{2k + d}$ and $C_0$ is a constant.
\end{Lemma}
\begin{proof}
We first note that it is sufficient to bound the error of estimating the image of $\DD_{\wtilde{f}_\ell}[0]$, the differential of $ \wtilde f_\ell(x) $ at 0, by the image of $ \DD_{\pi^*_{q_\ell, H_\ell}}[0] $, the differential of the local polynomial least-squares regression $\pi^*_{q_\ell, H_\ell}$.
Explicitly, if 
\[
\|{\DD_{\pi^*_{q_\ell, H_\ell}}[0] - \DD_{ \wtilde f_\ell}[0]}\|_{op} \leq \sqrt{d}\frac{C_0 \ln(1/\delta)}{n^{r_1}}
,\]
then, by using Lemma \ref{lem:Diff_Tf_bound} we get that
\[
\sin(\maxangle(T_0\wtilde f_\ell, T_0\pi^*_{q_\ell, H_\ell})) \leq \sqrt{d}\frac{C_0 \ln(1/\delta)}{n^{r_1}}
,\]
which for sufficiently large $ n $ yields
\[
\maxangle(T_0\wtilde f_\ell, T_0\pi^*_{q_\ell, H_\ell}) \leq 2\sqrt{d}\frac{C_0 \ln(1/\delta)}{n^{r_1}}
,\]
as required.

Therefore, it is sufficient to show that for any $\delta$ there is $N_\delta$ such that for all $n>N_\delta$ we have 
\[
\|{\DD_{\pi^*_{q_\ell, H_\ell}}[0] - \DD_{ \wtilde f_\ell}[0]}\|_{op} \leq \sqrt{d}\frac{C_0 \ln(1/\delta)}{n^{r_1}}
,\]
with probability of at least $1 - \delta$.
Let us reiterate the minimization problem by which we derive the approximant.
Namely, given a sample $ \{r_i\}_{i=1}^{n} $ drawn i.i.d from $ \text{Unif}(\MM_\sigma) $, and a coordinate system $ (q, H)\in \RR^D\times Gr(d,D) $ we look for a polynomial $ \pi^*_{q,H} $ minimizing
\begin{equation}\label{eq:Step2_sample}
J_2(\pi ~|~ q, H) = \frac{1}{N_{q,H}}\sum_{r_i \in U_{\textrm{ROI}}^{n}} \norm{r_i - \pi(x_i)}^2
,\end{equation}
where $x_i$ are the projections of $r_i - q$ onto $H$, and $U_{\textrm{ROI}}^{n}(q,H)$ is defined through a bandwidth $\epsilon_{n}$ as
\begin{equation}\label{eq:ROI_ell_sample}
U_{\textrm{ROI}}^{n}(q,H) = {\{r_i\in U_\textrm{ROI}~|~\dist(x_i, q)<\epsilon_{n}\}}
,\end{equation}
and $N_{q,H}$ denotes the number of samples in $U_{\textrm{ROI}}^{n}(q, H)$.
Explicitly,
\begin{equation}\label{eq:argmin2_sample}
\pi^*_{q_\ell, H_\ell} = \argmin_{\pi\in \Pi_{k-1}^{d\mapsto D}} J_2(\pi ~|~ q_\ell, H_\ell)
.\end{equation}
We demand that the bandwidth $\epsilon_{n}\to 0$ as ${n}\to\infty$ such that 
\begin{equation}\label{eq:bandwidthDecay_sample}
0<\lim_{{n}\rightarrow\infty}N^{1/(2k+1)}\cdot \epsilon_{n} < \infty
.\end{equation}
And, the approximation is defined through $\DD_{\pi^*_{q_\ell, H_\ell}}[0]$


From Lemma \ref{lem:noise_cov_bound} we can apply Theorem 3.2 from \cite{aizenbud2021VectorEstimation} that gives convergence rates for local polynomial regression of vector valued functions in our case. Thus, we have that for every direction in the basis $\{x^j\}_{j=1}^d\subset H_\ell$ and every $\delta$ there exists  $N_{\delta}$ such that for all $n > N_{\delta}$ we have
\[
\Pr(\|{\partial_{x^j} \pi^*_{q_\ell, H_\ell}(0) - \partial_{x^j} \wtilde f_\ell(0)}\| > \frac{C_0 \ln(1/\delta)^{r_1}}{n^{r_1}}) < \delta
,\]
where $r_1 = \frac{k-1}{2k + d}$ and $C_0$ is a constant. Notice that $r_1 \leq 1/2$, and thus 
\[
\Pr(\|{\partial_{x^j} \pi^*_{q_\ell, H_\ell}(0) - \partial_{x^j} \wtilde f_\ell(0)}\| > \frac{C_0 \ln(1/\delta)}{n^{r_1}}) < \delta
.\]
Thus, taking into account all $d$ directions of the basis to $H_\ell$ we get that there are $C$ and $N_\delta$ such that for all $n > N_\delta$  
\[
\Pr(\|{\partial_{x^j} \pi^*_{q_\ell, H_\ell}(0) - \partial_{x^j} \wtilde f_\ell(0)}\| > \frac{C_0 \ln(1/\delta)}{n^{r_1}}\textrm{ for any } 1 \leq j \leq d) < d\delta
,\]
and thus
\[
\Pr(\|{\DD_{\pi^*_{q_\ell, H_\ell}}[0] - \DD_{\wtilde f_\ell}[0]}\|_{op} > \sqrt{d}\frac{C_0 \ln(1/\delta)}{n^r}) < d\delta
,\]
as required.
\end{proof}

In order to use convergence rate results of local polynomial regression for vector valued functions as described in Theorems 3.1 and 3.2 of \cite{aizenbud2021VectorEstimation} in our case, we need to show that the noise distribution $\eta_\ell$ defined in \eqref{eq:ftilde_def} is such that $\|\cov(\eta_\ell)\|\leq \sqrt{c/D}$.
\begin{Lemma}\label{lem:noise_cov_bound}
 Let $H_\ell\in Gr(d, D)$, and let $f_\ell:H_\ell\simeq\RR^d \rightarrow \RR^{D-d}$, defined as in \eqref{eq:fl_def}. Let  $\eta_\ell$ defined in \eqref{eq:ftilde_def}. Denote 
$\alpha_\ell = \maxangle (T_0f_{\ell}, H_{\ell})$ and assume $\alpha_\ell<1/D^{1/4}$. Then, 
\[
\|\cov(\eta_\ell)\|_{op} \leq \sqrt{\frac{c\sigma}{D-d}}
\]
    
\end{Lemma}
\begin{proof}
For ease of notation, denote $\wtilde{D} = D-d$.
Since we are interested in bounding 
\begin{equation}\label{eq:cov_norm}
\|\cov(\eta_\ell) \|_{op} = \max_{\vec{x}\in \mathbb{S}_{\wtilde{D} - 1} } \vec{x}^T \cov(\eta_\ell)\vec{x}
\end{equation}
we note that 
\[
\vec{x}^T\cov(\eta_\ell)\vec{x} = \mbox{Var}(\eta_\ell\cdot \vec x)
.\]
Thus, rewriting \eqref{eq:cov_norm} we have
\begin{equation}\label{eq:cov_norm2}
    \|\cov(\eta_\ell)\|_{op}= \max_{\vec{x}\in \mathbb{S}_{\wtilde{D} - 1}} \vec{x}^T\cov(\eta_\ell)\vec{x} = \max_{\vec{x}\in \mathbb{S}_{\wtilde{D} - 1}}\mbox{Var}(\eta_\ell\cdot \vec{x}) \leq \max_{\vec{x}\in \mathbb{S}_{\wtilde{D} - 1}}\EE((\eta_\ell\cdot \vec{x})^2)
\end{equation}
Thus, our goal is to bound, for any $\vec{z}\in \mathbb{S}_{\wtilde{D} - 1}$  the expression $\EE((\eta_\ell\cdot \vec{z})^2)$. From the definition of $g$ and $\Omega$ in \eqref{eq:g_def} and \eqref{eq:Omega_def}, we have that
\begin{align}\label{eq:E(Z_z)}
\EE((\eta_\ell\cdot \vec{z})^2) &= \frac{\int\limits_{y\in\Omega(x)} (y \cdot \vec{z})^2 dy}{\int\limits_{y\in\Omega(x)} dy} = 
\frac{\int\limits_{\mathbb{S}_{\wtilde D - 1}}\int\limits_0^{g(x,\theta)} (\theta \cdot \vec{z})^2 r^2 r^{\wtilde D-1} drd\theta}{\int\limits_{\mathbb{S}_{\wtilde D - 1}}\int\limits_0^{g(x,\theta)} r^{\wtilde D-1} drd\theta} \notag\\
& = \frac{\wtilde D\int\limits_{\mathbb{S}_{\wtilde D - 1}}(\theta \cdot \vec{z})^2 g(x,\theta)^{\wtilde D+2} d\theta}{(\wtilde D+2)\int\limits_{\mathbb{S}_{\wtilde D - 1}}g(x,\theta)^{\wtilde D} drd\theta}     
,\end{align}
where $dr$ is the measure over the radial component, $r^{\wtilde D-1}$ is the Jacobian introduced by the change of variables and $d\theta$ is the measure over the $(\wtilde D-1)$-dimensional sphere.

Following the rationale of the proof of Lemma \ref{lem:bound_Df-Dfwtilde}, we split $\mathbb{S}_{\wtilde D-1}$ into $\Omega_1$ and $\Omega_2$ of \eqref{eq:Omegas_def}.
That is,
\begin{equation*}
    \begin{aligned}
    \Omega_1 &= \{\theta~|~0\leq \vec{z}^T\theta \leq \xi\}\\
    \Omega_2 &= \{\theta~|~\vec{z}^T\theta > \xi\}
    \end{aligned}
,\end{equation*}
for some $\xi$ to be chosen later.
Thus, denoting $z = \theta^T \vec{z}$ we rewrite \eqref{eq:E(Z_z)} as` 
\begin{align*}
\EE((\eta_\ell\cdot \vec{z})^2)   &=\frac{\wtilde D\left(\int\limits_{\Omega_1}z^2 g(x,\theta)^{\wtilde D+2} d\theta + \int\limits_{\Omega_2}z^2 g(x,\theta)^{\wtilde D+2} d\theta\right)}{(\wtilde D+2)\int\limits_{\mathbb{S}_{\wtilde D-1}}g(x,\theta)^{\wtilde D} d\theta} 
\\
&\leq  \frac{\wtilde D\left(\xi^2 \int\limits_{\Omega_1} g(x,\theta)^{\wtilde D+2} d\theta + \int\limits_{\Omega_2} g(x,\theta)^{\wtilde D+2} d\theta\right)}{(\wtilde D+2)\int\limits_{\mathbb{S}_{\wtilde D-1}}g(x,\theta)^{\wtilde D} d\theta} 
\end{align*}

Since the conditions of Lemma \ref{lem:g_bond_weak} are met, we have $\sigma \leq g(0,\theta)\leq \sigma + 4\sigma \alpha_\ell ^2$, and thus
\begin{align*}
\EE((\eta_\ell\cdot \vec{z})^2)
&\leq  \frac{\wtilde D\left(\xi^2 \int\limits_{\Omega_1} g(x,\theta)^{\wtilde D+2} d\theta + \int\limits_{\Omega_2} g(x,\theta)^{\wtilde D+2} d\theta\right)}{(\wtilde D+2)\int\limits_{\mathbb{S}_{\wtilde D-1}}g(x,\theta)^{\wtilde D} d\theta}
\\
&\leq  \frac{\wtilde D\sigma^2(1+4\alpha_\ell ^2)^2\left(\xi^2  \int\limits_{\Omega_1} g(x,\theta)^{\wtilde D} d\theta + \int\limits_{\Omega_2} g(x,\theta)^{\wtilde D} d\theta\right)}{(\wtilde D+2)\int\limits_{\mathbb{S}_{\wtilde D-1}}g(x,\theta)^{\wtilde D} d\theta}
\\
&\leq  \frac{\wtilde D\sigma^2(1+4\alpha_\ell ^2)^2}{\wtilde D+2}\left(\xi^2  +\frac{ \int\limits_{\Omega_2} g(x,\theta)^{\wtilde D} d\theta}{\int\limits_{\mathbb{S}_{\wtilde D-1}}g(x,\theta)^{\wtilde D} d\theta} \right)
\\
&\leq  \frac{\wtilde D\sigma^2(1+4\alpha_\ell ^2)^2}{\wtilde D+2}\left(\xi^2  +\frac{\sigma^{\wtilde D}(1+4\alpha_\ell ^2)^{\wtilde D} \int\limits_{\Omega_2} d\theta}{\sigma^{\wtilde D}\int\limits_{\mathbb{S}_{\wtilde D-1}} d\theta} \right)
\\
&\leq  \frac{\wtilde D\sigma^2(1+4\alpha_\ell ^2)^2}{\wtilde D+2}\left(\xi^2 +\frac{(1+4\alpha_\ell ^2)^{\wtilde D}}{\xi\sqrt{\wtilde D-1}} e^{-(\wtilde D-1)\xi^2/2}\right)
,\end{align*}
where the last inequality comes from Eq. \eqref{eq:Omega3_to_ball_ratio}.
Since $\alpha_\ell <1/D^{1/4}$ we have that $(1+4\alpha_\ell ^2)^{\wtilde D}$ is bounded by some constant $c$. Choosing $\xi = 2\sqrt{\frac{\log(\wtilde D-1)}{\wtilde D-1}}$ we have 
\begin{align*}
\EE((\eta_\ell\cdot \vec{z})^2)
&\leq  \frac{\wtilde D\sigma^2(1+4\alpha_\ell ^2)^2 }{\wtilde D+2}\left(4\frac{\log(\wtilde D-1)}{\wtilde D-1} +\frac{c}{2\sqrt{\log(\wtilde D-1)}} (\wtilde D-1)^{-2}\right)
\\
&\leq  c_1\sigma^2\frac{\log(\wtilde D-1)}{\wtilde D-1}\leq  \sqrt{\frac{c\sigma}{\wtilde D}}
\end{align*}
for some constants $c,c_1$. Combining with Eq. 
\eqref{eq:cov_norm2}, we conclude the proof.

\end{proof}


% For completeness, we present here Theorem 2.1 of \cite{aizenbud2021VectorEstimation}, which is an extension of the results in \cite{stone1980optimal}.
% \begin{Theorem}[From \cite{aizenbud2021VectorEstimation}]\label{thm:vector_LPR}
% Let $ (X,Y) $ be a pair of random variables such that $ X $ in $ \RR^d $ and $ Y $ in $ \RR^{\mathfrak{D}} $.
% 	Assume that the distribution of $ X $ is absolutely continuous and that its density is bounded away from zero and from infinity on $ U $, an open neighborhood of the origin of $ \RR^d $.
% 	Denote the regression function $ f:\RR^d\to\RR^{\mathfrak{D}} $ be 
% 	\[
% 	f = \EE(Y ~|~ X=x)
% 	\]
% 	Let the conditional distribution of $ Y $ given $ X=x $ be uniformly distributed on $ I(f(x)) $ a ball of finite radius centered around $ f(x) $.
% 	Let $ f\in \C^k $, $ \alpha = (\alpha_1, \ldots, \alpha_d) $ be an integer valued multi-index, $ \partial^\alpha = \partial^{\alpha_1}_{x^1}\cdots \partial^{\alpha_d}_{x^d} $.
% 	Given a sample $ \{(x_i, y_i)\}_{i=1}^{n} $ drawn i.i.d from $ (X, Y) $, we define the local polynomial regressor $ \pi^*  $ as
% 	\[
% 	\pi^*(x) = \argmin_{\pi\in \Pi_{k-1}^{d\mapsto \mathfrak{D}}} \sum_{\norm{x_i}\leq \epsilon_{n} }\norm{y_i - \pi(x_i)}^2
% 	,\]
% 	where $ \epsilon_n $ is such that $ 0 < \lim\limits_{n\to\infty}n^{\frac{1}{2k+d}} \epsilon_n < \infty$.
% 	Let $ L = \sum_{\abs{\alpha}\leq m}c_\alpha \partial^\alpha $, and define the operator
% 	\[
% 	T(f) = L[f](0)\in \RR^\mathfrak{D}
% 	,\]
% 	as well as its estimate by
% 	\[
% 	\widehat T_{n} = L[\pi^*](0)\in \RR^\mathfrak{D}
% 	.\]
% 	Then, for any $\varepsilon<0.5$ there are $C_\eps$ and $ N_\eps $ such that for any $ n > N_\eps $ we have
% 	\[
% 	\Pr(\|{\widehat T_{n} - T(f)}\| > \frac{C_\eps}{n^{r_m}}) \leq \varepsilon
% 	,\]
% 	for $ r_m =\frac{k-m}{2k+d} $.
	
% 	Moreover, there are some constants $C_0$ and $N_0$ such that 
% 	\[
% 	N_\eps = N_0 \ln \frac{1}{\varepsilon},
% 	\]
% 	and 
% 	\[
% 	C_\eps = C_0 \left(\ln \frac{1}{\varepsilon}\right)^{r_m}.
% 	\]
	
% \end{Theorem}


\begin{comment}
Theorem \ref{thm:stone} deals with the case of $ f:\RR^d\to \RR $, which in our case translates to $ D=d+1 $ (i.e., $ \wtilde f_\ell:\RR_d\to\RR $).
If this was indeed the case ($ D=d+1 $) then an immediate result would have been that for every direction in the basis $ \{x^j\}_{j=1}^d \subset H_\ell $, and every $ \varepsilon $ there exists $ C $ and $ N $ large enough such that for all $ n>N $ we have
\[
\Pr(|{\partial_{x^j} \pi^*_{q_\ell, H_\ell}(0) - \partial_{x^j} \wtilde f_\ell}(0)| > \frac{C}{n^r}) < \varepsilon
,\]
where $r = \frac{k-1}{2k+d}$.
In the case of $ D = d+1 $ we have $ \DD_{\wtilde f_\ell} = (\nabla\wtilde f_\ell)^T $ we get that there is $ C $ and $ N $ such that
\[
\Pr(|{\partial_{x^j} \pi^*_{q_\ell, H_\ell}(0) - \partial_{x^j} \wtilde f_\ell}(0)| > \frac{C}{n^r} \text{ for any } 1\leq j \leq d) < d\varepsilon
,\]
and thus
\[
\Pr(\|{\nabla \pi^*_{q_\ell, H_\ell}(0) - \nabla \wtilde f_\ell(0)}\| > \sqrt{d}\frac{C}{n^r}) < d\varepsilon
.\]
Removing the assumption that $ D = d+1 $ we could apply similar consideration for each coordinate of the $ D-d $ coordinates in the image of $\wtilde  f_\ell $ and achieve that
\[
\Pr(\|{\DD_{\pi^*_{q_\ell, H_\ell}}[0] - \DD_{ \wtilde f_\ell}[0]}\|_{op} > \sqrt{d}\frac{C}{n^r}) < (D-d)d\varepsilon
.\]
However, this requires that the number of samples we need depends on the ambient dimension $ D $ rather than just the intrinsic one $ d $, which seems excessive.
Thus, we know with failure probability of $ \eps(n, D, d) $ that 
\[
\|{\DD_{\pi^*_{q_\ell, H_\ell}}[0] - \DD_{ \wtilde f_\ell}[0]}\|_{op} \leq \sqrt{d}\frac{C}{n^r}
,\]
and using Lemma \ref{lem:Diff_Tf_bound} we get that
\[
\sin(\maxangle(T_0\wtilde f, T_0\pi^*_{q_\ell, H_\ell})) \leq \sqrt{d}\frac{C}{n^r}
,\]
which for sufficiently large $ n $ yields
\[
\maxangle(T_0\wtilde f, T_0\pi^*_{q_\ell, H_\ell}) \leq 2\sqrt{d}\frac{C}{n^r}
,\]
as required.

We wish to argue below that since the sampling assumption of Stone's theorem involves the domain and not the target of the function, the theorem could be generalized to arbitrary image dimension without requiring more samples.
However, since this is not in the focus of the current article (which is evidently wide enough), we do not handle this problem in a completely formal manner.
Thus, we will pose this as a conjecture that will be followed by a heuristic justification. 




Applying Conjecture \ref{conj:approx} for estimating the differential of $ \wtilde f_\ell $ we get in a similar manner to the reasoning above that
\[
\Pr(\|{\DD_{\pi^*_{q_\ell, H_\ell}}[0] - \DD_{ \wtilde f_\ell}[0]}\|_{op} > \sqrt{d}\frac{C}{n^r}) < d\varepsilon
.\]
Thus, we know with failure probability of $ \eps(n, d) $ that 
\[
\|{\DD_{\pi^*_{q_\ell, H_\ell}}[0] - \DD_{ \wtilde f_\ell}[0]}\|_{op} \leq \sqrt{d}\frac{C}{n^r}
,\]
and using Lemma \ref{lem:Diff_Tf_bound} we get that
\[
\sin(\maxangle(T_0\wtilde f, T_0\pi^*_{q_\ell, H_\ell})) \leq \sqrt{d}\frac{C}{n^r}
,\]
which for sufficiently large $ n $ yields
\[
\maxangle(T_0\wtilde f, T_0\pi^*_{q_\ell, H_\ell}) \leq 2\sqrt{d}\frac{C}{n^r}
,\]
as required.
\end{comment}
% --- end Section_ProofStep2_BoundTf_tilde_TF_hat.tex ---

\subsubsection{Bounding the distance of $q_\ell$ from $f_\ell(0)$}\label{sec:f0_est}
% --- begin Section_ProofStep2_f0_est.tex ---


Back to Theorem \ref{thm:Step2} proof road-map see Figure \ref{tikz:thm33_lemmas}.

\begin{Lemma}\label{lem:dist_to_f_l0_weak}
For $f_\ell$ defined in \eqref{eq:fl_def}. Denote 
$\alpha_\ell = \maxangle (T_0f_{\ell}, H_{\ell})$ and assume $\alpha_\ell<1/D^{1/4}$. Then,
for any $\delta$ there is $N$ such that for any number of samples $n > N$
\[
\|f_\ell(0)\| \leq \tau \alpha_\ell /48
\]
with probability of at least $1-\delta$.
\end{Lemma}
\begin{proof}
In this proof we will assume that $\alpha_\ell  \geq\frac{1}{D}$. The case when $\alpha_\ell <\frac{1}{D}$ will be treated at the end of this proof.
Using the triangle inequality, we have 
    \begin{equation}\label{eq:f_ell0_triangle}
        \|f_\ell(0)\| \leq \| \wtilde{f}_\ell(0) \| + \| \wtilde{f}_\ell(0)  - f_\ell(0)\|.
    \end{equation}
    and from Lemma \ref{lem:dist_q_l_and f_tilde0} we have that 
    \begin{equation}\label{eq:f_tilde0_bound_stone}
            \| \wtilde{f}_\ell(0)\| \leq \frac{C_2\ln\left(\frac{1}{\delta}\right)}{n^{r_0}}
    \end{equation}
    with probability of at least $1-\delta$.
    
Similar to \eqref{eq:f_tilde-f_integrals} of the proof of Lemma \ref{lem:bound_Df-Dfwtilde} we can write 
\begin{equation*}
\wtilde{f}_\ell(x) - f_\ell(x)=
\frac{\wtilde D\int\limits_{\mathbb{S}_{\wtilde D-1}}\theta g(x,\theta)^{\wtilde D+1} d\theta}{(\wtilde D+1)\int\limits_{\mathbb{S}_{\wtilde D-1}}g(x,\theta)^{\wtilde D} drd\theta}    ,\end{equation*}
where $\wtilde D = D - d$, and $\mathbb{S}_{\wtilde D - 1}$ is the $({\wtilde D - 2})$-dimensional unit sphere in $\RR^{\wtilde D - 1}$.
Following the rationale of the proof of Lemma \ref{lem:bound_Df-Dfwtilde}, we split $\mathbb{S}_{\wtilde D-1}$ into $\Omega_1$ and $\Omega_2$ of \eqref{eq:Omegas_def}.
That is,
\begin{equation*}
    \begin{aligned}
    \Omega_1 &= \{\theta~|~0\leq \vec{z}^T\theta \leq \xi\}\\
    \Omega_2 &= \{\theta~|~\vec{z}^T\theta > \xi\}
    \end{aligned}
,\end{equation*}
for some $\xi$ to be chosen later.
Thus, we have  
\begin{align*}
z^T\cdot (\wtilde{f}_\ell(x) - f_\ell(x))&=
\frac{\wtilde D\int\limits_{\mathbb{S}_{\wtilde D-1}}z g(x,\theta)^{\wtilde D+1} d\theta}{(\wtilde D+1)\int\limits_{\mathbb{S}_{\wtilde D-1}}g(x,\theta)^{\wtilde D} d\theta} 
\\
&=\frac{\wtilde D\left(\int\limits_{\Omega_1}z g(x,\theta)^{\wtilde D+1} d\theta + \int\limits_{\Omega_2}z g(x,\theta)^{\wtilde D+1} d\theta\right)}{(\wtilde D+1)\int\limits_{\mathbb{S}_{\wtilde D-1}}g(x,\theta)^{\wtilde D} d\theta} 
\\
&\leq  \frac{\wtilde D\left(\xi \int\limits_{\Omega_1} g(x,\theta)^{\wtilde D+1} d\theta + \int\limits_{\Omega_2} g(x,\theta)^{\wtilde D+1} d\theta\right)}{(\wtilde D+1)\int\limits_{\mathbb{S}_{\wtilde D-1}}g(x,\theta)^{\wtilde D} d\theta} 
\end{align*}
Since the conditions of Lemma \ref{lem:g_bond_weak} are met, we have $\sigma \leq g(0,\theta)\leq \sigma + 4\sigma \alpha_\ell ^2$, and thus
\begin{align*}
z^T\cdot (\wtilde{f}_\ell(x) - f_\ell(x))
&\leq  \frac{\wtilde D\left(\xi \int\limits_{\Omega_1} g(x,\theta)^{\wtilde D+1} d\theta + \int\limits_{\Omega_2} g(x,\theta)^{\wtilde D+1} d\theta\right)}{(\wtilde D+1)\int\limits_{\mathbb{S}_{\wtilde D-1}}g(x,\theta)^{\wtilde D} d\theta}
\\
&\leq  \frac{\wtilde D\sigma(1+4\alpha_\ell ^2)\left(\xi  \int\limits_{\Omega_1} g(x,\theta)^{\wtilde D} d\theta + \int\limits_{\Omega_2} g(x,\theta)^{\wtilde D} d\theta\right)}{(\wtilde D+1)\int\limits_{\mathbb{S}_{\wtilde D-1}}g(x,\theta)^{\wtilde D} d\theta}
\\
&\leq  \frac{\wtilde D\sigma(1+4\alpha_\ell ^2)}{\wtilde D+1}\left(\xi  +\frac{ \int\limits_{\Omega_2} g(x,\theta)^{\wtilde D} d\theta}{\int\limits_{\mathbb{S}_{\wtilde D-1}}g(x,\theta)^{\wtilde D} d\theta} \right)
\\
&\leq  \frac{\wtilde D\sigma(1+4\alpha_\ell ^2)}{\wtilde D+1}\left(\xi  +\frac{\sigma^{\wtilde D}(1+4\alpha_\ell ^2)^{\wtilde D} \int\limits_{\Omega_2} d\theta}{\sigma^{\wtilde D}\int\limits_{\mathbb{S}_{\wtilde D-1}} d\theta} \right)
\\
&\leq  \frac{\wtilde D\sigma(1+4\alpha_\ell ^2)}{\wtilde D+1}\left(\xi  +\frac{(1+4\alpha_\ell ^2)^{\wtilde D}}{\xi\sqrt{\wtilde D-1}} e^{-(\wtilde D-1)\xi^2/2}\right)
\end{align*}
Since $\alpha_\ell <1/D^{1/4}$ we have that $(1+4\alpha_\ell ^2)^{\wtilde D}$ is bounded by some constant $c$. Choosing $\xi = 2\sqrt{\frac{\log(\wtilde D-1)}{\wtilde D-1}}$ we have 
\begin{align*}
z^T\cdot (\wtilde{f}_\ell(x) - f_\ell(x))
&\leq  \frac{\wtilde D\sigma(1+\alpha_\ell )}{\wtilde D+1}\left(2\sqrt{\frac{\log(\wtilde D-1)}{\wtilde D-1}} +\frac{c}{2\sqrt{\log(\wtilde D-1)}} (\wtilde D-1)^{-2}\right)
\\
&\leq  c_1\sigma\sqrt{\frac{\log(\wtilde D-1)}{\wtilde D-1}} 
\end{align*}
for some constant $c_1$. Recalling that $M=\frac{\tau}{\sigma}>C_\tau \sqrt{D \log D}$ we have,
\begin{equation*}
    \|\wtilde{f}_\ell(x) - f_\ell(x)\| \leq c_1\sigma\sqrt{\frac{\log(\wtilde D-1)}{\wtilde D-1}} = \sigma C_\tau \sqrt{D\log D} \frac{c_1}{C_\tau D}\leq \tau \frac{c_1}{C_\tau D}.
\end{equation*}
and thus, for $\alpha_\ell \geq 1/D$, and for large enough $\C_\tau$ we have that 
\begin{equation*}
    \|\wtilde{f}_\ell(x) - f_\ell(x)\|\leq \frac{\tau \alpha_\ell }{2\cdot 48}.
\end{equation*}



Combining this with  \eqref{eq:f_ell0_triangle} and \eqref{eq:f_tilde0_bound_stone}, we have that for any $\delta>0$, for $\alpha_\ell >\frac{1}{D}$, and for number of samples $n>N$ large enough, 
\[
\|f_\ell(0)\| \leq \tau \alpha_\ell /48,
\]
with probability of at least $1-\delta$.

For $\alpha_\ell<\frac{1}{D}$ using Lemma \ref{lem:dist_to_f_l0_srtong_small_alpha} concludes the proof.
\end{proof}




\begin{Lemma}\label{lem:dist_to_f_l0_srtong_small_alpha}
For $f_\ell$ defined in \eqref{eq:fl_def}. Denote 
$\alpha_\ell = \maxangle (T_0f_{\ell}, H_{\ell})$ and assume $\alpha_\ell<1/D$. Then,
for any $\delta$ there is $N$ such that for any number of samples $n > N$, we have 
\[
\| f_\ell(0)\| \leq D C_{1} \alpha_\ell ^2 + \frac{C_2\ln\left(\frac{1}{\delta}\right)}{n^{r_0}}.
\]
\end{Lemma}
\begin{proof}
Using the triangle inequality, we have 
\begin{equation}\label{eq:f_ell0_triangle_l2}
    \|f_\ell(0)\| \leq \| \wtilde{f}_\ell(0) \| + \| \wtilde{f}_\ell(0)  - f_\ell(0)\|.
\end{equation}
and Lemma \ref{lem:dist_q_l_and f_tilde0} we have that 
\begin{equation}\label{eq:f_tilde0_bound_stone_l2}
        \| \wtilde{f}_\ell(0)\| \leq \frac{C_2\ln\left(\frac{1}{\delta}\right)}{n^{r_0}}
\end{equation}
with probability of at least $1-\delta$.

Now we focus on bounding $\| \wtilde{f}_\ell(0)  - f_\ell(0)\|$

  From \eqref{eq:f_tilde-f_integrals} we have that 
\begin{equation*}
\wtilde{f}_\ell(x) - f_\ell(x)=
\frac{\wtilde D\int\limits_{\mathbb{S}_{\wtilde D-1}}\theta g(x,\theta)^{\wtilde D+1} d\theta}{(\wtilde D+1)\int\limits_{\mathbb{S}_{\wtilde D-1}}g(x,\theta)^{\wtilde D} drd\theta}     
=
\frac{\wtilde D\int\limits_{\mathbb{S}_{\wtilde D-1}}\theta (g(x,\theta)^{\wtilde D+1} - \sigma^{\wtilde D+1}) d\theta}{(\wtilde D+1)\int\limits_{\mathbb{S}_{\wtilde D-1}}g(x,\theta)^{\wtilde D} drd\theta}     
\end{equation*}
or, looking at some direction $\vec{z}$ we have 
\begin{align*}
z^T\cdot (\wtilde{f}_\ell(x) - f_\ell(x))&=
\frac{\wtilde D\int\limits_{\mathbb{S}_{\wtilde D-1}}z (g(x,\theta)^{\wtilde D+1}- \sigma^{\wtilde D+1}) d\theta}{(\wtilde D+1)\int\limits_{\mathbb{S}_{\wtilde D-1}}g(x,\theta)^{\wtilde D} d\theta} 
\\
&\leq
\frac{\wtilde D}{(\wtilde D+1)}\sigma((1+4\alpha_\ell^2)^{\wtilde D+1} - 1)
\end{align*}

We note that for $x<1/n$ the following holds
\[
(1+x)^n-1<x((1+1/n)^n-1)/(1/n)<nx(e-1)<2nx.
\]
Using the above observation, and the fact that $4\alpha_\ell^2 <\wtilde D+1$ we have that 
\begin{align*}
z^T\cdot (\wtilde{f}_\ell(x) - f_\ell(x))&\leq
\frac{8\sigma\wtilde D}{(\wtilde D+1)}(\wtilde D+1)\alpha_\ell^2 \leq 8\sigma D \alpha_\ell^2
\end{align*}
Combining this with  \eqref{eq:f_ell0_triangle_l2} and \eqref{eq:f_tilde0_bound_stone_l2}, we have that for any $\delta>0$, for $\alpha_\ell >\frac{1}{D}$, and for number of samples $n>N$ large enough, 
\[
\|f_\ell(0)\| \leq 8\sigma D \alpha_\ell^2 + \frac{C_2\ln\left(\frac{1}{\delta}\right)}{n^{r_0}},
\]
with probability of at least $1-\delta$.
\end{proof}

\begin{Lemma}\label{lem:dist_q_l_and f_tilde0}
For any $\delta>0$ there is $N_\delta$ such that for any $n>N_\delta$, we have 
\[
\| \wtilde{f}_\ell(0)\| \leq \frac{C \ln\left(\frac{1}{\delta}\right)}{n^{r_0}}
\] 
with probability at least $1-\delta$, where $r_0 = \frac{k}{2k+d}$.
\end{Lemma}
\begin{proof}
First we note that 
\[
\wtilde f_\ell(0) = \wtilde f_{\ell - 0.5}(0) - \pi^*_{q_{\ell-1}, H_\ell}(0)
.\]
Then, from applying Theorem \ref{thm:vector_LPR} on $\wtilde{f}_{\ell-0.5}(0)$, and noting that $r_0 <1/2$,  we get the desired bound.
\end{proof}
% --- end Section_ProofStep2_f0_est.tex ---


\subsubsection{Bounding the error induced by the shifted origin}\label{sec:ShiftedCenter_err}
% --- begin Section_ProofStep2_BoundTf_move_origin.tex ---
\begin{figure}
	\centering
	\includegraphics[width=0.8\textwidth]{Figures/Manifold_ball_2.pdf}
	\caption{Illustration of the difference between $ p $ and $ f_0(0) $.}
	\label{fig:f0(0)vs_p}
\end{figure}


Back to Theorem \ref{thm:Step2} proof road-map see Figure \ref{tikz:thm33_lemmas}.

\begin{Lemma}\label{lem:shift_ang_general}
    Let $G_0$ be a $d$ dimensional linear space in $\RR^D$, and let $g_0:G_0 \rightarrow \RR^{D-d}$, such that the graph of $g_0$ is a manifold with reach bounded by $\tau$. Assume that $\maxangle (T_0g_0, G_0) \leq \alpha$. Let  $G_1$ be a $d$ dimensional linear space in $\RR^D$, such that $\maxangle(G_0,G_1)\leq \beta$. Define $g_1:G_1 \rightarrow \RR^{D-d}$ as the function who's graph coincides with the graph of $g_0$. 

    Then, for $\alpha \leq \pi/16$, $\beta\leq \beta_c$ where $\beta_c$ is some constant dependent only on $c_{\pi/4}$ of Lemma \ref{lem:M_is_locally_a_fuinction_clean}, and  $\|g_0(0)\| \leq \frac{3\tau}{4 \cdot 16} $, we have 
    \[
    \maxangle (G_1, T_0g_1)) \leq \maxangle(G_1, T_0g_0)  + \frac{ 8\|g_0(0)\|}{\tau}
    \]
\end{Lemma}

We first need a supporting lemma that will show us that $g_{1}$ exists, and specifically, $g_{1}(0)$ exist.
\begin{Lemma} \label{lem:rough_w_bound}
Under the conditions of Lemma \ref{lem:shift_ang_general}, $g_{1}(0)$ exists and 
\[
\|P_{T_0 g_0} (o_1 + (0, g_{1}(0))_{G_1} - \tilde o_1) \| \leq \tau/2
\]
\end{Lemma}
\begin{proof}
We begin with defining the coordinate system $(\tilde o_1, G_1)$ with $\tilde o_1 = (0, g_0(0))_{G_0}$. Let $\wtilde g_1:(\wtilde o_1,G_1) \simeq \RR^d \to \RR^{D-d}$ be the function defined in Lemma \ref{lem:M_is_locally_a_fuinction_clean}, such that 
\[
\Gamma_{\wtilde g_1} =  \MM \cap \textrm{Cyl}(\wtilde o_1,c_{\pi/4}\tau, \tau/2)
\]
From Lemma \ref{lem:M_is_locally_a_fuinction_clean} we know that $\wtilde g_1$ is defined for any $x\in \RR^d$ such that $\norm{x} \leq c_{\pi/4}\tau$.
Now, we denote $x_o = P_{G_1}(o_1 - \tilde o_1)$, the projection of $o_1$ onto the affine space defined by $(\tilde o_1, G_1)$.
From the assumptions we know that $\norm {\tilde o_1 - o_1} \leq  \frac{3\tau}{4 \cdot 16}$. 
Since $\maxangle (G_1,G_0) \leq \beta$ from Lemma \ref{lem:angle_space_to_vec} we have that $\norm{x_o} \leq \frac{3\tau}{4 \cdot 16} \sin \beta$. 
Thus, for sufficiently small $\beta$, depending only on $c_{\pi/4}$, $\|x_o\| \leq c_{\pi/4}
\tau$, and $\wtilde{g}_1(x_o)$ is therefore defined (by Lemma \ref{lem:M_is_locally_a_fuinction_clean}). Since $\wtilde g_1$ identifies with $g_1$ up to some shift in the domain and target, it follows that $g_1(0)$ is well defined.

Next we bound  $ \|P_{T_0g_0} (o_1 + (0, g_{1}(0))_{G_1} - \tilde o_1) \| $.
Since $\wtilde{g}_1(0) = 0$, from Lemma \ref{lem:f_bound_circle_H_clean} and the triangle inequality for maximal angles between flats
we have that 
\[
\|\wtilde{g}_1(x_o)\| \leq \tau\cos(\alpha + \beta)-\sqrt{\tau^2-(\norm{x_o} + \tau\sin(\alpha+\beta))^2}
\]
Substituting $\|x_o\|$ in the right hand side we set 
\begin{align}
\|\wtilde{g}_1(x_o)\| &\leq \tau\cos(\alpha + \beta)-\sqrt{\tau^2-(\frac{3\tau}{4 \cdot 16} \sin \beta + \tau\sin(\alpha+\beta))^2} \notag\\
&= \tau \left( \cos(\alpha + \beta)-\sqrt{1-(\frac{3}{4 \cdot 16} \sin \beta + \sin(\alpha+\beta))^2} \right)    
\end{align}
Since  
\begin{multline}
 \|P_{T_0g_0} (o_1 + (0, g_{1}(0))_{G_1} - \tilde o_1) \| \leq \norm{o_1 + (0, g_{1}(0))_{G_1} - \tilde o_1)}   = \sqrt{\|\wtilde{g}_1(x_o)\|^2 + \|x_o\|^2} 
\\
= \tau \sqrt{\left(\frac{3}{4 \cdot 16}\right)^2 \sin^2 \beta + \left( \cos(\alpha + \beta)-\sqrt{1-(\frac{3}{4 \cdot 16} \sin \beta + \sin(\alpha+\beta))^2} \right)^2 }
\end{multline}
which, for small enough $\beta$ and fixed $\alpha$ is smaller than $0.5 \tau$.
\begin{figure}
    \centering
    \includegraphics{./Figures/RoughWBound.pdf}
    \caption{Illustration of an angle change of a coordinate system. We have $(o_1, G_0)$ as some coordinate system. Locally we look at $\MM$ (marked in solid red) as a graph of $g_0:(o_1, G_0)\to G_0\perp$. The point $\tilde o_1$ equals $g_0(0)$. Let $G_1$ be some rotated coordinate system and describe $\MM$ as a local graph of $g_1:(o_1, G_1)\to G_1^\perp$.}
    \label{fig:G0G1}
\end{figure}
\end{proof}


Next we prove Lemma \ref{lem:shift_ang_general}

\begin{proof}[proof of Lemma \ref{lem:shift_ang_general}] 
We first note that from Lemma \ref{lem:rough_w_bound} $(0,g_1(0))_{G_1}\in \MM$ exists.
Then, we denote by $o_1$ the origin, $\tilde o_1 = o_1 + (0, g_0(0))_{G_0}$ and $\bar o_1 = o_1 + (0, g_1(0))_{G_1}$ (see Figure \ref{fig:G0G1}).
Thus, We can write $\bar o_1 = \tilde o_1 + w_x + w_y$, where 
\[
\begin{array}{c}
w_x = P_{T_0g_0}(\bar o_1 - \tilde o_1) \in T_0g_0 \\ 
w_y = \bar o_1 - w_x - \tilde o_1 = P_{T_0g_0^\perp}(\bar o_1-\tilde o_1) \in T_0g_0^\perp.    
\end{array}
\]
In this case, since $w_x$ and $w_y$ are orthogonal, we have 
\begin{equation}\label{eq:dist_f0_w}
    \|\bar o_1 - \tilde o_1\|^2 = \|w_x\|^2 + \|w_y\|^2.
\end{equation}
From Lemma \ref{lem:f_bound_circle_no_func_clean}, we have that  $\|w_y\| \leq \tau - \sqrt{\tau^2 - \|w_x\|^2}$. Since $\|w_x\|\leq \tau$, and by Remark \ref{rem:taylor_sqrt1-x2_clean} we have,
\begin{equation}\label{eq:wy_bound}
    \|w_y\| \leq \tau - \sqrt{\tau^2 - \|w_x\|^2} \leq \frac{\|w_x\|^2}{\tau}.
\end{equation}
Thus, in order to bound \eqref{eq:dist_f0_w}, we only need to bound $\|w_x\|$. Recall that  $\maxangle(G_0,G_1)\leq \beta$. By Lemma \ref{lem:angle_space_to_vec} we have that 
\[
\beta \geq \min\limits_{v\in G_0^\perp} \angle (v,\bar o_1).
\]
Taking the cosine of both sides we have
\begin{equation}\label{eq:cos_bound_wx_norm}
    \cos\left(\beta\right)\leq \max\limits_{v\in G_0^\perp} \cos(\angle (v,\bar o_1)) =\max\limits_{v\in G_0^\perp} \frac{<v,\bar o_1>}{\|\bar o_1\|}.
\end{equation}
    
For any unit vector $v\in G_0^\perp$ we have
\begin{equation}\label{eq:vw_prod_bound}
 <v,\bar o_1> = <v,\tilde o_1 + w_x + w_y> \leq \|\tilde o_1\| + <v,w_x> + \|w_y\|.   
\end{equation}
We also note that 
\[
    \|\bar o_1\| \geq \|w_x +w_y\| - \|\tilde o_1\| ,
\]
and since $w_x$ and $w_y$ are orthogonal, we have
\begin{equation}\label{eq:w_norm_bound}
\|\bar o_1\| \geq \sqrt{\|w_x\|^2 + \|w_y\|^2} - \|\tilde o_1\| \geq \|w_x\| - \|\tilde o_1\|.    
\end{equation}
rewriting Eq. \eqref{eq:cos_bound_wx_norm} using Eq. \eqref{eq:vw_prod_bound} and Eq. \eqref{eq:w_norm_bound}, we have
\[
    \cos\left(\beta\right)\leq \max\limits_{v\in H_0^\perp} \frac{<v,w>}{\|w\|} \leq \max\limits_{v\in H_0^\perp} \frac{\|\tilde o_1\| + <v,w_x> + \|w_y\| }{\|w_x\| - \|\tilde o_1\|}
\]
Since $v\in G_0^\perp$ and $w_x\in T_0g_0$, from Lemma \ref{lem:angle_space_perp_space} we have
\[
    <v,w_x> \leq \|w_x\| \cos(\frac{\pi}{2} - \alpha) = \|w_x\| \sin(\alpha),
\]
and thus
\begin{equation}\label{eq:wx_bound1}
    \cos\left(\beta\right) \leq
 \frac{\|\tilde o_1\| + \|w_x\| \sin(\alpha) + \|w_y\| }{\|w_x\| - \|\tilde o_1\|}.
\end{equation}
Since $\|w_y\| \leq \frac{\|w_x\|^2}{\tau}$ we can rewrite Eq. \eqref{eq:wx_bound1}
\begin{equation}\label{eq:order2_ineq_wx}
-\frac{1}{\tau}\|w_x\|^2+\|w_x\|\left(\cos\left(\beta\right) - \sin( \alpha)\right) - \|\tilde o_1\| \left( \cos\left(\beta\right) + 1 \right) \leq 0.    
\end{equation}
Since we have that $\alpha \leq \pi/16$ and $\beta \leq \pi/10 $, we have  $1 \geq \cos(\beta)\geq 0.95$ and $\cos (\pi/2-\alpha) \leq 0.2$, and then
\begin{equation}\label{eq:order2_ineq_wx_2}
-\frac{1}{\tau}\|w_x\|^2+0.75 \|w_x\|  -2 \|\tilde o_1\| \leq 0 
\end{equation}
The left hand side of Eq. \eqref{eq:order2_ineq_wx_2} is a second degree polynomial in $\|w_x\|$. The roots of this polynomial are:
\[
\begin{array}{cc}
    \xi_1 =  \frac{3\tau}{8}\left(1 - \sqrt{1-\frac{128}{9\tau} \|\tilde o_1\|}\right) \\
    \xi_2 = \frac{3\tau}{8} \left(1 + \sqrt{1-\frac{128}{9\tau} \|\tilde o_1\|}\right).
\end{array}
\]
by Remark \ref{rem:taylor_sqrt1-x2_clean} we have,
\begin{equation} \label{eq:x112_bound}
\begin{array}{ll}
    \xi_1 \leq& \frac{3\tau}{8} \left[1 - 1+\frac{128}{9\tau} \|\tilde o_1\|\right] = \frac{16}{3} \|\tilde o_1\| \\
    \xi_2 \geq &\frac{3\tau}{8} \left[1 + 1-\frac{128}{9\tau} \|\tilde o_1\|\right] = \frac{3\tau}{4} - \frac{16}{3} \|\tilde o_1\|.
\end{array}
\end{equation}
We have from Eq. \eqref{eq:order2_ineq_wx} that $\|w_x\|\leq \xi_1$ or $\|w_x\|\geq \xi_2$, or, from Eq. \eqref{eq:x112_bound} we have 
\[
\|w_x\|\leq \frac{16}{3} \|\tilde o_1\|  \mbox{~~~or~~~~} \|w_x\|\geq \frac{3\tau}{4} - \frac{16}{3} \|\tilde o_1\|.
\]
By Lemma \ref{lem:rough_w_bound}, and for $\|\tilde o_1\| \leq \frac{3\tau}{4 \cdot 16} $  we have that $\|w_x\| \leq \tau/2 \leq \frac{3\tau}{4} - \frac{16}{3} \|\tilde o_1\|$, and thus, we have   
\[
\|w_x\|\leq \frac{16}{3} \|\tilde o_1\|.
\]
From Eq. \eqref{eq:dist_f0_w} and \eqref{eq:wy_bound} we have that 
\[
    \|\tilde o_1-w\|^2 = \|w_x\|^2 + \|w_y\|^2 \leq \frac{16}{3}\|\tilde o_1\|^2 + \frac{16\|\tilde o_1\|^2 }{3\tau^2},
\]
and since $\tau > 1$ we have $\|\tilde o_1-w\|^2 \leq 32/3 \|\tilde o_1\|^2$

Finally, from Corollary 3 in \cite{boissonnat2017reach} we conclude that 
\[
\sin \frac{\maxangle(T_0g_0,T_0g_1(w))}{2} \leq \frac{\sqrt{32/3} \|\tilde o_1\|}{2\tau} \leq \frac{4 \|\tilde o_1\|}{2\tau} = \frac{ 2\|\tilde o_1\|}{\tau}.
\]
Since $x/2 \leq \sin x$
\[
\maxangle(T_0g_0,T_0g_1) \leq \frac{ 8\|\tilde o_1\|}{\tau}. 
\]
and since 
\[
\maxangle (T_0g_0, G_1) \leq \maxangle (T_0g_0, G_0) + \maxangle (G_0, G_1) \leq \alpha + \beta
\]
we have  
\[
\maxangle (T_0g_1, G_1) \leq \maxangle (T_0g_0, G_1) + \maxangle(T_0g_0,T_0g_1)\leq \maxangle(T_0g_0(0),G_1)  + \frac{ 8\|\tilde o_1\|}{\tau} \left(\leq \alpha+\beta +\frac{ 8\|\tilde o_1\|}{\tau}\right)
\]
\end{proof}

% --- end Section_ProofStep2_BoundTf_move_origin.tex ---

% --- end main_appendix.tex ---

\end{appendices}


\end{document} 

